\documentclass[11pt,oneside,openany]{book}
\usepackage[T1]{fontenc}
\usepackage{lmodern}
\usepackage{microtype}
\usepackage[a4paper,margin=28mm,headheight=14pt]{geometry}
\usepackage{amsmath,amssymb,amsthm,mathtools}
\usepackage[hidelinks,bookmarksopen,bookmarksnumbered]{hyperref}
\hypersetup{pdftitle={Variational Methods, Convex Analysis, and Optimality},
  pdfauthor={MathGenome},pdfsubject={Unified notes, P1-06}}

\newcommand{\R}{\mathbb R}
\newcommand{\N}{\mathbb N}
\newcommand{\C}{\mathbb C}
\newcommand{\norm}[1]{\lVert #1\rVert}
\newcommand{\pair}[2]{\langle #1,#2\rangle}
\DeclareMathOperator{\ri}{ri}
\DeclareMathOperator{\dom}{dom}
\DeclareMathOperator{\epi}{epi}
\DeclareMathOperator*{\argmin}{arg\,min}
\allowdisplaybreaks[1]
\setlength{\emergencystretch}{2em}
\theoremstyle{plain}
\newtheorem{theorem}{Theorem}[chapter]
\newtheorem{lemma}[theorem]{Lemma}
\newtheorem{proposition}[theorem]{Proposition}
\newtheorem{corollary}[theorem]{Corollary}
\theoremstyle{definition}
\newtheorem{definition}[theorem]{Definition}
\newtheorem{example}[theorem]{Example}
\theoremstyle{remark}
\newtheorem{remark}[theorem]{Remark}

\begin{document}
\frontmatter
\hypersetup{pageanchor=false}
\begin{titlepage}
\thispagestyle{empty}
\vspace*{31mm}
{\large MathGenome \enspace / \enspace P1-06\par}
\vspace{7mm}
\hrule height 0.8pt
\vspace{11mm}
{\Huge\bfseries Variational Methods,\\[2mm] Convex Analysis,\\[2mm] and Optimality\par}
\vspace{13mm}
{\large Unified notes\par}
\vspace{4mm}
{\normalsize Convex separation, conjugacy, and dual attainment\\
Constrained optimality, weak compactness, and variational existence\par}
\vfill
{\large 8 November 2026\par}
\end{titlepage}
\hypersetup{pageanchor=true}
\tableofcontents

\chapter*{Scope and prerequisites}
\addcontentsline{toc}{chapter}{Scope and prerequisites}

The first three chapters establish convex geometry, continuous-dual
calculus, and existence in finite-dimensional, reflexive, and Hilbert
spaces. The next three use these tools for perturbation duality, linear
and conic programs, and smooth optimality conditions. The final four
develop classical variation, weak lower semicontinuity and direct
methods for Sobolev energies, flux duality, and compact fixed-point
methods. The source map links the 69 groups to specific parts of the
text and to the initial 30 result contracts.

Throughout, locally convex spaces are real and Hausdorff unless stated
otherwise. For such a space \(X\), its continuous dual is \(X'\), with
pairing \(\langle x',x\rangle=x'(x)\). In finite dimensions we use
Euclidean coordinates when gradients and Hessians are convenient.
A proper extended-real convex function has no \(-\infty\) value and
is finite at some point. Cones contain the origin. For a closed convex cone \(K\), the positive
dual cone is \(K^*=\{y:\langle y,x\rangle\ge 0\text{ for all }x\in K\}\).
For the smooth minimization convention \(g_i\le0,\ h_j=0\), we use
\(L=f+\sum_i\lambda_i g_i+\sum_j\nu_j h_j\), with \(\lambda_i\ge0\).
The sign of each program dual will also be checked directly from its gap.
In the complex Hilbert variational problem the form is conjugate linear
in the unknown and linear in the test argument, and the right-hand
functional is linear in that argument.

The analytic prerequisites include the algebraic Hahn--Banach
extension theorem for sublinear bounds, Banach--Alaoglu, Hilbert-space
Riesz representation, the contraction theorem, finite-dimensional
linear algebra, Taylor's theorem, the implicit-function theorem, and
existence and continuous dependence for linear ordinary differential
equations. We use standard Lebesgue integration, H\"older's inequality,
the \(L^p\) duality and reflexivity theorems for the Lebesgue spaces in
question, density and mollification, Arzelà--Ascoli, and the first-order extension and
trace theorems on bounded Lipschitz open sets
\cite[Chapter 9]{Brezis}\cite[Unit U5]{MathGenomeP102}. The compactness
needed for the direct method is proved with these inputs. Green's
formula on a regular planar domain is used in the isoperimetric
example. The relevant hypotheses of these results are recorded at
their uses.

Rockafellar and Boyd--Vandenberghe are the principal references for
finite-dimensional convexity and duality. Ekeland--T\'emam and Brezis
supply the function-space versions. Dacorogna and Liberzon guide the
classical and Sobolev variational arguments, while Nocedal--Wright
and Mangasarian--Fromovitz guide the smooth multiplier theory.
Applications that require a separate theory, such as general
relaxation, elliptic regularity, and numerical algorithms, are located
by their interfaces.

\mainmatter
\chapter{Convex Geometry and Separation}

This chapter develops the geometry used in convex duality and in the
first-order analysis of constraints. In finite dimension the proofs may
use any Euclidean coordinates. A prime denotes the continuous dual, and
\(\langle y,x\rangle=y(x)\) for \(y\in X'\), \(x\in X\).

\section{Affine hulls and relative interiors}\label{v01-01}

\begin{definition}
An affine set contains every affine combination of its points, where
the coefficients sum to one. A convex set contains all convex
combinations, which are affine combinations with nonnegative
coefficients. A convex cone contains zero and is closed under addition
and nonnegative scalar multiplication. The affine hull, convex hull
and conic hull are the smallest sets of the respective types containing
a given set. A face of a convex set \(C\) is a convex subset \(F\)
such that \(x,y\in C\) and \((1-t)x+ty\in F\) for some \(0<t<1\)
imply \(x,y\in F\). An extreme point \(x\) is one for which \(\{x\}\)
is a face. For a nonempty finite-dimensional convex set \(C\), its
relative interior \(\operatorname{ri}C\) is its interior in
\(\operatorname{aff}C\).
\end{definition}

\begin{lemma}[Hull descriptions and elementary operations]
\label{lem:hull-descriptions}
The affine, convex and conic hulls are the sets of finite combinations
of their respective types. Intersections, affine images, affine
inverse images, Cartesian products and Minkowski sums preserve convexity.
A linear image of a convex cone is a convex cone.
\end{lemma}

\begin{proof}
Combining two finite expressions of the specified type with coefficients
of that type again gives such an expression, by multiplying the
coefficients. The conic hull allows the empty combination, equal to
zero. This proves closure under the defining operations. Every set
containing the generators and closed under those operations contains
the displayed expressions. For intersections and inverse images,
the defining membership conditions hold in every factor or every
preimage. Images preserve affine combinations. Products and sums
preserve convex combinations by taking the same two weights in
each coordinate.
\end{proof}

\begin{theorem}[Relative interiors and half-open segments]
\label{thm:relative-segment}
Every nonempty convex set \(C\) in finite dimension has nonempty
relative interior. If \(a\in\operatorname{ri}C\) and
\(b\in\overline C\), then \((1-t)a+tb\in\operatorname{ri}C\)
for \(0\leq t<1\).
\end{theorem}

\begin{proof}
If the affine hull \(H\) has dimension \(d\), choose \(d+1\)
affinely independent points in \(C\). They exist by extending an
affinely independent list until its affine span contains \(C\).
Their simplex lies in \(C\), and its barycenter has strictly positive
barycentric coordinates. By continuity of those coordinates a
neighborhood of the barycenter in \(H\) lies in the simplex. In
dimension zero the singleton is open in \(H\).

For the segment assertion choose \(r>0\) such that the relative
ball \(a+B_H(0,r)\) is in \(C\), where \(B_H(0,r)\) denotes a
ball in the direction space of \(H\). Fix \(0<t<1\), and take
\(b_0\in C\) with \(t\lVert b-b_0\rVert<(1-t)r/2\).
For every direction-space vector \(h\) with
\(\lVert h\rVert<(1-t)r/2\),
\[
 (1-t)a+tb+h
  =(1-t)\left(a+\frac{t(b-b_0)+h}{1-t}\right)+t b_0 \in C.
\]
This gives a relative ball about the desired point. The endpoint
\(t=0\) is the hypothesis.
\end{proof}

These conventions agree with \cite[\S\S1--3 and 6]{Rockafellar}
and \cite[\S2.1]{Boyd}.

\section{Relative interior calculus}\label{v01-02}

\begin{lemma}[An extendible-segment test]\label{lem:ri-extendible}
For nonempty convex \(C\) in finite dimension and \(x\in C\),
one has \(x\in\operatorname{ri}C\) exactly when for each \(y\in C\)
there is an \(s>0\) such that \(x+s(x-y)\in C\).
\end{lemma}

\begin{proof}
A relative ball at \(x\) gives the extension for sufficiently small
\(s\). Conversely choose \(y\in\operatorname{ri}C\) by
Theorem~\ref{thm:relative-segment}. The extension expresses \(x\)
as a convex combination of \(y\) and \(x+s(x-y)\) in which the
coefficient of \(y\) is positive. The segment theorem applies.
\end{proof}

\begin{theorem}[Relative interior identities]\label{thm:ri-calculus}
For a nonempty convex set \(C\) in finite dimension,
\(\overline{\operatorname{ri}C}=\overline C\) and
\(\operatorname{ri}\overline C=\operatorname{ri}C\).
For a finite family of convex sets \(C_i\) with a common
relative-interior point,
\[
 \operatorname{ri}\Big(\bigcap_i C_i\Big)=\bigcap_i\operatorname{ri}C_i,
 \qquad
 \overline{\bigcap_i C_i}=\bigcap_i\overline{C_i}.
\]
For an affine map \(A\) between finite-dimensional spaces,
\(\operatorname{ri}(AC)=A(\operatorname{ri}C)\).
If \(A(E)\) meets \(\operatorname{ri}D\), where \(D\) is convex,
then \(\operatorname{ri}(A^{-1}D)=A^{-1}(\operatorname{ri}D)\).
For nonempty convex \(C,D\) in the same space,
\(\operatorname{ri}(C+D)=\operatorname{ri}C+\operatorname{ri}D\).
\end{theorem}

\begin{proof}
Fix \(a\in\operatorname{ri}C\). Approaching any \(b\in\overline C\)
along the half-open segment from \(a\) proves the first closure
identity. The affine hull of \(\overline C\) is that of \(C\),
since an affine subspace is closed in finite dimension. If
\(x\in\operatorname{ri}\overline C\) and \(x\neq a\), a short extension
\(x+s(x-a)\) lies in \(\overline C\). The segment theorem puts
\(x\) in \(\operatorname{ri}C\). The other inclusion follows
from the relative-ball definition.

For the intersection choose a common point \(a\) in all
\(\operatorname{ri}C_i\). For \(x\) in all those relative interiors,
extend the segment from any \(y\in\bigcap_i C_i\) beyond \(x\)
in every \(C_i\). The minimum of the finitely many positive extension
parameters works by convexity. If \(x\in\operatorname{ri}(\bigcap_i C_i)\),
extend from \(a\) beyond \(x\) inside the intersection and use
the half-open segment theorem in each \(C_i\). The closure identity
follows by approaching any point of \(\bigcap_i\overline{C_i}\)
on the segment from \(a\).

The affine map \(A\) maps \(\operatorname{aff}C\) onto
\(\operatorname{aff}(AC)\) and is open between those affine hulls.
Indeed its linear part restricts to a linear isomorphism from a
complement of the kernel onto the direction space of the image.
This gives \(A(\operatorname{ri}C)\subset\operatorname{ri}(AC)\).
If \(y\in\operatorname{ri}(AC)\) is not \(Aa\), extend the segment
from \(Aa\) beyond \(y\) to \(y_1\in AC\). Lift \(y_1=Ab\) with
\(b\in C\). The corresponding point strictly between \(a,b\) is
in \(\operatorname{ri}C\) and maps to \(y\). The case \(y=Aa\) is
immediate.

For inverse images choose \(a\) with \(Aa\in\operatorname{ri}D\).
If \(Ax\in\operatorname{ri}D\), extend from each \(Ay\in D\) beyond
\(Ax\). By affinity this is an extension from \(y\) beyond \(x\)
in \(A^{-1}D\), so Lemma~\ref{lem:ri-extendible} applies. Conversely
an extension beyond \(x\) from \(a\) in the inverse image, followed
by the half-open segment theorem in \(D\), shows \(Ax\in\operatorname{ri}D\).
Finally the relative interior of a product is the product of the
relative interiors, since the affine hull is the product of the
affine hulls and relative product neighborhoods characterize
interior. Apply the image identity to addition.
\end{proof}

For comparison see \cite[Theorems 6.3--6.7]{Rockafellar}. The image
identity concerns relative interiors. Closedness of linear images is
addressed in Section~\ref{v01-07}.

\section{Finite convex hulls, extreme points, and Helly's theorem}\label{v01-03}

\begin{theorem}[Carath\'eodory]\label{thm:caratheodory}
If \(S\) has affine dimension at most \(d\), each point of its convex
hull is a convex combination of at most \(d+1\) points of \(S\).
\end{theorem}

\begin{proof}
Start with a finite convex combination and discard zero weights.
If its support is affinely dependent, choose numbers
\(c_i\), not all zero, with \(\sum_i c_i x_i=0\) and \(\sum_i c_i=0\).
Some \(c_i\) is positive. If the weights are \(\lambda_i\), set
\(t=\min_{c_i>0}\lambda_i/c_i\). The new weights
\(\lambda_i-tc_i\) remain nonnegative, sum to one and represent the
same point. At least one is zero. Iteration leaves affinely
independent support and therefore at most \(d+1\) points.
\end{proof}

\begin{lemma}[Relative boundary support in Euclidean space]
\label{lem:euclidean-boundary-support}
Let \(C\) be a nonempty closed convex set in finite dimension.
If \(x\in C\setminus\operatorname{ri}C\), a nonzero linear functional
on the direction space of \(\operatorname{aff}C\) supports \(C\) at \(x\).
It is strictly smaller than its value at \(x\) at some point of \(C\).
\end{lemma}

\begin{proof}
Choose exterior points \(z_j\) of \(C\) in its affine hull converging
to \(x\). A closest \(c_j\in C\) to \(z_j\) exists, because
a distance-minimizing sequence can be restricted to a closed bounded
ball. For \(u\in C\), minimizing squared distance along
\(c_j+t(u-c_j)\) for \(0\le t\le 1\) gives
\((z_j-c_j)\cdot(u-c_j)\le0\) by the right derivative at zero.
Also \(\lVert c_j-z_j\rVert\le \lVert x-z_j\rVert\), so \(c_j\to x\).
A subsequence of the unit vectors in direction \(z_j-c_j\)
converges to a unit vector \(n\) in the direction space of the
affine hull. Thus \(n\cdot(u-x)\le0\) for all \(u\in C\).
If equality held throughout \(C\), its affine hull would lie in
a proper hyperplane with normal \(n\), which is impossible.
\end{proof}

\begin{theorem}[Extreme point representation]
\label{thm:extreme-representation}
A nonempty compact convex set in finite dimension is the convex hull
of its extreme points.
\end{theorem}

\begin{proof}
Induct on affine dimension. A singleton is its own extreme point.
If \(x\) lies on the relative boundary, the support functional in
Lemma~\ref{lem:euclidean-boundary-support} cuts out a nonempty compact
convex set of smaller affine dimension containing \(x\).
The equality condition for a supporting linear inequality is
preserved by decomposing a convex combination with strictly positive
weights, so this set is a face of \(C\). Each of its extreme points
is extreme in \(C\), and the induction hypothesis represents \(x\)
using them. If \(x\) is in the relative interior and the affine
dimension is positive, a line in the affine hull through \(x\)
intersects \(C\) in a compact interval with distinct endpoints.
The endpoints are on the relative boundary by maximality, and their
convex combinations include \(x\).
\end{proof}

\begin{lemma}[Radon's partition]\label{lem:radon-partition}
Any \(d+2\) points in \(\mathbb R^d\), indexed with repetitions
allowed, can be partitioned into two nonempty index sets whose
convex hulls meet.
\end{lemma}

\begin{proof}
Take a nonzero affine dependence \(c\) of the points. Because the
coefficients sum to zero, their positive and negative parts have the
same strictly positive sum. The affine dependence gives equality
of the two convex combinations obtained by normalizing those
parts. Indices with zero coefficient may be put in either block.
\end{proof}

\begin{theorem}[Finite Helly]\label{thm:helly}
Let \(C_1,\dots,C_m\) be convex subsets of \(\mathbb R^d\) with
\(m\ge d+1\). If every subfamily of \(d+1\) members has nonempty
intersection, then the whole family has nonempty intersection.
\end{theorem}

\begin{proof}
Induct on \(m\), the case \(m=d+1\) being the hypothesis.
For \(m>d+1\) the induction hypothesis supplies, for each
\(i=1,\dots,d+2\), a point \(x_i\) in every \(C_j\) except
possibly \(C_i\). Take a common point of the two convex hulls
given by Lemma~\ref{lem:radon-partition}. For any \(j\), at least
one block omits the index \(j\), and all the points of that block
belong to \(C_j\). The common point is in \(C_j\) by convexity.
\end{proof}

See \cite[\S\S17, 18, and 21]{Rockafellar}.

\section{Separation and supporting hyperplanes}\label{v01-04}

\begin{theorem}[Open and compact--closed separation]
\label{thm:geometric-separation}
Let \(A,B\) be disjoint nonempty convex subsets of a real
Hausdorff locally convex topological vector space. If \(A\)
is open there are a nonzero continuous linear functional \(l\)
and \(\alpha\in\mathbb R\) with
\(l(a)<\alpha\le l(b)\) for all \(a\in A,\ b\in B\).
If \(A\) is compact and \(B\) is closed there is such a functional
with \(\sup_A l<\inf_B l\).
\end{theorem}

\begin{proof}
In the open case \(A-B\) is open, convex and excludes \(0\).
Fix \(z_0\in A-B\) and let \(U=A-B-z_0\). This open convex set
contains \(0\). Its gauge
\[
 p(v)=\inf\{t>0:v\in tU\}
\]
is finite since \(U\) is absorbing, and is nonnegative,
positively homogeneous and subadditive. For subadditivity, add
expressions \(v\in sU,\ w\in tU\), and use convexity to put
\(v+w\) in \((s+t)U\) before taking infima. One has \(U=\{p<1\}\):
membership in \(tU\) for \(t<1\) implies membership in \(U\)
by convexity, while openness gives a slight dilation of any point
of \(U\) that remains in \(U\). Since \(-z_0\notin U\),
\(p(-z_0)\ge1\). Define a linear functional on the span of
\(-z_0\) by \(l(t(-z_0))=t\,p(-z_0)\). For \(t\ge0\) this is
dominated by \(p\) by homogeneity, and for \(t<0\) it is dominated
by the nonnegative \(p\). The real Hahn--Banach extension theorem
extends it with \(l\le p\). It is continuous because on
\(U\cap(-U)\) both \(l\) and \(-l\) are bounded by one.
For \(z\in A-B\),
\[
 l(z)=l(z_0)+l(z-z_0)<-p(-z_0)+1\le0.
\]
Consequently \(l(a)<l(b)\) for all pairs, and
\(\alpha=\sup_A l\) is finite and at most \(\inf_B l\).
Because \(A\) is open and \(l\neq0\), each \(a\) has a nearby
point in \(A\) of larger \(l\)-value, giving \(l(a)<\alpha\).

In the compact--closed case \(A-B\) is closed. To see this,
for \(z\notin A-B\), every \(a-z\) with \(a\in A\) lies outside
\(B\). Continuity of subtraction gives neighborhoods of \(a\)
and \(z\) whose differences avoid \(B\). A finite subcover
of \(A\), and the intersection of the corresponding
neighborhoods of \(z\), show \(z\) has a neighborhood outside
\(A-B\). Choose an open convex neighborhood \(V\) of \(0\)
such that \(-V\) avoids \(A-B\), using local convexity.
Then \(A+V\) is open convex and disjoint from \(B\).
The open case gives \(l\neq0\) with
\(\sup_{A+V}l\le\inf_B l\). Choose \(v\in V\) with \(l(v)>0\).
This gives \(\sup_A l+l(v)\le\inf_B l\), the required gap.
\end{proof}

\begin{corollary}[Finite-dimensional support and proper separation]
\label{thm:proper-separation}
Every boundary point of a nonempty closed convex set in a
positive-dimensional Euclidean space admits a supporting hyperplane.
A relative boundary point admits one within the affine hull.
Two nonempty convex sets \(C,D\) in finite dimension have disjoint
relative interiors exactly when there exists a linear functional
\(l\) with \(l(c)\le l(d)\) for all pairs, and strict inequality
for at least one pair.
\end{corollary}

\begin{proof}
For ambient support when the affine hull is full, use
Lemma~\ref{lem:euclidean-boundary-support}. If the affine hull
is proper, an ambient hyperplane containing it gives support,
and the relative assertion is the lemma.

For proper separation put \(Q=D-C\).
Theorem~\ref{thm:ri-calculus} shows that the two relative
interiors meet exactly when \(0\in\operatorname{ri}Q\).
If \(0\notin\overline Q\), apply
Theorem~\ref{thm:geometric-separation} to the singleton \(0\)
and \(\overline Q\). If \(0\in\overline Q\setminus\operatorname{ri}Q\),
then \(0\) is a relative boundary point of \(\overline Q\)
by Theorem~\ref{thm:ri-calculus}. Reverse the sign of the
support functional of the lemma to obtain \(l\ge0\) on
\(\overline Q\) and a strict inequality at some point.
Extend this functional from the direction space of
\(\operatorname{aff}Q\) to the ambient space. Density of \(Q\)
gives a point of strict inequality in \(Q\). These two cases
give the asserted separation. Conversely, if \(0\in\operatorname{ri}Q\),
a relative ball around \(0\) in \(Q\) forces any linear \(l\ge0\)
on \(Q\) to vanish on the direction space of its affine hull.
It cannot then give any strict inequality.
\end{proof}

See \cite[\S11]{Rockafellar}, \cite[Chapter I, \S1]{EkelandTemam}
and \cite[\S1.2]{Brezis}. Proper separation concerns all pairs
and at least one strict inequality. The compact--closed result
provides a uniform gap.

\section{Halfspaces, dual cones, and polars}\label{v01-05}

\begin{definition}
For \(S\subset X\) define the support function
\(\sigma_S(l)=\sup_{x\in S}\langle l,x\rangle\) on \(X'\).
The negative polar is \(S^\circ=\{l:\sigma_S(l)\le0\}\).
For a cone \(K\) the positive dual cone is
\(K^*=-K^\circ\). For a subset of \(X'\) its polar in \(X\)
is defined by the same nonpositive pairing inequality.
\end{definition}

\begin{theorem}[Halfspace representation and cone bipolarity]
\label{thm:cone-bipolar}
A closed convex set in a real Hausdorff locally convex space
is the intersection of the closed continuous-linear halfspaces
containing it. For any subset \(S\) of that space,
\(S^{\circ\circ}\) is its closed convex conic hull. In particular
a closed convex cone is recovered by taking the positive dual
cone twice, using the dual pairing.
\end{theorem}

\begin{proof}
For a point \(x\) outside a nonempty closed convex set, separate
the singleton \(x\) strictly from that set by
Theorem~\ref{thm:geometric-separation}. Reversing the sign if
necessary gives a closed halfspace containing the set and excluding
\(x\). For the empty set one may use the empty halfspace
\(0\le -1\).

Let \(K\) be the closed convex conic hull of \(S\). A continuous
linear functional nonpositive on \(S\) is nonpositive on \(K\)
by combination and closure, so \(K^\circ=S^\circ\).
For \(x\notin K\), strict separation gives
\(\sup_K l<l(x)\). Since \(K\) contains \(0\) and is a cone,
finiteness of \(\sup_K l\) forces \(l\le0\) on \(K\) and
\(\sup_K l=0\). Thus \(l\in S^\circ\) and \(l(x)>0\),
excluding \(x\) from the bipolar. Every \(x\in K\) is
in the bipolar by the definition of the polar. Negating the
functional twice changes the polar convention to the positive
dual convention.
\end{proof}

See \cite[Theorem 11.5 and \S\S13--15]{Rockafellar}
and \cite[\S\S2.5--2.6]{Boyd}.

\section{Recession directions and convex sublevels}\label{v01-06}

\begin{definition}
For nonempty closed convex \(C\subset E\) define
\[
 \operatorname{rec}C
   =\{d:x+td\in C\text{ for all }x\in C,\ t\ge0\}.
\]
Its lineality space is \(\operatorname{rec}C\cap(-\operatorname{rec}C)\).
\end{definition}

\begin{proposition}[Rays, escape, and recession]\label{prop:rec-set}
The recession set is a closed convex cone. For a given \(d\), a
ray \(x_0+[0,\infty)d\) in \(C\) from one point suffices for
\(d\in\operatorname{rec}C\). In finite dimension a nonempty closed
convex set is unbounded exactly when its recession cone contains
a nonzero vector.
\end{proposition}

\begin{proof}
A ray from \(x_0\) gives, for any \(x\in C\), \(t\ge0\),
and \(s>t\), the point
\[
  (1-t/s)x+(t/s)(x_0+sd)=x+td+(t/s)(x_0-x)\in C.
\]
Let \(s\to\infty\) and use closedness. The defining conditions
are closed in \(d\) and are preserved by addition and nonnegative
scaling, giving the cone statement. For the unbounded case take
\(x_j\in C\) with \(r_j=\lVert x_j-x_0\rVert\to\infty\).
A subsequence of \((x_j-x_0)/r_j\) converges to a unit vector \(d\).
For fixed \(t\ge0\), the points \(x_0+t(x_j-x_0)/r_j\)
are eventually in \(C\) by convexity. Closedness gives the ray
from \(x_0\) in direction \(d\). Conversely such a nonzero ray
is unbounded.
\end{proof}

\begin{proposition}[Recession of a function]\label{prop:rec-function}
Let \(f:E\to(-\infty,+\infty]\) be proper, convex and lower
semicontinuous in finite dimension. For \(x\) with \(f(x)<\infty\),
\[
 f^\infty(d)
  :=\inf\{r\in\mathbb R:(d,r)\in\operatorname{rec}(\operatorname{epi} f)\}
  =\sup_{t>0}\frac{f(x+td)-f(x)}{t}.
\]
The infimum of an empty fiber is \(+\infty\). The quotients are
nondecreasing in \(t\), and the supremum is independent of \(x\).
For any nonempty sublevel \(C_\alpha=\{f\le\alpha\}\) with
\(\alpha\in\mathbb R\),
\(\operatorname{rec}C_\alpha=\{d:f^\infty(d)\le0\}\).
\end{proposition}

\begin{proof}
The epigraph is nonempty, closed and convex. By
Proposition~\ref{prop:rec-set}, membership of \((d,r)\) in
its recession cone is equivalent to the ray from
\((x,f(x))\) being in the epigraph, which says exactly that
\(r\) bounds all the displayed quotients. Convexity along
the ray proves their monotonicity. The quotient at \(t=1\)
is real or \(+\infty\), so the supremum is never \(-\infty\).
For \(d\in\operatorname{rec}C_\alpha\), start at \(x\in C_\alpha\)
and bound the quotient by \((\alpha-f(x))/t\). Monotonicity
and \(t\to\infty\) show \(f^\infty(d)\le0\).
Conversely this last bound keeps the whole ray from each
\(x\in C_\alpha\) in the sublevel.
\end{proof}

In particular, a nonempty sublevel is unbounded exactly when
there is a nonzero \(d\) with \(f^\infty(d)\le0\), by
Proposition~\ref{prop:rec-set}.

See \cite[\S8, in particular Theorem 8.5]{Rockafellar}.

\section{When linear images are closed}\label{v01-07}

\begin{theorem}[A closed-image criterion]\label{thm:closed-image}
Let \(C\) be nonempty closed convex and \(A:E\to F\) linear
between finite-dimensional spaces. If
\(\operatorname{rec}C\cap\ker A\) is contained in the lineality
space of \(\operatorname{rec}C\), then \(AC\) is closed.
\end{theorem}

\begin{proof}
Let \(L=\operatorname{rec}C\cap\ker A\). The hypothesis implies
\(L=-L\), so it is a subspace and \(C+L=C\).
For a convergent sequence \(Ax_j\) with \(x_j\in C\), replace
\(x_j\) by its component in \(L^\perp\), which is still in \(C\)
and has the same image. Fix \(x_0\in C\cap L^\perp\), which
exists by the same replacement. If the replaced sequence is
unbounded, pass to a subsequence and normalize \(x_j-x_0\)
as in Proposition~\ref{prop:rec-set}. The unit limit belongs
to \(\operatorname{rec}C\cap L^\perp\). Boundedness of
\(A(x_j-x_0)\) gives \(Ad=0\) for this limit \(d\).
Hence \(d\in L\cap L^\perp=\{0\}\), a contradiction.
A bounded subsequence of representatives converges to some
\(x\in C\), which maps to the original limit.
\end{proof}

\begin{example}[A nonclosed image of a closed cone]
The cone of positive semidefinite \(2\)-by-\(2\) real symmetric
matrices is closed, as it is the intersection of the closed
inequalities \(u^{\mathsf T}Mu\ge0\).
Its linear image taking \(M=\left(\begin{smallmatrix}a&b\\b&c\end{smallmatrix}\right)\)
to \((a,b)\) is
\[
 \{(a,b):a>0\}\cup\{(0,0)\}.
\]
Indeed \(a<0\) is impossible. At \(a=0\), applying the quadratic
form to \((t,1)\) for all real \(t\) gives \(b=0\).
If \(a>0\), choose \(c=b^2/a\) and complete the square.
Thus \((0,1)\) belongs to the closure of the image and not to the image.
\end{example}

See \cite[\S9, Theorem 9.1]{Rockafellar}.

\section{Polyhedra and finite elimination}\label{v01-08}

\begin{definition}
A polyhedron is an intersection of finitely many affine closed
halfspaces in a finite-dimensional real space. An affine
equality can be represented by two inequalities.
\end{definition}

\begin{theorem}[Fourier--Motzkin elimination]\label{thm:poly-image}
A coordinate projection of a polyhedron is a polyhedron.
Every affine image of a polyhedron is therefore a polyhedron
and is closed.
\end{theorem}

\begin{proof}
To eliminate a scalar variable \(t\), normalize its nonzero
coefficients so that each affected row is either
\(t\ge L_i(x)\) or \(t\le U_j(x)\), with affine \(L_i,U_j\).
The rows with zero coefficient constrain only \(x\).
Given \(x\) satisfying those rows, an admissible \(t\) exists
exactly when \(L_i(x)\le U_j(x)\) for all pairs. When both
lists are nonempty, the largest lower bound is at most the
smallest upper bound and either gives an admissible choice.
When one list is empty, choose \(t\) at or beyond the extremal
bound of the other list, or choose any \(t\) when both are
empty. This gives finitely many affine inequalities in \(x\).
Iterating proves the coordinate projection assertion.
The graph over the polyhedron of any affine map is a
polyhedron in the product space, with equalities specifying
the graph. Projecting to the target proves the assertion.
\end{proof}

Compare \cite[\S\S19--20 and 22]{Rockafellar} and
\cite[\S13.2]{NocedalWright}.

\section{Finite cones and theorems of the alternative}\label{v01-09}

\begin{lemma}[Finite generation gives closedness]
\label{lem:finitely-generated-closed}
The cone generated by finitely many vectors in finite dimension is closed.
\end{lemma}

\begin{proof}
In a nonnegative expression \(\sum_i \lambda_i a_i\), discard zero
weights. If the remaining vectors are dependent, choose a linear
dependence \(c_i\) with at least one positive coefficient, reversing
the sign of the dependence if necessary. Subtract \(t c_i\) from
the weights, with \(t=\min_{c_i>0}\lambda_i/c_i\).
All weights stay nonnegative and one vanishes, while the sum is
unchanged. Iteration gives an independent support, possibly empty.
For a convergent sequence of cone elements, choose such expressions.
A subsequence uses the same support because there are finitely many
supports. The span of that independent list is closed and its
coordinate map has a continuous inverse on the span.
The coefficients converge to nonnegative coefficients for the
limit. If the support is empty the sequence on that subsequence is zero.
\end{proof}

\begin{theorem}[Farkas' alternative]\label{thm:farkas}
For a real matrix \(A\) and a real vector \(b\) exactly one
of the following holds:
\[
  Ax=b \text{ for some }x\ge0,\qquad
  A^{\mathsf T}y\ge0,\ b^{\mathsf T}y<0 \text{ for some }y.
\]
Inequalities on coordinate vectors are componentwise.
\end{theorem}

\begin{proof}
The image of the nonnegative orthant is the closed cone
generated by the columns of \(A\), by
Lemma~\ref{lem:finitely-generated-closed}.
If \(b\) is outside that cone, separate it strictly using
Theorem~\ref{thm:geometric-separation}, with the singleton on
the lower side. The separating linear functional is
nonnegative on the cone, since otherwise scaling one vector
makes the infimum \(-\infty\). The infimum is then zero.
Its Euclidean coordinate vector \(y\) gives the stated inequalities.
Both alternatives cannot hold, since
\(b^{\mathsf T}y=x^{\mathsf T}A^{\mathsf T}y\ge0\) in the first.
\end{proof}

\begin{theorem}[Gordan's alternative]\label{thm:gordan}
For a real matrix \(B\), exactly one of the following holds:
\[
  Bd<0\text{ for some }d,\qquad
  B^{\mathsf T}\mu=0
     \text{ for some }\mu\ge0,\ \mu\ne0 .
\]
\end{theorem}

\begin{proof}
If \(B\) has no rows, the strict system is vacuous and the
second alternative is impossible. Otherwise, failure of the strict
system means the image subspace of \(B\) is disjoint from the
open negative orthant. Separate them by
Theorem~\ref{thm:geometric-separation}, putting the orthant on
the lower side. A linear functional bounded below on a subspace
vanishes on it. The supremum on the negative orthant is at least
zero by approaching the origin, so the functional is nonpositive
there. By letting any one negative coordinate grow arbitrarily
large in magnitude while keeping the others fixed, all its
coefficients must be nonnegative. Nonzeroness of the separating
functional gives the required \(\mu\).
If both alternatives held, \(\mu^{\mathsf T}Bd\) would be
zero and strictly negative, a contradiction.
\end{proof}

See \cite[\S22]{Rockafellar} and \cite[\S5.8.3]{Boyd}.
Both results may be applied to a subspace of directions by
choosing coordinates on that subspace.

\section{Tangent and normal cones of polyhedra}\label{v01-10}

\begin{definition}
For \(x\in C\) in a Euclidean space, the contingent tangent cone
\(T_C(x)\) consists of all limits of \((x_k-x)/t_k\)
with \(x_k\in C\) and \(t_k>0\) tending to zero.
For convex \(C\), its normal cone at \(x\) is
\(N_C(x)=\{l:\langle l,y-x\rangle\le0\ \forall y\in C\}\).
\end{definition}

\begin{theorem}[Active linear constraints]\label{thm:poly-normal}
Let \(P=\{x:a_i\cdot x\le b_i,\ e_j\cdot x=d_j\}\)
be a polyhedron and \(x\in P\). Write
\(I(x)=\{i:a_i\cdot x=b_i\}\).
Then
\[
\begin{split}
 T_P(x)&=\{v:a_i\cdot v\le0\ (i\in I(x)),\ e_j\cdot v=0\ (\forall j)\},\\
 N_P(x)&=\left\{\sum_{i\in I(x)}\lambda_i a_i+\sum_j\nu_j e_j:
                   \lambda_i\ge0,\ \nu_j\in\mathbb R\right\}.
\end{split}
\]
No strict feasibility condition is required.
\end{theorem}

\begin{proof}
Each active inequality and equality gives the stated restrictions
on contingent limits by dividing by \(t_k\) and taking the limit.
Conversely, if \(v\) satisfies those restrictions, \(x+tv\)
satisfies the active inequalities and the equalities for all
\(t\ge0\). Each inactive inequality has a positive slack at
\(t=0\), so finitely many such inequalities all hold for
sufficiently small \(t>0\). This realizes \(v\) as a tangent.

For convex \(P\), \(N_P(x)=T_P(x)^\circ\). The normal inequality
passes to tangent limits, and every \(y-x\) with \(y\in P\)
is tangent along the segment from \(x\) to \(y\).
Let \(K\) be the cone generated by the active \(a_i\) and
the vectors \(e_j,-e_j\). Then the formula for \(T_P(x)\)
gives \(T_P(x)=K^\circ\). The cone \(K\) is closed by
Lemma~\ref{lem:finitely-generated-closed}, and
Theorem~\ref{thm:cone-bipolar} gives \(N_P(x)=K\).
Regroup the coefficients of \(e_j,-e_j\) into the free \(\nu_j\).
\end{proof}

This exact normal formula is used below for polyhedral perturbations,
LP and convex QP. See \cite[\S22]{Rockafellar} and
\cite[\S\S12.3 and 13.2]{NocedalWright}.
\chapter{Convex Functions, Conjugates, and Subgradients}

\section{Epigraphs, lower semicontinuity, and convex operations}
\label{v02-01}

\begin{definition}
A proper function \(f:X\to(-\infty,+\infty]\) is finite somewhere.
Its effective domain, epigraph, and real sublevel sets are
\[
 \operatorname{dom}f=\{x:f(x)<+\infty\},\qquad
 \operatorname{epi} f=\{(x,t)\in X\times\mathbb R:f(x)\le t\},
\]
with sublevel set \(\{x:f(x)\le a\}\) at each \(a\in\mathbb R\).
Convexity means that the epigraph is convex. Lower semicontinuity
means that \(\{x:f(x)>a\}\) is open for every real \(a\).
\end{definition}

\begin{proposition}[Epigraph criteria]\label{prop:epigraph-criteria}
Convexity is equivalent to the two-point convexity inequality
for points in the effective domain. It implies finite Jensen
inequalities. Lower semicontinuity is equivalent both to
closed real sublevels and to a closed epigraph.
\end{proposition}

\begin{proof}
Applying epigraph convexity to \((x,f(x))\) and \((y,f(y))\)
gives the two-point inequality. Conversely the inequality
for finite endpoints implies convexity for points of the
epigraph with any real heights above the endpoint values.
Inducting on the number of positive weights proves
\[
 f\Big(\sum_{i=1}^m\lambda_i x_i\Big)
 \le \sum_{i=1}^m\lambda_i f(x_i),
 \qquad \lambda_i>0,\quad \sum_i\lambda_i=1 .
\]
Weights equal to zero can be discarded. If a positive weight
has infinite function value, the inequality holds with infinite
right-hand side.

The sublevel formulation is the complement of the defining
open sets. If the sublevels are closed and \((x,t)\) lies
outside the epigraph, choose \(r\) real with \(t<r<f(x)\).
The product \(\{f>r\}\times(-\infty,r)\) is an open
neighborhood disjoint from the epigraph. Conversely a real
sublevel is the inverse image of the epigraph under the
continuous map \(x\mapsto(x,a)\).
\end{proof}

\begin{proposition}[Convex operations]\label{prop:convex-operations}
Positive weighted sums, pointwise suprema, and pullback by an
affine map preserve convexity. The perspective
\((x,t)\mapsto t f(x/t)\) on \(t>0\), extended by \(+\infty\)
on \(t\le0\), is convex. If \(F\) is convex on \(X\times Y\)
and \(v(x)=\inf_y F(x,y)\) is proper, then \(v\) is convex.
Positive sums, suprema and continuous affine pullback preserve
lower semicontinuity for proper lower-semicontinuous inputs.
The open perspective and inf-projection can lose lower semicontinuity.
\end{proposition}

\begin{proof}
At points where the result is finite, sum the defining
inequalities with positive weights, take their supremum, or
use affinity to pull them back. For a convex combination of
\((x_i,t_i)\) with \(t_i>0\), writing
\(t=(1-s)t_1+st_2\) gives
\[
 \frac{(1-s)x_1+sx_2}{t}
  =\frac{(1-s)t_1}{t}\frac{x_1}{t_1}
    +\frac{st_2}{t}\frac{x_2}{t_2}.
\]
Apply convexity and multiply by \(t\).
For inf-projection, at finite endpoint values choose \(y_i\)
giving \(F(x_i,y_i)<v(x_i)+\epsilon\). The convex combination
of the two \(y_i\) witnesses the convex bound with error
\(\epsilon\) for the projected function. Let \(\epsilon\downarrow0\).

For lower semicontinuity of a sum, if the sum at \(x\) exceeds
\(a\), choose real lower thresholds for the finitely many
summands at \(x\) whose positive weighted sum exceeds \(a\).
Intersect the corresponding open strict upper-level sets.
For a supremum, a strict upper-level set is the union of
those of its summands, and affine pullback uses inverse images.

The perspective of the constant zero function has value \(0\)
for \(t>0\) and \(+\infty\) at \(t=0\). For inf-projection let
\(F\) be the indicator of the set
\(\{(x,y)\in\mathbb R^2:x>0,\ y\ge1/x\}\).
Here the indicator of a set is \(0\) on it and \(+\infty\)
outside it. The set is convex by convexity of \(1/x\) for
\(x>0\), and closed because \(1/x\to+\infty\) as \(x\downarrow0\).
The projected indicator is that of \((0,\infty)\), which
fails lower semicontinuity at zero.
\end{proof}

See \cite[\S\S4--5 and 7]{Rockafellar} and
\cite[Chapter 3]{Boyd}.

\section{Continuity on relative interiors}\label{v02-02}

\begin{theorem}[Upper bounds and local Lipschitz continuity]
\label{thm:convex-local-lipschitz}
A proper convex function on a normed space, bounded above on
a neighborhood of \(a\), is finite and Lipschitz on a smaller
neighborhood of \(a\). In finite dimension a proper convex
function is locally Lipschitz on \(\operatorname{ri}(\operatorname{dom}f)\)
for the relative topology of its affine hull.
\end{theorem}

\begin{proof}
Suppose \(f\le M\) on \(B(a,R)\). Reflection through \(a\)
and the midpoint inequality give
\(f(x)\ge 2f(a)-M=:m\) there. If \(x,y\in B(a,R/4)\)
are distinct, put \(t=\lVert y-x\rVert<R/2\) and
\(z=x+(R/2)(y-x)/t\in B(a,R)\).
Since \(y=(1-2t/R)x+(2t/R)z\), convexity gives
\(f(y)-f(x)\le 2t(M-m)/R\). Swapping \(x,y\)
proves the Lipschitz bound.

For \(a\) in the relative interior of the effective domain,
work in its affine hull. In positive dimension a small cube
around \(a\) has all vertices in the effective domain. Its
convex combinations give a finite upper bound throughout
the cube. Apply the first part in the direction space of the
affine hull. In dimension zero the claim concerns a single
point.
\end{proof}

The argument at a relative-interior point uses only finite
values in the affine hull. See \cite[Theorem 10.1]{Rockafellar}
and \cite[Chapter I, \S2]{EkelandTemam}.

\section{Weak lower semicontinuity and Mazur's lemma}
\label{v02-03}

\begin{theorem}[Weak closure and convex approximation]
\label{thm:convex-weak-closed}
In a real normed space a norm-closed convex set is weakly
closed, and a norm-lower-semicontinuous convex function is
weakly lower semicontinuous. If \(x_n\rightharpoonup x\),
there are finite convex combinations \(z_n\) of members of
\(\{x_k:k\ge n\}\) with \(z_n\to x\) in norm.
\end{theorem}

\begin{proof}
The weak topology is the topology generated by all continuous
linear functionals. Each of their norm-closed halfspaces is
weakly closed. Theorem~\ref{thm:cone-bipolar} expresses any
norm-closed convex set as an intersection of such halfspaces.
Apply this to the norm-closed convex sublevels to obtain
weak lower semicontinuity.

Let \(C_n\) be the norm closure of the convex hull of the tail
from \(x_n\). It is convex because taking closures preserves
convexity by continuity of the vector operations. It is weakly
closed by the first assertion and contains the weak limit \(x\)
because it contains every member of the tail.
Choose a finite convex combination \(z_n\) in that tail with
\(\lVert z_n-x\rVert<1/n\).
\end{proof}

See \cite[Theorem 3.7 and Corollary 3.9]{Brezis}.

\section{Conjugates and the Fenchel--Young inequality}\label{v02-04}

\begin{definition}
For a proper \(f\) on a real locally convex space \(X\), its
conjugate on the continuous dual is
\[
 f^*(y)=\sup_{x\in\operatorname{dom}f}
             \bigl(\langle y,x\rangle-f(x)\bigr).
\]
For \(C\subset X\), write \(\delta_C\) for its extended-real
indicator, \(0\) on \(C\) and \(+\infty\) elsewhere.
\end{definition}

\begin{proposition}[Conjugacy and examples]\label{prop:conjugate-basic}
The conjugate is convex, is lower semicontinuous for the weak-star
topology \(\sigma(X',X)\), and satisfies
\(f(x)+f^*(y)\ge\langle y,x\rangle\).
For a nonempty \(C\), \(\delta_C^*=\sigma_C\).
If \(C\) is also closed convex, then \(\sigma_C^*=\delta_C\) in the
same pairing. On a normed space the conjugate of \(x\mapsto\lVert
x\rVert\) is the indicator of the closed dual unit ball.
For a real symmetric positive semidefinite matrix \(Q\), the
conjugate of \(q(x)=x^{\mathsf T}Qx/2\) is
\[
 q^*(y)=
 \begin{cases}
  y^{\mathsf T}Q^\dagger y/2,& y\in\operatorname{ran}Q,\\
  +\infty,&y\notin\operatorname{ran}Q,
 \end{cases}
\]
where \(Q^\dagger\) is the inverse on \((\ker Q)^\perp\) and
zero on \(\ker Q\).
\end{proposition}

\begin{proof}
The defining supremum is over continuous affine functions of
\(y\) for the weak-star topology, giving convexity and lower
semicontinuity. A finite point of \(f\) supplies a finite
lower bound for this supremum at each \(y\). Using \(x\) as a
competitor proves Fenchel--Young when \(f(x)\) is finite, and
the remaining case holds with infinite left-hand side.
For the indicator the defining supremum is \(\sigma_C\).
For \(x\in C\) one has
\(\langle y,x\rangle\le\sigma_C(y)\) for every \(y\), and
\(y=0\) gives equality. For \(x\notin C\) strict separation
gives \(y\) with \(\langle y,x\rangle>\sigma_C(y)\).
Scaling \(y\) makes the supremum infinite.

If \(\lVert y\rVert_*\le1\), the pairing is at most
\(\lVert x\rVert\) and \(x=0\) gives conjugate value zero.
For \(\lVert y\rVert_*>1\), choose \(x\) with
\(\langle y,x\rangle>\lVert x\rVert\), reversing its sign if
needed, and scale \(x\).
For the quadratic, symmetry gives
\(\operatorname{ran}Q=(\ker Q)^\perp\). If \(y\) has a
nonzero pairing on the kernel, scale a kernel vector with
positive pairing to get value \(+\infty\). Otherwise solve
\(Qz=y\) with \(z=Q^\dagger y\), and write
\[
 \langle y,x\rangle-\tfrac12x^{\mathsf T}Qx
 =\tfrac12z^{\mathsf T}Qz-\tfrac12(x-z)^{\mathsf T}Q(x-z).
\]
Its maximum is the desired value, attained at \(z\).
\end{proof}

See \cite[\S\S12--16]{Rockafellar}, \cite[\S3.3]{Boyd},
and \cite[Chapter I, \S\S3--4]{EkelandTemam}.

\section{Fenchel--Moreau and vertical separation}\label{v02-05}

The biconjugate in the continuous-dual pairing is
\(f^{**}(x)=\sup_{y\in X'}(\langle y,x\rangle-f^*(y))\).
For this supremum a term with \(f^*(y)=+\infty\) is \(-\infty\).

\begin{theorem}[Fenchel--Moreau]\label{thm:fenchel-moreau}
If \(X\) is real Hausdorff locally convex and \(f\) is proper
lower-semicontinuous convex, then \(f=f^{**}\). Equivalently
\(f\) is the pointwise supremum of its continuous affine
minorants. At least one such minorant exists.
\end{theorem}

\begin{proof}
Let \(x_0\) be a finite point of \(f\), and take a real
\(r_0<f(x_0)\). In \(X\times\mathbb R\) strictly separate
\((x_0,r_0)\) from the closed nonempty convex epigraph by
Theorem~\ref{thm:geometric-separation}. A continuous linear
functional on this product has the form
\((z,t)\mapsto l(z)+s t\). Orient the separation so that
\[
 l(x_0)+s r_0<\alpha\le l(z)+s t
       \qquad ((z,t)\in\operatorname{epi} f).
\]
Since the epigraph contains arbitrarily large \(t\) over
\(x_0\), one must have \(s\ge0\). Comparing to the point
\((x_0,f(x_0))\) shows \(s>0\). Hence
\(m_0(z)=(\alpha-l(z))/s\) is a continuous affine minorant.

Now take any \(x\) and real \(r<f(x)\). Separate \((x,r)\)
strictly in the same orientation. If the height coefficient
\(s>0\), the resulting affine minorant has value greater
than \(r\) at \(x\). If \(s=0\), separation gives
\(l(x)<\alpha\le l(z)\) for all \(z\in\operatorname{dom}f\).
For each \(k>0\),
\[
 m_k(z)=m_0(z)+k(\alpha-l(z))
\]
is still a minorant, while \(m_k(x)\to+\infty\).
Thus the supremum of the minorants is \(f\), including at
infinite points. If \(m(z)=\langle y,z\rangle+b\le f(z)\),
then \(f^*(y)\le-b\). It follows that
\(f^{**}(x)\ge\langle y,x\rangle+b\) for each such \(m\).
The reverse inequality is Fenchel--Young.
\end{proof}

In particular a proper lower-semicontinuous convex function
has proper conjugate. The argument pairs \(X\) directly
with \(X'\), with no reflexivity hypothesis. See
\cite[Chapter I, \S\S3--4]{EkelandTemam} and
\cite[Theorems 12.1--12.2]{Rockafellar}.

\section{Subgradients and Fenchel equality}\label{v02-06}

\begin{definition}
At \(x\in\operatorname{dom}f\), the subdifferential
\(\partial f(x)\) consists of all \(y\in X'\) satisfying
\[
 f(z)\ge f(x)+\langle y,z-x\rangle\qquad(z\in X).
\]
At other points it is empty. For a function on \(X'\) the
subdifferential in the pairing uses test elements \(x\in X\),
acting by evaluation on \(X'\).
\end{definition}

\begin{theorem}[Stationarity and equality]\label{thm:fenchel-equality}
For a proper convex function, \(x\in\operatorname{dom}f\)
minimizes \(f\) if and only if \(0\in\partial f(x)\).
For \(x\in C\) and \(C\) nonempty convex,
\(\partial\delta_C(x)=N_C(x)\), where
\(N_C(x)=\{y\in X':\langle y,z-x\rangle\le0\ \text{for all }z\in C\}\).
For \(x\in\operatorname{dom}f\),
\[
 y\in\partial f(x)\quad\Longleftrightarrow\quad
 f(x)+f^*(y)=\langle y,x\rangle.
\]
If \(f\) is lower semicontinuous on a Hausdorff locally convex
space, these equivalent statements are also equivalent to
\(x\in\partial f^*(y)\) in the pairing.
\end{theorem}

\begin{proof}
The first assertion is the support inequality with zero
slope. For the indicator, points outside \(C\) impose no
restriction, and points inside impose the normal inequality.
A subgradient satisfies
\(\langle y,z\rangle-f(z)\le\langle y,x\rangle-f(x)\)
for all finite \(z\), with equality at \(z=x\). This
characterizes equality in the defining supremum, and the
argument is reversible. Applying the same characterization
to \(f^*\) exchanges \(x,y\) and replaces \(f\) by \(f^{**}\).
Theorem~\ref{thm:fenchel-moreau} identifies \(f^{**}=f\).
\end{proof}

See \cite[Theorem 23.5]{Rockafellar} and
\cite[Chapter I, \S5]{EkelandTemam}.

\section{Directional derivatives, differentiability, and monotonicity}
\label{v02-07}

\begin{proposition}[Directional quotients and support]
\label{prop:directional-support}
For a proper convex function and a finite point \(x\), define
\[
 f'(x;d)=\lim_{t\downarrow0}
     \frac{f(x+td)-f(x)}{t}
    =\inf_{t>0}\frac{f(x+td)-f(x)}{t}.
\]
This limit is understood in the extended reals. If \(f\) is
finite and continuous in a neighborhood of \(x\) in a real
normed space, it is a finite continuous sublinear function of
\(d\), and
\[
 \partial f(x)=\{y\in X':\langle y,d\rangle\le f'(x;d)
                              \text{ for all }d\},
 \qquad
 f'(x;d)=\max_{y\in\partial f(x)}\langle y,d\rangle .
\]
In finite dimension \(\partial f(x)\) is nonempty at each
relative-interior point of the effective domain.
\end{proposition}

\begin{proof}
For \(0<t<s\), the point \(x+td\) lies on the segment from
\(x\) to \(x+sd\). Convexity makes the difference quotient
nondecreasing in \(t\). If \(f\) is continuous near \(x\),
Theorem~\ref{thm:convex-local-lipschitz} yields a local
Lipschitz constant \(L\), giving \(|f'(x;d)|\le L\lVert
d\rVert\). Positive homogeneity follows by rescaling \(t\).
The midpoint inequality at \(x+2td_1,x+2td_2\) gives
subadditivity on taking the limit. These estimates imply
continuity by bounding each difference of two directional
values using the direction difference.
A subgradient lies below every directional quotient
and hence lies below their infima. Conversely, a continuous
linear functional below those infima lies below the
quotients at \(t=1\), giving the subgradient inequality.

For \(d\neq0\), define a linear functional on the span of
\(d\) with value \(f'(x;d)\) at \(d\). The bound at negative
multiples holds since sublinearity gives
\(f'(x;d)+f'(x;-d)\ge0\).
Hahn--Banach extends it dominated by \(f'(x;\cdot)\), and the
Lipschitz bound on this sublinear function implies continuity
of the extension. This gives a subgradient attaining the
maximum for direction \(d\). For \(d=0\), a dominated
extension from the zero subspace gives a subgradient and
attainment. At a finite-dimensional relative-interior point,
apply this argument in the direction space of the affine
hull using Theorem~\ref{thm:convex-local-lipschitz}. Extend
the resulting linear functional to the ambient space.
The inequality off that affine hull holds with \(+\infty\)
on its left, so the extension remains a subgradient.
\end{proof}

\begin{corollary}[Smooth support and monotonicity]
\label{cor:convex-differentiability}
At an interior point where a convex function is Fr\'echet
differentiable, its derivative is the unique subgradient.
In finite dimension a proper convex function with a singleton
subdifferential at an interior point of its effective domain
is Fr\'echet differentiable there. Every convex subdifferential
is monotone:
\[
 y_i\in\partial f(x_i)\ (i=1,2)
 \quad\Longrightarrow\quad
 \langle y_1-y_2,x_1-x_2\rangle\ge0 .
\]
For a differentiable function on an open convex set,
convexity is equivalent to the gradient supporting inequality
throughout that set.
\end{corollary}

\begin{proof}
Differentiability makes the directional derivative linear,
so the functionals lying below it in every direction equal
the derivative by applying the inequality to both \(d\) and
\(-d\). Conversely, the directional support formula makes
every directional derivative equal to the singleton slope
\(y\). Suppose Fr\'echet differentiability fails and take
nonzero \(h_n\to0\) with normalized linearization errors
bounded away from zero in absolute value. Along a subsequence
\(d_n=h_n/\lVert h_n\rVert\) tends to a unit vector \(d\).
Local Lipschitz continuity bounds the difference between the
directional quotients for \(d_n,d\) at \(t_n=\lVert h_n\rVert\)
by \(L\lVert d_n-d\rVert\).
The quotient for \(d\) converges to \(\langle y,d\rangle\),
and \(y(d_n)\to y(d)\), a contradiction.

Adding the support inequalities at \(x_1,x_2\) gives
monotonicity. If the gradient is supporting at every point
of an open convex set, apply its support at
\(z=(1-t)x+ty\) to both \(x\) and \(y\). The weighted sum
of the linear terms vanishes, giving the convexity inequality.
The converse is the unique-subgradient assertion.
\end{proof}

At boundary points the directional derivative can be infinite.
For example \(f(t)=-\sqrt t\) on \(t\ge0\), extended by
\(+\infty\) otherwise, has \(f'(0;1)=-\infty\), and the
support inequality gives an empty subdifferential at zero.
See \cite[Theorems 23.2, 23.4 and 25.1]{Rockafellar}
and \cite[Chapter I, \S5]{EkelandTemam}.

\section{Proximal maps and the Moreau envelope}\label{v02-08}

Let \(H\) be a real Hilbert space and identify its continuous
dual with \(H\) by the inner product. For proper lower-semicontinuous
convex \(f\) and \(\lambda>0\) set
\[
 e_\lambda f(x)
  =\inf_{u\in H}\left(f(u)+\frac{\lVert u-x\rVert^2}{2\lambda}\right).
\]
The proximal map \(P_\lambda^f(x)\) will denote the unique
minimizer in this formula.

\begin{theorem}[Proximal existence and characterization]
\label{thm:prox-characterization}
The envelope is finite everywhere and its infimum is attained
at a unique point \(p=P_\lambda^f(x)\). For \(p\in\operatorname{dom}f\),
\[
 p=P_\lambda^f(x)\quad\Longleftrightarrow\quad
   (x-p)/\lambda\in\partial f(p).
\]
The map \(P=P_\lambda^f\) is firmly nonexpansive:
\(\lVert P x-P y\rVert^2\le (Px-Py,x-y)\).
Both \(P\) and \(I-P\) are nonexpansive.
\end{theorem}

\begin{proof}
Theorem~\ref{thm:fenchel-moreau} supplies a continuous
affine lower bound for \(f\). Adding the positive quadratic
gives a finite lower bound for the minimized expression,
while a finite point of \(f\) gives an upper bound.
Writing that expression as \(F_x\), convexity and the
quadratic identity give
\[
 F_x((u+v)/2)\le
 \tfrac12 F_x(u)+\tfrac12 F_x(v)
   -\frac{\lVert u-v\rVert^2}{8\lambda}.
\]
For a finite minimizing sequence this implies the Cauchy
property because the midpoint value is at least the
infimum. Completeness and lower semicontinuity yield a
minimizer, and the same inequality yields uniqueness.

At a minimizer \(p\), for finite \(z\) and \(0<t\le1\),
minimality along \(p+t(z-p)\) and convexity of \(f\) yield
\[
 0\le t(f(z)-f(p))
       +\frac{t}{\lambda}(p-x,z-p)
       +\frac{t^2}{2\lambda}\lVert z-p\rVert^2.
\]
Divide by \(t\) and let \(t\downarrow0\) to obtain the
subgradient inequality for \((x-p)/\lambda\).
Conversely this inequality and expansion of the square
give \(F_x(z)\ge F_x(p)+\lVert z-p\rVert^2/(2\lambda)\).

For \(p=Px,\ q=Py\), apply monotonicity of \(\partial f\)
to the two proximal subgradients. Multiplying by \(\lambda\)
gives
\((x-y,p-q)\ge\lVert p-q\rVert^2\).
Cauchy--Schwarz gives nonexpansiveness of \(P\).
Expanding
\(\lVert(x-p)-(y-q)\rVert^2\) and using the same inequality
gives the bound by \(\lVert x-y\rVert^2\) for \(I-P\).
\end{proof}

\begin{theorem}[Envelope gradient and Moreau decomposition]
\label{thm:moreau}
The Moreau envelope \(e_\lambda f\) is convex and Fr\'echet differentiable,
with gradient \((x-P_\lambda^f(x))/\lambda\), which is
\(1/\lambda\)-Lipschitz. With the same proximal-map convention,
\[
 x=P_\lambda^f(x)
      +\lambda P_{1/\lambda}^{f^*}(x/\lambda),
 \qquad
 e_\lambda f(x)+e_{1/\lambda}f^*(x/\lambda)
     =\frac{\lVert x\rVert^2}{2\lambda}.
\]
\end{theorem}

\begin{proof}
Joint convexity of \(f(u)+\lVert u-x\rVert^2/(2\lambda)\)
and Proposition~\ref{prop:convex-operations} give convexity
of the finite inf-projection. Put \(p=P_\lambda^f(x)\) and
\(p_h=P_\lambda^f(x+h)\).
Testing the envelope at \(x+h\) with \(p\), and at \(x\)
with \(p_h\), gives
\[
 \frac{(x+h-p_h,h)}{\lambda}
      -\frac{\lVert h\rVert^2}{2\lambda}
 \le e_\lambda f(x+h)-e_\lambda f(x)
 \le \frac{(x-p,h)}{\lambda}
      +\frac{\lVert h\rVert^2}{2\lambda}.
\]
Since \(I-P_\lambda^f\) is nonexpansive,
\(\lVert(x+h-p_h)-(x-p)\rVert\le\lVert h\rVert\).
Both bounds establish the asserted derivative and its
Lipschitz constant.

Put \(q=(x-p)/\lambda\).
The proximal characterization and
Theorem~\ref{thm:fenchel-equality} give \(p\in\partial f^*(q)\).
Here \(f^*\) is proper lower-semicontinuous convex on \(H\):
weak-star lower semicontinuity implies norm lower
semicontinuity under the Riesz identification.
As \(p=\lambda(x/\lambda-q)\), applying the proximal
characterization to \(f^*\) at \(x/\lambda\) gives
\(q=P_{1/\lambda}^{f^*}(x/\lambda)\) and the decomposition.
Finally, the two minimizing expressions sum to
\[
 f(p)+f^*(q)+\frac{\lambda}{2}\lVert q\rVert^2
        +\frac{\lVert p\rVert^2}{2\lambda}
 = (p,q)+\frac{\lambda}{2}\lVert q\rVert^2
        +\frac{\lVert p\rVert^2}{2\lambda}
 = \frac{\lVert x\rVert^2}{2\lambda},
\]
by Fenchel equality and \(p+\lambda q=x\).
\end{proof}

The argument uses the supporting and conjugacy principles of
\cite[\S\S12 and 23--24]{Rockafellar} and
\cite[Chapter I, \S5]{EkelandTemam}. For an indicator of a nonempty
closed convex set the proximal map reduces to nearest-point projection.
\chapter{Minimization, Weak Compactness, and Hilbert-Space Problems}
\chaptermark{Weak Compactness and Hilbert-Space Problems}

\section{Weakly convergent subsequences in reflexive spaces}
\label{v03-01}

For a real normed space \(X\), write \(X'\) for its continuous dual.
The weak topology on \(X\) uses evaluations by \(X'\), whereas the weak-star
topology on \(X'\) uses evaluations at points of \(X\). The canonical map
\(J_X:X\to X''\) is isometric by Hahn--Banach. Reflexivity means that this
map is onto. In particular the weak and weak-star topologies on the dual
of a reflexive space agree.

\begin{theorem}[Weak subsequences in a reflexive Banach space]
\label{thm:reflexive-weak-subsequence}
Every bounded sequence in a real reflexive Banach space \(X\) has a weakly
convergent subsequence. Also each closed norm ball is weakly compact.
\end{theorem}
\begin{proof}
For weak compactness, \(J_X\) identifies a closed ball centered at zero
with the corresponding ball of \(X''\). Under this map the weak topology
on \(X\) corresponds to the weak-star topology of \(X''\) as the dual
of \(X'\). Apply Banach--Alaoglu, and use translations for other centers.

Let \((x_n)\) be bounded and put
\(Y=\overline{\operatorname{span}}\{x_n:n\ge1\}\). Rational linear
combinations show that \(Y\) is separable. To see it is reflexive, let
\(F\in Y''\) and let \(R:X'\to Y'\) be restriction. Hahn--Banach
makes \(R\) surjective. Since \(X\) is reflexive, \(F\circ R\) is evaluation
at some \(x\in X\). Every functional vanishing on \(Y\) vanishes at \(x\),
so separation of the closed subspace \(Y\) gives \(x\in Y\).
Consequently \(F\) is evaluation at \(x\) on \(Y'\).

The dual unit ball \(B\subset Y'\) is weak-star compact by Banach--Alaoglu.
For a countable dense sequence \((y_j)\) in \(Y\), the metric
\[
 d(f,g)=\sum_{j\ge1}2^{-j}\min\{1,|(f-g)(y_j)|\}
\]
induces the weak-star topology on \(B\). Indeed convergence in this metric
is convergence at every \(y_j\), and the bound \(\|f-g\|\le2\) extends
the evaluation control to every \(y\in Y\). Conversely a finite number of
evaluations controls this metric up to an arbitrarily small tail.
Thus \(B\) is compact metric and has a countable weak-star dense subset
\(D\). Reflexivity of \(Y\) makes \(D\) weakly dense in \(B\).
The norm-closed convex hull of \(D\) is weakly closed by separation.
Since it contains \(D\), it contains \(B\), and it is contained in \(B\).
Hence rational convex combinations of points of \(D\) are norm dense in
\(B\).

Take a diagonal subsequence of \((x_n)\) on which evaluation by every
one of these countably many functionals converges, using scalar boundedness
at each extraction. The uniform bound on \(\|x_n\|\) and norm density
make evaluation along this subsequence Cauchy for every \(f\in Y'\)
after rescaling into \(B\). The limits define a bounded linear functional
on \(Y'\), represented by a point \(y\in Y\) by reflexivity. This is weak
convergence in \(Y\). It is weak convergence in \(X\) as well, since every
functional in \(X'\) restricts to one in \(Y'\).
\end{proof}

This argument distinguishes weak-star compactness from sequential
compactness of weakly bounded sequences. See
\cite[Theorems 3.16--3.18 and Section 3.6]{Brezis}.

\section{The direct method and its hypotheses}
\label{v03-02}

An extended-real objective is proper if it is never \(-\infty\) and has
a finite value somewhere. A function is coercive on a subset \(C\) of a
normed space if each real sublevel in \(C\) is bounded. A set is weakly
sequentially closed if it contains the limits of its weakly convergent
sequences. Weak sequential lower semicontinuity of \(J\) on \(C\) means
\(J(x)\le\liminf_n J(x_n)\) for weakly convergent sequences in \(C\)
whose limit \(x\) belongs to \(C\).

\begin{theorem}[Three forms of the direct method]
\label{thm:direct-method}
Let \(J\) be proper on a nonempty set \(C\).
\begin{enumerate}
\item If \(C\) is compact and \(J\) is lower semicontinuous, then \(J\)
attains its infimum. More generally it suffices that for one
finite-energy point \(x_0\), the sublevel
\(\{x\in C:J(x)\le J(x_0)\}\) is compact and \(J\) is lower
semicontinuous on it.
\item In finite dimension, a lower semicontinuous coercive \(J\) on a
closed \(C\) attains its infimum.
\item In a reflexive real Banach space, suppose \(C\) is weakly
sequentially closed and \(J\) is bounded below, coercive and weakly
sequentially lower semicontinuous. Then \(J\) attains its infimum.
If for some \(\alpha\in\R\) the sublevel \(\{x\in C:J(x)\le\alpha\}\)
is nonempty and weakly compact, and \(J\) is weakly lower semicontinuous
on it, the conclusion follows from the first assertion in the weak topology.
\end{enumerate}
\end{theorem}
\begin{proof}
For the first assertion it suffices to minimize on the indicated sublevel
\(K\), since points outside it have strictly greater value than \(J(x_0)\).
Put \(m=\inf_K J\). If \(m>-\infty\), the closed subsets
\(\{x\in K:J(x)\le m+1/n\}\) are nonempty and nested. Compactness yields a
point of value \(m\). If \(m=-\infty\), use the nonempty closed subsets
with thresholds \(-n\) to obtain a point where \(J=-\infty\), contradicting
properness. For the second assertion the finite sublevel is closed in the
ambient space and bounded, hence compact. For the third, choose a
minimizing sequence in a bounded sublevel, take a weakly convergent
subsequence by Theorem~\ref{thm:reflexive-weak-subsequence}, and use
sequential closedness and \(J(x)\le\liminf_n J(x_n)\) at the weak limit.
The weakly compact version is the first assertion in the weak topology.
\end{proof}

\begin{example}[Norm continuity alone]
On the closed unit ball of \(\ell^2(\mathbb N;\mathbb R)\) let
\[
 J(x)=1-\|x\|^2+\sum_{j\ge1}\frac{x_j^2}{j}.
\]
The series term is the quadratic form of a bounded diagonal operator,
so \(J\) is norm continuous. It is positive at every point, including
the sphere, where the series is positive. But \(J(e_n)=1/n\), so the
infimum is zero and is not attained. Since \(e_n\) converges weakly to
zero and \(J(0)=1\), the missing hypothesis is weak lower
semicontinuity.
\end{example}

See \cite[Chapter II, Section 1]{EkelandTemam},
\cite[Section 27]{Rockafellar}, and \cite[Chapter 3]{Dacorogna}.

\section{Convex minimizers and quantitative strong convexity}
\label{v03-03}

\begin{proposition}[Local minima, uniqueness, and first-order testing]
\label{prop:v03-03-01}
Let \(C\) be convex and \(J:C\to(-\infty,\infty]\) be proper and convex.
Every finite local minimizer of \(J\) on \(C\) is a global minimizer.
If \(J\) is strictly convex between distinct finite-value points, it has
at most one minimizer. If \(J\) extends to a differentiable convex
function on an open convex neighborhood of \(C\) in a Euclidean space,
then \(x_*\in C\) minimizes \(J\) on \(C\) exactly when
\[
 \langle \nabla J(x_*),y-x_*\rangle\ge0 \qquad(y\in C).
\]
\end{proposition}
\begin{proof}
If \(J(y)<J(x_*)\), each sufficiently short segment from \(x_*\)
toward \(y\) remains feasible and has smaller value by convexity.
Two distinct minimizers would give a midpoint of strictly smaller value
under strict convexity. At a differentiable minimum, the right
derivative of \(t\mapsto J(x_*+t(y-x_*))\) is nonnegative at \(0\).
Conversely, monotonicity of secant slopes for this one-dimensional
convex function gives
\(J(y)-J(x_*)\ge\langle\nabla J(x_*),y-x_*\rangle\).
\end{proof}

\begin{theorem}[Strong-convexity bounds]
\label{thm:v03-03-02}
Let \(J:\mathbb R^d\to\mathbb R\) be differentiable and
\(\mu\)-strongly convex for \(\mu>0\), meaning
\(J-\mu\|\cdot\|^2/2\) is convex. If \(x_*\) minimizes \(J\), then
\[
 \frac{\mu}{2}\|x-x_*\|^2
 \ \le\ J(x)-J(x_*)
 \ \le\ \frac{\|\nabla J(x)\|^2}{2\mu}\qquad(x\in\mathbb R^d).
\]
\end{theorem}
\begin{proof}
Apply the supporting inequality to \(J-\mu\|\cdot\|^2/2\) to obtain
\[
 J(y)\ge J(x)+\langle\nabla J(x),y-x\rangle
                  +\frac{\mu}{2}\|y-x\|^2.
\]
With the support point \(x_*\), stationarity gives the lower bound.
With support point \(x\), the right side has infimum
\(J(x)-\|\nabla J(x)\|^2/(2\mu)\) over \(y\), by completing the square.
In particular it has this lower bound at \(y=x_*\).
\end{proof}

See \cite[Section 27]{Rockafellar}, \cite[Sections 4.2 and 9.1]{Boyd},
and \cite[Theorem 2.5]{NocedalWright}.

\section{Ekeland's metric variational principle}
\label{v03-04}

\begin{theorem}[Ekeland]
\label{thm:v03-04-01}
Let \((M,d)\) be a complete metric space, and let
\(f:M\to(-\infty,\infty]\) be proper, lower semicontinuous and
bounded below. For \(\epsilon,\lambda>0\) suppose
\(f(x_0)\le\inf_M f+\epsilon\). There exists \(x_*\in M\) such that
\[
 f(x_*)\le f(x_0),\qquad d(x_0,x_*)\le\lambda,\qquad
 f(y)>f(x_*)-\frac{\epsilon}{\lambda}d(y,x_*)\quad(y\ne x_*).
\]
\end{theorem}
\begin{proof}
Set \(\alpha=\epsilon/\lambda\). For finite-valued \(x\) let
\(S(x)=\{y:f(y)+\alpha d(y,x)\le f(x)\}\).
It is nonempty and closed. The triangle inequality shows that
\(S(y)\subset S(x)\) whenever \(y\in S(x)\).
Choose \(x_{n+1}\in S(x_n)\) with
\(f(x_{n+1})\le\inf_{S(x_n)} f+2^{-n}\).
Then
\[
 \alpha\sum_{j=n}^{k-1}d(x_{j+1},x_j)
 \le f(x_n)-f(x_k).
\]
The decreasing finite values converge to some \(m\).
Thus \((x_n)\) is Cauchy, with limit \(x_*\). The nested closed sets
give \(x_*\in S(x_n)\) for all \(n\). Membership in \(S(x_0)\) gives
the value bound and
\(\alpha d(x_0,x_*)\le f(x_0)-f(x_*)\le\epsilon\).
For \(y\in S(x_*)\) we have \(y\in S(x_n)\) for all \(n\), and
hence \(f(y)\ge f(x_{n+1})-2^{-n}\). Passing to the limit gives
\(f(y)\ge m\). This also applies to \(x_*\), while lower
semicontinuity gives \(f(x_*)\le m\). The definition of \(S(x_*)\)
now forces \(d(y,x_*)=0\). Therefore \(S(x_*)=\{x_*\}\).
\end{proof}

\begin{corollary}[Approximate stationarity]
\label{cor:v03-04-02}
If \(M=X\) is a Banach space, \(d\) is its norm metric, and \(f\) is
Fr\'echet differentiable at \(x_*\) from the theorem, then
\(\|Df(x_*)\|\le\epsilon/\lambda\).
\end{corollary}
\begin{proof}
For \(v\in X\), apply the last inequality to \(x_*+t v\)
for \(t>0\). Divide by \(t\) and pass to zero, obtaining
\(Df(x_*)v\ge-\alpha\|v\|\). Apply this to \(-v\) as well.
\end{proof}

See \cite[Chapter I, Section 6]{EkelandTemam}.

\section{Projection onto a closed convex set}
\label{v03-05}

\begin{theorem}[Hilbert projection]
\label{thm:hilbert-projection}
In a real Hilbert space \(H\), let \(C\) be nonempty, closed, and convex.
For every \(x\in H\) there is a unique nearest point \(P_C x\in C\).
The variational characterization is
\[
 p=P_Cx
 \quad\Longleftrightarrow\quad
 p\in C,\quad \langle x-p,y-p\rangle\le0\quad(y\in C).
\]
The projection is firmly nonexpansive:
\[
 \|P_Cx-P_Cz\|^2\le \langle P_Cx-P_Cz,x-z\rangle \qquad(x,z\in H).
\]
The complex Hilbert-space version uses the real part of the inner
product, and convexity is over real parameters.
\end{theorem}
\begin{proof}
Take \(p_n\in C\) with \(\|x-p_n\|\to d(x,C)\).
The midpoint of \(p_n,p_k\) is in \(C\), and the parallelogram law gives
\[
 \|p_n-p_k\|^2
 =2\|x-p_n\|^2+2\|x-p_k\|^2
       -4\left\|x-\frac{p_n+p_k}{2}\right\|^2.
\]
The last squared norm is at least \(d(x,C)^2\), so \(p_n\) is Cauchy.
Its limit is a nearest point in \(C\). Applying the identity to two
nearest points proves uniqueness. For a nearest point \(p\), differentiate
the squared distance at the feasible segment \(p+t(y-p)\) from the right.
This gives the displayed inequality. Conversely, expanding
\(\|x-y\|^2\) using \(x-y=(x-p)-(y-p)\) gives minimality from that
inequality. Test the characterizations at \(p=P_Cx\) and \(q=P_Cz\)
with \(q\) and \(p\), respectively. Adding yields
\(\langle x-z,p-q\rangle\ge\|p-q\|^2\).
Cauchy--Schwarz in particular makes \(P_C\) nonexpansive.
Apply the real proof to the underlying real Hilbert space for the
complex version.
\end{proof}

See \cite[Section 5.1]{Brezis}.

\section{Stampacchia's variational inequality}
\label{v03-06}

\begin{theorem}[Bounded coercive bilinear form]
\label{thm:stampacchia}
Let \(H\) be a real Hilbert space and \(C\subset H\) be nonempty, closed
and convex. Suppose \(a:H\times H\to\mathbb R\) is bilinear with
\[
 |a(u,v)|\le M\|u\|\|v\|,\qquad a(v,v)\ge \alpha\|v\|^2
 \quad(\alpha>0)
\]
and \(\ell\in H'\). There is a unique \(u\in C\) such that
\[
 a(u,v-u)\ge\ell(v-u)\qquad(v\in C).
\]
When \(a\) is symmetric, this inequality is equivalent to minimizing
\(\frac12 a(v,v)-\ell(v)\) on \(C\).
\end{theorem}
\begin{proof}
The zero-dimensional case is immediate. Otherwise \(M\ge\alpha>0\).
By Riesz representation write \(a(u,v)=\langle Au,v\rangle\) and
\(\ell(v)=\langle b,v\rangle\). The bounded linear operator \(A\)
satisfies \(\|A\|\le M\). Choose \(0<\tau<2\alpha/M^2\). For \(h\in H\),
\[
 \|(I-\tau A)h\|^2
 \le (1-2\tau\alpha+\tau^2M^2)\|h\|^2 .
\]
Thus \(u\mapsto P_C(u-\tau Au+\tau b)\) is a contraction on the
complete metric space \(C\) by Theorem~\ref{thm:hilbert-projection}.
Its unique fixed point satisfies exactly the asserted inequality by
the projection characterization. In the symmetric case, put
\(E(v)=\frac12 a(v,v)-\ell(v)\). For \(v\in C\),
\[
 E(v)-E(u)=a(u,v-u)-\ell(v-u)+\tfrac12 a(v-u,v-u).
\]
This proves minimality from the inequality. Conversely the feasible
segment from \(u\) toward \(v\) gives the inequality by taking the
right derivative of \(E\).
\end{proof}

See \cite[Theorem 5.6]{Brezis} and
\cite[Chapter II, Section 3]{EkelandTemam}.

\section{Lax--Milgram over the real and complex fields}
\label{v03-07}

\begin{theorem}[Lax--Milgram]
\label{thm:lax-milgram}
Let \(H\) be a real or complex Hilbert space. Let \(a\) be a bounded
bilinear form in the real case, or a bounded sesquilinear form,
conjugate-linear in the first slot and linear in the second, in the
complex case. Let \(\ell\) be a continuous linear functional in the
test variable. If
\(\operatorname{Re} a(v,v)\ge\alpha\|v\|^2\) for all \(v\) and
\(\alpha>0\), there is exactly one \(u\in H\) such that
\[
 a(u,v)=\ell(v)\quad(v\in H),\qquad
 \|u\|\le \frac{\|\ell\|}{\alpha}.
\]
For real symmetric or complex Hermitian \(a\), the solution is the
unique minimizer of
\(E(v)=\tfrac12 a(v,v)-\operatorname{Re}\ell(v)\).
\end{theorem}
\begin{proof}
Over \(\mathbb R\), take \(C=H\) in Theorem~\ref{thm:stampacchia}
and test with \(u+w\) and \(u-w\). Over \(\mathbb C\), apply this
real result to the underlying real Hilbert space, to
\(\operatorname{Re}a\) and \(\operatorname{Re}\ell\).
Testing the resulting equality at \(v\) and \(iv\) recovers the real
and imaginary parts of \(a(u,v)=\ell(v)\) because both expressions
are linear in \(v\). For the estimate test the equation with \(v=u\)
and use
\[
 \alpha\|u\|^2\le\operatorname{Re}\ell(u)\le\|\ell\|\|u\|.
\]
Applying the same bound to the difference of solutions proves
uniqueness. In the symmetric or Hermitian case,
\(E(u+w)-E(u)=\frac12 a(w,w)\) at a solution, which gives strict
minimality. Conversely the real derivative along \(w\), and also along
\(iw\) in the complex case, yields the variational equation.
\end{proof}

Coercivity alone provides the equation and the estimate. An energy
representation uses the symmetric part of a real form or the
Hermitian part of a complex form, which explains the additional
hypothesis for the energy conclusion. See
\cite[Corollary 5.8]{Brezis} and
\cite[Chapter II, Section 3]{EkelandTemam}.
\chapter{Perturbation Duality and Convex Constraints}

\section{Value functions and the dual bound}
\label{v04-01}

Let \(X,Y\) be real locally convex Hausdorff spaces. A proper
convex function \(\Phi\) on \(X\times Y\) defines the perturbation value
\[
 v(u)=\inf_{x\in X}\Phi(x,u),\qquad \inf\varnothing=+\infty .
\]
We write \(X'\) and \(Y'\) for the continuous duals, with the pairing
\(\langle y,u\rangle=y(u)\). To prove convexity of the marginal value,
take finite candidates at two perturbations and apply convexity to
the two pairs. Infimizing over the candidates gives the convexity
inequality when the endpoint values are finite. If either value is
\(-\infty\) and the other admits a finite candidate, choose arbitrarily
negative candidates at the former endpoint. Every interior convex
combination then has value \(-\infty\). This last observation is useful
for checking properness before applying duality.

\begin{theorem}[Perturbation duality]\label{thm:perturbation-duality}
Suppose \(v\) is proper and \(p=v(0)\in\R\). Then
\[
 v^*(y)=\Phi^*(0,y),\qquad
 d=\sup_{y\in Y'}[-\Phi^*(0,y)]=v^{**}(0)\le p.
\]
The equality \(d=p\) holds exactly when \(v\) is lower
semicontinuous at zero. A functional \(y\) attains the dual value
with \(d=p\) exactly when \(y\in\partial v(0)\).
\end{theorem}

\begin{proof}
For each \(u\) the supremum of \(\langle y,u\rangle-\Phi(x,u)\)
over \(x\) is \(\langle y,u\rangle-v(u)\). Taking the remaining
supremum over \(u\) gives the first formula. The second follows
by evaluating the double conjugate at zero. In particular,
\(v^*(y)\ge -p\). The supporting inequality
\(v(u)\ge p+\langle y,u\rangle\) holds for all \(u\) precisely
when \(v^*(y)=-p\).

The function \(v^{**}\) is lower semicontinuous as a supremum
of continuous affine functions and satisfies \(v^{**}\le v\).
Thus \(v^{**}(0)=v(0)\) implies lower semicontinuity of \(v\)
at zero. Conversely, suppose \(v\) is lower semicontinuous at
zero and fix \(r<p\). The point \((0,r)\) lies outside the
closed convex set \(\overline{\operatorname{epi}v}\).
Strictly separating it from that set, as in
Section~\ref{v01-04}, gives \(a\in Y'\), \(\alpha\in\R\)
and \(\gamma\in\R\) such that
\[
 \alpha r<\gamma\le a(u)+\alpha t
 \qquad((u,t)\in\operatorname{epi}v).
\]
The upward vertical ray above \((0,p)\) gives \(\alpha\ge0\),
and the point \((0,p)\) forces \(\alpha>0\).
Hence the continuous affine function \((\gamma-a(u))/\alpha\)
minorizes \(v\) and has value greater than \(r\) at zero.
Its value at zero bounds \(v^{**}(0)\) from below, by the
definition of the conjugate. Letting \(r\uparrow p\) proves
the claimed equality.
\end{proof}

\begin{lemma}[Local upper bound and supporting functional]\label{lem:local-value-support}
Let \(v\) be a convex marginal value on a real normed space.
Assume \(v(0)\in\R\), and \(v(u)\le M<\infty\) for
\(\|u\|<r\) with \(r>0\). Then \(v\) is proper, is
Lipschitz in a neighborhood of zero, and has a continuous
subgradient at zero.
\end{lemma}

\begin{proof}
If \(v(u)=-\infty\), choose \(\epsilon>0\) so that
\(-\epsilon u\) lies in the ball. It admits a finite candidate.
Convexity with these two endpoints forces \(v(0)=-\infty\).
Thus \(v\) is proper and finite in the ball. Apply
Theorem~\ref{thm:convex-local-lipschitz} to obtain the local
Lipschitz conclusion. Convexity at the midpoint of \(u,-u\)
also gives \(v(u)\ge 2v(0)-M\) in the ball.

For each \(w\) the slope
\[
 q(w)=\lim_{t\downarrow0}\frac{v(tw)-v(0)}{t}
\]
exists by monotonicity of the secant slopes. Applying the bounds
on a smaller ball to the rays \(\pm w\) gives
\(|q(w)|\le C\|w\|\) for a constant \(C\) independent of \(w\).
For example, take the slope at \(t=r/(2\|w\|)\) to bound
the slope from above, and use the opposite ray for the lower
bound. The equality \(q(sw)=s q(w)\) for \(s>0\) follows by
rescaling \(t\), and the midpoint inequality at \(2tw,2tz\)
gives \(q(w+z)\le q(w)+q(z)\). Hahn--Banach yields a linear
functional \(y\) with \(y(w)\le q(w)\) for every \(w\). Applying
this to both signs of \(w\) proves continuity. Monotonicity
of the secant slopes yields \(q(u)\le v(u)-v(0)\) whenever
\(v(u)\) is finite. This proves the supporting inequality
everywhere.
\end{proof}

The two conclusions of Theorem~\ref{thm:perturbation-duality}
are logically distinct. Lower semicontinuity at a single
point controls the value. A continuous affine minorant
through the graph point controls dual attainment.
Compare \cite[Chapter III, Section 2]{EkelandTemam}
and \cite[Sections 29--30]{Rockafellar}.

\section{Finite-dimensional Fenchel--Rockafellar duality}
\label{v04-02}

\begin{theorem}[Relative-interior qualification]\label{thm:finite-fenchel}
Let \(X,Y\) be finite-dimensional real spaces, let
\(A:X\to Y\) be linear, and let \(f,g\) be proper convex
functions on \(X,Y\), respectively. Suppose
\[
 A(\operatorname{ri}\operatorname{dom}f)
 \cap\operatorname{ri}\operatorname{dom}g\ne\varnothing,\qquad
 p=\inf_x \{f(x)+g(Ax)\}\in\R .
\]
Then
\begin{equation}\label{eq:fenchel-dual}
 p=\max_{y\in Y'}\{-f^*(-A^*y)-g^*(y)\}.
\end{equation}
\end{theorem}

\begin{proof}
Set \(\Phi(x,u)=f(x)+g(Ax+u)\), and let \(v\) be its marginal
value. The set of perturbations admitting finite candidates
is \(S=\operatorname{dom}g-A(\operatorname{dom}f)\).
The relative-interior image and sum formulas of
Section~\ref{v01-02} give \(0\in\operatorname{ri}S\). If
\(v(u)=-\infty\) for some \(u\in S\), a small negative multiple
of \(u\) belongs to \(S\). Combining a finite candidate
there with arbitrarily negative candidates at \(u\) contradicts
\(v(0)=p\in\R\). Thus \(v\) is proper with domain \(S\).
In the linear space \(H=\operatorname{aff}S\), zero is an
interior point of its domain. The finite-dimensional local
upper bound for convex functions (Section~\ref{v02-02})
and Lemma~\ref{lem:local-value-support} supply a subgradient
on \(H\). Extend that functional to \(Y\). Its inequality
still holds, since \(v=+\infty\) outside \(H\).

Writing \(z=Ax+u\), we obtain
\[
 \Phi^*(0,y)=\sup_{x,z}
 \{\langle -A^*y,x\rangle-f(x)+\langle y,z\rangle-g(z)\}
 = f^*(-A^*y)+g^*(y).
\]
Both suprema in this separated expression are bounded below
by finite candidates. Theorem~\ref{thm:perturbation-duality}
therefore yields \eqref{eq:fenchel-dual}.
\end{proof}

The qualification permits a lower-dimensional perturbation
space, even when \(S=\{0\}\). The statement includes proper
closed convex functions. The proof uses no closure assumption
on \(f,g\) beyond the qualification in this theorem.
For the finite-dimensional duality forms see
\cite[Theorems 31.1--31.3]{Rockafellar}.

\section{Banach-space Fenchel duality}
\label{v04-03}

\begin{theorem}[Continuity qualification]\label{thm:banach-fenchel}
Let \(X,Y\) be real Banach spaces and \(A:X\to Y\)
bounded linear. Let \(f,g\) be proper convex functions and
assume there is \(x_0\in\operatorname{dom}f\) such that
\(g\) is finite and continuous at \(Ax_0\), in the norm
topology of \(Y\). If the primal value \(p\) is finite,
then \eqref{eq:fenchel-dual} holds with attainment in the
continuous dual \(Y'\).
\end{theorem}

\begin{proof}
Use the same perturbation. The candidate \(x_0\) gives
\(v(u)\le f(x_0)+g(Ax_0+u)\), which is bounded above on some
ball around zero. By Lemma~\ref{lem:local-value-support}
there is a continuous subgradient at zero. The computation
of \(\Phi^*(0,y)\) uses only the separated suprema and
remains valid. Apply Theorem~\ref{thm:perturbation-duality}.
\end{proof}

Continuity here is in the ambient perturbation space.
Relative-interior qualifications belong to the finite-dimensional
statement. Primal attainment may be obtained, for example,
from coercivity, weak sequential lower semicontinuity and
reflexivity as in Section~\ref{v03-02}. The proof of dual
attainment itself uses the finite primal \emph{value}.
Compare \cite[Chapter III, Propositions 2.1--2.4 and
Theorem 4.1]{EkelandTemam}.

\section{Exact subdifferential rules}
\label{v04-04}

\begin{theorem}[Sum and affine chain rules]\label{thm:exact-subdiff}
Let \(f,g\) be proper convex functions on a real Banach
space. At every point where both are finite,
\[
 \partial(f+g)(x)=\partial f(x)+\partial g(x)
\]
if \(g\) is finite and norm-continuous at a point of
\(\operatorname{dom}f\). In finite dimension the
qualification may instead be
\(\operatorname{ri}\operatorname{dom}f\cap
\operatorname{ri}\operatorname{dom}g\ne\varnothing\).

Let \(B(x)=Ax+b\), where \(A\) is bounded linear between
real Banach spaces and \(g\) is proper convex on the target.
At \(x\) with \(g(Bx)<\infty\),
\[
 \partial(g\circ B)(x)=A^*\partial g(Bx)
\]
if \(g\) is finite and norm-continuous at some \(Bx_0\).
In finite dimension the qualification may instead be
\(B(X)\cap\operatorname{ri}\operatorname{dom}g\ne\varnothing\).
\end{theorem}

\begin{proof}
Adding supporting inequalities gives the easy inclusion
for the sum. For the converse, let \(\ell\in\partial(f+g)(x)\).
The objective \(f-\ell+g\) attains its finite minimum at \(x\).
Apply Theorem~\ref{thm:banach-fenchel} or
Theorem~\ref{thm:finite-fenchel} with \(A=I\) and \(f\)
replaced by \(f-\ell\). The attained dual value is
\(- (f-\ell)^*(-y)-g^*(y)\) for a functional \(y\).
The two Fenchel--Young inequalities at \(x\) have
nonnegative gaps whose sum is zero. Hence
\(-y\in\partial(f-\ell)(x)\) and \(y\in\partial g(x)\).
This decomposes \(\ell\).

Composition of a supporting inequality with \(B\) gives
the easy inclusion for the chain rule. For the converse,
take \(\ell\in\partial(g\circ B)(x)\). Apply the same
duality theorem to the linear function \(-\ell\) and
\(z\mapsto g(z+b)\), with the given \(A\).
The primal minimum is finite and attained at \(x\).
The conjugate of \(-\ell\) equals zero at \(-\ell\)
and \(+\infty\) elsewhere, by scaling a vector at which
two different linear functionals differ.
Finite optimal dual value therefore forces \(A^*y=\ell\).
Fenchel--Young equality for the translated \(g\) gives
\(y\in\partial g(Bx)\).
\end{proof}

\begin{lemma}[Supporting inequality over a polyhedron]\label{lem:polyhedral-support}
Let \(f\) be proper convex on a finite-dimensional space
\(X\), and write
\[
 P=\{x:a_i(x)\le b_i\ (1\le i\le m),\
                 e_j(x)=d_j\ (1\le j\le k)\}
\]
for a polyhedron, with linear \(a_i,e_j\).
If \(\operatorname{ri}\operatorname{dom}f\cap P\ne\varnothing\)
and \(p=\inf_P f\in\R\), then some \(\lambda_i\ge0,\nu_j\in\R\)
satisfy
\[
 f(x)+\sum_i\lambda_i(a_i(x)-b_i)
       +\sum_j\nu_j(e_j(x)-d_j)\ge p
 \qquad(x\in\operatorname{dom}f).
\]
For \(n=\sum_i\lambda_i a_i+\sum_j\nu_j e_j\), the bound
can be written \(f(x)+n(x)-\max_P n\ge p\).
\end{lemma}

\begin{proof}
Work first on \(H=\operatorname{aff}\operatorname{dom}f\),
using affine coordinates. The epigraph of \(f\) has nonempty
interior in \(H\times\R\), since \(f\) is locally bounded
above at an interior domain point. Its interior is disjoint
from \((P\cap H)\times(-\infty,p)\).
The open-set separation theorem gives a nonzero affine
separator. Each point of the epigraph is in the closure
of its interior: combine it with a fixed interior point,
and let the weight of that point tend to zero. Thus the
separating inequality extends to the entire epigraph,
and in a choice of affine coordinates it reads
\[
 a(x)+\alpha t\ge\gamma\ge a(z)+\alpha s
 \quad\bigl((x,t)\in\operatorname{epi}f,\
        z\in P\cap H,\ s<p\bigr).
\]
The upward epigraph ray forces \(\alpha\ge0\).
If \(\alpha=0\), a point \(x_0\in
\operatorname{ri}\operatorname{dom}f\cap P\)
would minimize \(a\) over a neighborhood in \(H\).
Then the linear part of \(a\) vanishes on the direction
space of \(H\), contradicting nonzero separation.
Consequently \(\alpha>0\). With \(q=a/\alpha\), taking
\(s\uparrow p\) gives
\(f(x)+q(x)\ge p+\sup_{P\cap H}q\).
This supremum is finite. The affine image of the
polyhedron \(P\cap H\) under \(q\) is a nonempty closed
subset of \(\R\) by Section~\ref{v01-08}, so it attains
its supremum at some \(z_0\).

Use the polyhedral normal cone formula
(Section~\ref{v01-10}) at \(z_0\) within \(H\).
It represents the linear part of \(q\) on the direction
space of \(H\) as the combination of the restrictions of
\(a_i,e_j\) with \(\lambda_i\ge0\), with positive
coefficients allowed only for active inequalities.
The same combination \(n\) on \(X\) satisfies
\(n(x)-n(z_0)=q(x)-q(z_0)\) for \(x\in H\).
Since \(z_0\) is feasible and the inactive coefficients
are zero, \(n(z)-n(z_0)\le0\) on all of \(P\), and the
expression on the left equals the linear combination
of residuals in the statement, at any \(z\in X\).
This proves the bound and the asserted maximum.
\end{proof}

\begin{theorem}[One polyhedral summand]\label{thm:polyhedral-sum}
Let \(f,g\) be proper convex functions on a
finite-dimensional real space, with \(g\) polyhedral,
meaning that \(\operatorname{epi}g\) is a polyhedron.
If \(\operatorname{ri}\operatorname{dom}f\cap
\operatorname{dom}g\ne\varnothing\), then
\(\partial(f+g)(x)=\partial f(x)+\partial g(x)\)
at every common finite point.
\end{theorem}

\begin{proof}
First take \(g=I_P\), the function zero on a polyhedron
\(P\) and \(+\infty\) elsewhere. For \(\ell\in
\partial(f+I_P)(x)\), apply Lemma~\ref{lem:polyhedral-support}
to \(f-\ell\). Its infimum over \(P\) equals \(f(x)-\ell(x)\).
At \(x\), the resulting bound forces \(n(x)=\max_P n\).
Thus \(n\in N_P(x)\), where \(N_P(x)\) consists of the
linear functionals \(n\) satisfying \(n(z-x)\le0\) on \(P\).
Rearranging the bound gives \(\ell-n\in\partial f(x)\).

For general polyhedral \(g\) and
\(\ell\in\partial(f+g)(x)\), minimize
\(F(z,t)=f(z)-\ell(z)+t\) over
\(P=\operatorname{epi}g\). The point \((x,g(x))\)
is a minimizer, and \(\operatorname{ri}\operatorname{dom}F
=(\operatorname{ri}\operatorname{dom}f)\times\R\)
meets \(P\). The indicator case gives
\(0\in \partial F(x,g(x))+N_P(x,g(x))\).
Testing the two vertical directions in the subgradient
inequality for \(F\), and then the directions with
fixed \(t\), gives
\(\partial F(x,g(x))=(\partial f(x)-\ell)\times\{1\}\).
A normal \((a,-1)\) to \(\operatorname{epi}g\) at
\((x,g(x))\) is exactly \(a\in\partial g(x)\), by testing
\((z,g(z))\) for every finite \(g(z)\).
This decomposes \(\ell\). The reverse inclusion in
both cases follows by addition of the supporting inequalities.
\end{proof}

For convex sets \(C,D\) and \(x\in C\cap D\),
applying the sum rule to indicators gives
\(N_{C\cap D}(x)=N_C(x)+N_D(x)\) under its corresponding
qualification. In particular when \(D\) is a polyhedron
it suffices that \(D\) meet \(\operatorname{ri}C\).
See \cite[Theorems 23.8--23.9]{Rockafellar} and
\cite[Chapter I, Section 5]{EkelandTemam}.

\section{Saddles, certificates, and sensitivity}
\label{v04-05}

\begin{theorem}[Fenchel certificate and saddle]\label{thm:fenchel-saddle}
For proper closed convex \(f,g\) and bounded linear \(A\),
let \(x\) be primal feasible with finite objective and let
\(f^*(-A^*y)\) and \(g^*(y)\) be finite. The primal--dual gap
at \((x,y)\) vanishes precisely when
\[
 -A^*y\in\partial f(x),\qquad y\in\partial g(Ax).
\]
This is equivalent to the saddle condition for
\[
 \mathcal L(z,w)=f(z)+\langle w,Az\rangle-g^*(w)
 \qquad(z\in\operatorname{dom}f,\ w\in\operatorname{dom}g^*),
\]
namely \(\mathcal L(x,w)\le\mathcal L(x,y)
\le\mathcal L(z,y)\) for all \(z,w\) in these domains.
\end{theorem}

\begin{proof}
The gap is the sum of
\[
 f(x)+f^*(-A^*y)+\langle y,Ax\rangle,\qquad
 g(Ax)+g^*(y)-\langle y,Ax\rangle.
\]
Each term is nonnegative and its vanishing is exactly
the indicated subgradient inequality.
Equivalently the first term asserts that \(x\)
minimizes \(z\mapsto f(z)+\langle y,Az\rangle\),
and the second says that \(y\) maximizes
\(w\mapsto\langle w,Ax\rangle-g^*(w)\) at the value
\(g(Ax)\). Fenchel--Moreau (Section~\ref{v02-05})
identifies this supremum with \(g(Ax)\).
These are the two saddle inequalities. A zero gap
also proves that \(x\) and \(y\) attain the primal
and dual optima, by weak duality.
\end{proof}

\begin{proposition}[Value sensitivity]\label{prop:fenchel-sensitivity}
Let \(v\) and \(p=v(0)\) satisfy Theorem~\ref{thm:perturbation-duality}.
Every dual maximizer with value \(p\) gives the global bound
\(v(u)\ge p+\langle y,u\rangle\).
If \(v\) is finite and Fr\'echet differentiable at zero
in a normed space, the unique supporting functional is \(Dv(0)\).
\end{proposition}

\begin{proof}
The bound is the subgradient conclusion of the theorem.
For the differentiable case, on a small ball both
signs of a direction can be used. Comparing
\(v(tw)-v(0)\ge t\langle y,w\rangle\) for \(t>0\)
with its counterpart at \(-tw\) shows
\(\langle y,w\rangle=Dv(0)w\).
Conversely the secant slope along each ray is
nondecreasing, so \(Dv(0)w\le v(w)-v(0)\) at each
finite \(v(w)\). Thus \(Dv(0)\) is a supporting
functional.
\end{proof}

Compare \cite[Chapter III, Sections 2--3]{EkelandTemam}
and \cite[Sections 5.4 and 5.6]{Boyd}.

\section{Slater's condition and affine inequalities}
\label{v04-06}

\begin{theorem}[Refined Slater condition]\label{thm:refined-slater}
On a finite-dimensional real space, minimize a convex
objective \(f\), subject to convex inequalities \(g_i(x)\le0\)
and affine constraints \(a_i(x)\le b_i,\ e_j(x)=d_j\).
Write \(D=\operatorname{dom}f\cap\bigcap_i\operatorname{dom}g_i\)
and let \(P\) be the polyhedron of affine constraints.
Assume the primal value \(p\) is finite and there exists
\(x_0\in\operatorname{ri}D\cap P\) with \(g_i(x_0)<0\)
for every convex inequality in the first group.
Then the Lagrange dual, with its infimum over \(D\),
attains value \(p\). The witness may satisfy some
affine inequalities at equality.
\end{theorem}

\begin{proof}
Keep the affine constraints in \(P\) and perturb only the
first group. Define
\[
 w(u)=\inf\{f(x): x\in D\cap P,\ g_i(x)\le u_i\ \text{for all }i\}.
\]
It is a convex marginal value, is coordinatewise
nonincreasing, and \(w(0)=p\). Strict feasibility of
\(x_0\) gives the finite upper bound \(f(x_0)\) on
a neighborhood of zero. Lemma~\ref{lem:local-value-support}
gives a subgradient \(\eta\) at zero. Since
\(w(t e_i)\le w(0)\) for \(t>0\), its supporting
inequality forces \(\eta_i\le0\). Set
\(\lambda_i=-\eta_i\ge0\).
For any \(x\in D\cap P\), take \(u=(g_i(x))_i\).
The inequalities \(w(u)\le f(x)\) and
\(w(u)\ge p-\sum_i\lambda_i g_i(x)\) give
\(F(x)\ge p\), where \(F=f+\sum_i\lambda_i g_i\)
on \(D\) and \(F=+\infty\) elsewhere.
At a primal minimizing sequence,
\(F(x)\le f(x)\) by feasibility.
Consequently \(\inf_{x\in P}F(x)=p\).
The argument includes the case of an empty first group,
in which \(F\) is just the objective on \(D\).

Apply Lemma~\ref{lem:polyhedral-support} to \(F\) and \(P\),
using the same relative-interior witness.
It gives nonnegative multipliers for the affine
inequalities and free multipliers for the equalities,
with the full Lagrangian bounded below by \(p\) on \(D\).
Conversely, at any feasible point and for any admissible
multipliers, the Lagrangian is at most \(f(x)\).
Infimizing and using weak duality proves that the
constructed multipliers attain \(p\).
\end{proof}

This argument separates the two roles of feasibility.
The strict inequalities furnish an open perturbation
neighborhood. The affine constraints enter through
a polyhedron and its supporting normals.
Compare \cite[Sections 5.2.3 and 5.3.2]{Boyd}
and \cite[Sections 20 and 28]{Rockafellar}.

\section{Convex KKT conditions and boundary examples}
\label{v04-07}

\begin{theorem}[Convex KKT optimality]\label{thm:convex-kkt}
For the convex problem in Theorem~\ref{thm:refined-slater},
a feasible finite point \(x_*\) is a global minimizer
if some admissible multipliers make it minimize
the Lagrangian on \(D\), and the products of each
inequality multiplier and its constraint residual vanish
at \(x_*\).
Conversely, if the hypotheses of the refined Slater
theorem hold and \(x_*\) attains the optimum, there
are such multipliers.
\end{theorem}

\begin{proof}
Under the stated conditions, the Lagrangian at \(x_*\)
equals \(f(x_*)\). At every other feasible point it is at
most the objective there, so its global minimality
proves the desired inequality of objective values.
For necessity, take the dual multipliers from
Theorem~\ref{thm:refined-slater}. They give
\[
 p=\inf_D\mathcal L\le\mathcal L(x_*)\le f(x_*)=p.
\]
Thus \(x_*\) minimizes the Lagrangian. The inequalities
contributing to the difference \(\mathcal L(x_*)-f(x_*)\)
are all nonpositive, so their sum vanishes exactly
when each product vanishes.
\end{proof}

The global minimum of the Lagrangian is equivalently
\(0\in\partial(\mathcal L+I_D)(x_*)\), with the function
extended by \(+\infty\) outside \(D\).
For differentiable data on an open convex domain containing
\(x_*\), the same condition is the usual gradient
stationarity equation, by the convex supporting inequality.
The four conditions are primal feasibility, nonnegative
inequality multipliers, stationarity, and complementarity.

\begin{example}[Equal values without a dual maximizer]\label{ex:nonattained-dual}
Minimize \(x\) on \(\R\) subject to \(x^2\le0\).
The primal value is zero, attained at zero.
For \(\lambda>0\),
\(\inf_x(x+\lambda x^2)=-1/(4\lambda)\), by completing
the square. For \(\lambda=0\) the infimum is \(-\infty\).
Thus the dual value is zero and every finite multiplier
gives a strictly smaller value.
\end{example}

\begin{example}[A positive duality gap]\label{ex:positive-gap}
Minimize \(e^{-x}\) on \(y>0\) subject to \(x^2/y\le0\).
The constraint function is convex as the perspective of
the square, and every feasible point has \(x=0\).
Thus \(p=1\). For each \(\lambda\ge0\), the Lagrangian
\(e^{-x}+\lambda x^2/y\) is nonnegative and tends to zero
along \(x\to+\infty,\ y=x^3\). Its infimum is zero,
so the attained dual value is zero.
\end{example}

These examples give separate failures when the qualification
is absent. See \cite[Section 5.5 and Exercises 5.21--5.22]{Boyd}.
\chapter{Linear, Quadratic, and Conic Programs}

\section{Linear programs, attainment, and duality}
\label{v05-01}

\begin{theorem}[Linear programming duality]\label{thm:linear-programming-duality}
For a real matrix \(A\) and vectors \(b,c\), consider
\[
 \inf\{c^{\mathsf T}x:Ax=b,\ x\ge0\},\qquad
 \sup\{b^{\mathsf T}y:A^{\mathsf T}y\le c\}.
\]
If the primal program is feasible and has finite value,
both optima are attained and their values agree.
For any feasible pair, the gap equals
\(\sum_i x_i(c-A^{\mathsf T}y)_i\); it is zero
exactly when every summand is zero.
\end{theorem}

\begin{proof}
For a feasible pair the gap identity follows from \(Ax=b\).
The terms are nonnegative, proving weak duality.
Suppose the primal has finite infimum \(p\).
The set
\[
 \{(Ax,c^{\mathsf T}x+t):x\ge0,\ t\ge0\}
\]
is a linear image of a polyhedron and is closed by
Section~\ref{v01-08}. A minimizing sequence shows that
\((b,p)\) belongs to this image. A preimage satisfies
\(Ax=b,\ x\ge0\) and \(c^{\mathsf T}x+t=p\).
The definition of \(p\) gives \(c^{\mathsf T}x=p\)
and \(t=0\).

At a primal minimizer \(x_*\), the vector \(-c\)
is normal to the feasible polyhedron. By the exact
normal cone formula of Section~\ref{v01-10}, there are
a free vector \(\xi\) and \(\mu\ge0\) such that
\[
 -c=A^{\mathsf T}\xi-\mu,\qquad
 \mu_i=0\quad\text{if }(x_*)_i>0.
\]
With \(y=-\xi\) this says \(c-A^{\mathsf T}y=\mu\),
so \(y\) is dual feasible. The gap identity and
\(\mu_i(x_*)_i=0\) give the dual maximum.
\end{proof}

\begin{corollary}[Feasibility alternatives]\label{cor:lp-alternatives}
If both programs are feasible, both have finite attained
optima. If precisely one is feasible, that side is unbounded
in its optimizing direction. Both sides can be infeasible.
\end{corollary}

\begin{proof}
With both sides feasible, a single dual feasible value
bounds the primal infimum below, and any primal feasible
value bounds it above. Apply the theorem.
If the primal is feasible and the dual is infeasible,
the theorem rules out a finite primal value.
If the primal is infeasible and \(y_0\) is dual feasible,
Farkas' alternative (Section~\ref{v01-09}) gives \(z\)
with \(A^{\mathsf T}z\ge0\) and \(b^{\mathsf T}z<0\).
Then \(y_0-tz\) is dual feasible for every \(t\ge0\),
and \(b^{\mathsf T}(y_0-tz)\to+\infty\).
For two infeasible sides take the one-dimensional data
\(A=0,\ b=1,\ c=-1\).
\end{proof}

Free variables can be written as differences of
nonnegative variables, and an affine inequality can be
turned into an equality by a nonnegative slack.
These substitutions preserve feasibility and objective
values by construction. They also allow the theorem to
be used for a general finite polyhedral linear program.
See \cite[Theorem 13.1]{NocedalWright} and
\cite[Section 4.3]{Boyd}.

\section{Convex quadratic programs and reduced Hessians}
\label{v05-02}

\begin{theorem}[Convex quadratic programming]\label{thm:convex-qp}
Let \(Q\) be symmetric positive semidefinite, and let \(P\)
be a nonempty polyhedron. If
\[
 p=\inf_{x\in P}\{\tfrac12 x^{\mathsf T}Qx+c^{\mathsf T}x\}
\]
is finite, it is attained. At any feasible \(x_*\), global
optimality is equivalent to the normal cone condition
\[
 -(Qx_*+c)\in N_P(x_*).
\]
For an affine description of \(P\), the latter condition is
equivalent to the existence of KKT multipliers with
complementarity on the inequalities.
\end{theorem}

\begin{proof}
Let \(R=Q^{1/2}\) and use the closed polyhedron
\(S=\{(Rx,c^{\mathsf T}x):x\in P\}\).
The objective on \(S\) is \(\tfrac12\|u\|^2+t\)
in the coordinates \((u,t)\).
The direction \((0,-1)\) cannot be a recession direction
of \(S\), since otherwise the infimum would be \(-\infty\).
Write \(S\) as an intersection of finitely many halfspaces
\(\alpha_i\cdot u+\beta_i t\le \gamma_i\), representing
an equality by two inequalities when necessary.
If all \(\beta_i\ge0\), moving down vertically would
preserve every inequality. Hence some \(\beta_i<0\),
which gives a lower affine bound \(t\ge \ell(u)\).
In any finite sublevel of the objective, the inequality
\(\tfrac12\|u\|^2+\ell(u)\le M\) bounds \(u\),
and then \(\ell(u)\le t\le M-\tfrac12\|u\|^2\)
bounds \(t\). A minimizing sequence lies in a nonempty
compact sublevel, since \(S\) is closed. Its limit point
with value \(p\) lifts to a minimizer in \(P\).

If \(x_*\) minimizes, differentiate the objective along
each feasible segment from \(x_*\); the derivative
\((Qx_*+c)\cdot(x-x_*)\) is nonnegative.
Conversely the exact quadratic expansion at \(x_*\)
has this linear part plus a nonnegative quadratic
part. Thus the normal condition implies optimality.
The formula in Section~\ref{v01-10} decomposes the
normal into a free equality combination and a
nonnegative combination of active inequality
normals. Extending the active coefficients by zero
gives the multiplier and complementarity conditions.
\end{proof}

\begin{proposition}[Equality constraints and the reduced Hessian]\label{prop:reduced-hessian}
Let \(A\) have full row rank and \(Q\) be symmetric.
Assume \(u^{\mathsf T}Qu>0\) for each nonzero
\(u\in\ker A\). For every \(b,c\) the system
\[
 \begin{pmatrix}Q&A^{\mathsf T}\\ A&0\end{pmatrix}
 \begin{pmatrix}x\\ \nu\end{pmatrix}
 =
 \begin{pmatrix}-c\\ b\end{pmatrix}
\]
has a unique solution. Its \(x\) is the unique global
minimum of the quadratic objective on \(Ax=b\).
\end{proposition}

\begin{proof}
In the homogeneous system \(Au=0\) and
\(Qu+A^{\mathsf T}\eta=0\). Pairing the latter with \(u\)
gives \(u^{\mathsf T}Qu=0\), so \(u=0\). Full row rank
then gives \(\eta=0\), proving invertibility.
For any other feasible \(x+u\), the objective difference
is \((Qx+c)^{\mathsf T}u+\tfrac12u^{\mathsf T}Qu\).
The first term vanishes because \(Qx+c=-A^{\mathsf T}\nu\)
and \(Au=0\). The second is strictly positive for
nonzero \(u\), proving the claim.
\end{proof}

Positivity of the reduced Hessian is positivity of the
quadratic form only on feasible differences.
Global semidefiniteness alone still permits an infeasible
constraint system or an unbounded linear objective on
the kernel of \(Q\). Compare
\cite[Sections 16.1 and 16.3]{NocedalWright} and
\cite[Section 4.4]{Boyd}.

\section{Conic duality and relative strict feasibility}
\label{v05-03}

Let \(X,Y\) be finite-dimensional real inner-product
spaces and \(K\subset X\) a closed convex cone.
Its positive dual is
\(K^*=\{s:\langle s,x\rangle\ge0\text{ for all }x\in K\}\).
The programs associated to \(A:X\to Y\), \(b\in Y\),
and \(c\in X\) are
\begin{equation}\label{eq:conic-primal-dual}
 \inf_{x\in K,\ Ax=b}\langle c,x\rangle,\qquad
 \sup_{c-A^*y\in K^*}\langle b,y\rangle .
\end{equation}
For a feasible pair the gap is
\(\langle c-A^*y,x\rangle\ge0\).
Thus zero gap at a feasible pair is equivalent to
the complementary pairing \(\langle c-A^*y,x\rangle=0\).

\begin{theorem}[Conic Slater condition]\label{thm:conic-slater}
If the primal value in \eqref{eq:conic-primal-dual}
is finite and \(Ax_0=b\) for some
\(x_0\in\operatorname{ri}K\), then the dual maximum
is attained and equals the primal value.
\end{theorem}

\begin{proof}
Let \(V=\operatorname{span}K\) and \(W=A(V)\).
On \(W\) consider the convex value
\(v(u)=\inf\{\langle c,x\rangle:x\in K,\ Ax=b+u\}\).
Choose a linear right inverse \(R:W\to V\) to \(A|_V\).
Since \(x_0\) is relatively interior in \(V\),
\(x_0+Ru\in K\) for all sufficiently small \(u\).
These feasible perturbations bound \(v\) above
near zero by \(\langle c,x_0+Ru\rangle\).
Lemma~\ref{lem:local-value-support} gives a supporting
linear functional at zero in \(W\), represented by
\(y\in Y\) after extension to \(Y\).
Writing \(p=v(0)\), the support gives
\[
 \langle c-A^*y,x\rangle\ge p-\langle y,b\rangle
 \qquad(x\in K).
\]
Taking \(x=0\) gives \(\langle y,b\rangle\ge p\).
Scaling any \(x\in K\) by arbitrarily large
positive numbers forces \(c-A^*y\in K^*\).
Weak duality gives the opposite inequality
\(\langle y,b\rangle\le p\), proving the result.
\end{proof}

\begin{corollary}[Strict feasibility on the dual side]\label{cor:dual-slater}
If the dual value is finite and \(c-A^*y_0\in
\operatorname{ri}K^*\) for some \(y_0\), the primal
minimum is attained and equals that dual value.
\end{corollary}

\begin{proof}
Rewrite the dual as a minimization of
\(-\langle b,y\rangle\) over
\((y,s)\in Y\times K^*\) with \(A^*y+s=c\).
The pair \((y_0,c-A^*y_0)\) belongs to the
relative interior of the product cone.
Apply Theorem~\ref{thm:conic-slater} to that problem.
The product cone has positive dual \(\{0\}\times K\)
by the bipolar theorem of Section~\ref{v01-05}.
With multiplier \(z\in X\), feasibility of the
new dual says
\((-b,0)-(Az,z)\in\{0\}\times K\), or \(Az=-b\)
and \(-z\in K\). Its maximized objective is
\(\langle c,z\rangle\). Changing variables to
\(x=-z\) gives exactly the claimed primal attainment.
\end{proof}

A relative interior witness on a given side provides
attainment on the opposite side under the finite-value
hypothesis. See \cite[Section 5.9]{Boyd}.

\section{Standard cones and quadratic epigraphs}
\label{v05-04}

\begin{proposition}[Self-dual cones]\label{prop:standard-self-dual-cones}
The nonnegative orthant is self-dual under the Euclidean
pairing. The Lorentz cone
\[
 \mathcal Q=\{(t,u):t\ge \|u\|_2\}
\]
is self-dual under the product Euclidean pairing.
The cone \(\mathbb S^n_+\) of positive semidefinite
symmetric matrices is self-dual under
\(\langle S,T\rangle=\operatorname{tr}(ST)\).
\end{proposition}

\begin{proof}
For the orthant, nonnegative vectors pair nonnegatively
with each other, and testing on coordinate vectors
recovers nonnegativity. For the Lorentz cone,
if \(a\ge\|b\|\) and \(t\ge\|u\|\) then
\(at+b\cdot u\ge at-\|b\|\|u\|\ge0\).
If \(a<\|b\|\) and \(b\ne0\), the test
\((1,-b/\|b\|)\in\mathcal Q\) pairs negatively.
If \(b=0\) and \(a<0\), use \((1,0)\).
For positive semidefinite matrices \(S,T\),
diagonalize \(T=\sum_j t_j v_j v_j^{\mathsf T}\)
with \(t_j\ge0\). The trace pairing equals
\(\sum_j t_j v_j^{\mathsf T}Sv_j\ge0\).
Conversely, testing \(vv^{\mathsf T}\) recovers
\(v^{\mathsf T}Sv\ge0\) for every \(v\).
\end{proof}

\begin{proposition}[A Schur complement with a singular block]\label{prop:singular-schur}
Let \(R\) be symmetric positive semidefinite. Then
\[
 \begin{pmatrix}t&w^{\mathsf T}\\w&R\end{pmatrix}\succeq0
 \quad\Longleftrightarrow\quad
 w\in\operatorname{ran}R,\qquad t\ge w^{\mathsf T}R^\dagger w ,
\]
where \(R^\dagger\) acts as the inverse on
\(\operatorname{ran}R\) and as zero on \(\ker R\).
\end{proposition}

\begin{proof}
If the block matrix is positive semidefinite, test
\((1,sv)\) for \(v\in\ker R\) and all \(s\in\R\).
The resulting affine expression in \(s\) is
nonnegative, so \(w\cdot v=0\).
Thus \(w\in(\ker R)^\perp=\operatorname{ran}R\).
Put \(z=R^\dagger w\). For arbitrary \(r,v\) the
quadratic form is
\[
 r^2(t-z^{\mathsf T}Rz)+(v+rz)^{\mathsf T}R(v+rz).
\]
Taking \(v=-rz\) gives necessity of \(t\ge z^{\mathsf T}Rz\).
This condition also makes the expression nonnegative
for all \(r,v\), proving sufficiency.
Finally \(z^{\mathsf T}Rz=w^{\mathsf T}R^\dagger w\).
\end{proof}

A constraint \(\|Bx+d\|_2\le a^{\mathsf T}x+\beta\)
is therefore an affine Lorentz-cone constraint.
For the squared norm epigraph
\(t\ge\|Bx+d\|_2^2\), take the block matrix above
with \(w=Bx+d\) and \(R=I\).
An affine cone constraint can be written in standard
form by introducing a cone variable and an equality
that identifies it with the affine expression. Any free
variables can be included in a product cone with
a full linear space. These substitutions give the
usual second-order and semidefinite representations.
See \cite[Sections 2.6 and 4.4--4.6 and Appendix A]{Boyd}.

\section{Closure, weak infeasibility, and nonattainment}
\label{v05-05}

\begin{proposition}[Nonclosed image and weak infeasibility]\label{prop:sdp-weak-infeasibility}
The projection of \(\begin{pmatrix}a&b\\b&c\end{pmatrix}\in\mathbb S^2_+\)
to \((a,b)\) has image \(\{(a,b):a>0\}\cup\{(0,0)\}\).
The target \((0,1)\) is weakly infeasible. It lies in the closure of the
image, and there is no strict separating conic Farkas certificate.
\end{proposition}

\begin{proof}
For \(a>0\), every \(b\) is possible
by taking \(c=b^2/a\). For \(a=0\), positivity
forces \(b=0\), by testing \((r,1)\) and varying \(r\).
The image is consequently
\(\{(a,b):a>0\}\cup\{(0,0)\}\).
The target \((0,1)\) is infeasible but is approached
by the images of
\(\begin{pmatrix}1/n&1\\1&n\end{pmatrix}\).
The matrices are positive semidefinite by the
Schur complement. A continuous linear functional
nonnegative on the image remains nonnegative on
its closure. Hence no such functional can be
strictly negative on this infeasible target,
as a strict Farkas certificate would require.
\end{proof}

\begin{proposition}[Strict feasibility with an unattained primal infimum]\label{prop:sdp-nonattainment}
The problem of minimizing \(a\) subject to
\(X=\begin{pmatrix}a&1\\1&c\end{pmatrix}\succeq0\) is strictly feasible.
Its finite primal infimum is unattained and equals its attained dual maximum.
\end{proposition}

\begin{proof}
The point with \(a=c=2\) is positive definite.
Every feasible \(a\) is positive, and the matrices
in the preceding example show that the infimum
is zero. It is unattained. The conic dual with
equality \(X_{12}=1\) has feasible multiplier
\(y=0\), since the cost matrix is
\(\operatorname{diag}(1,0)\). Its value is zero,
so the dual maximum is attained.
\end{proof}

\section{Finite matrix minimax}
\label{v05-06}

\begin{theorem}[Finite matrix game]\label{thm:matrix-minimax}
For a real \(m\times n\) payoff matrix \(M\) with
\(m,n\ge1\), let \(\Delta_m,\Delta_n\) be the
probability simplexes. Then
\[
 \max_{p\in\Delta_m}\min_{q\in\Delta_n}p^{\mathsf T}Mq
 =
 \min_{q\in\Delta_n}\max_{p\in\Delta_m}p^{\mathsf T}Mq .
\]
There are maximizing and minimizing mixed strategies
forming a saddle at the common value.
\end{theorem}

\begin{proof}
For fixed \(p\), the minimum over \(q\) is
\(\min_j(M^{\mathsf T}p)_j\), since every mixed
payoff is a convex combination of the pure payoffs.
Thus the negative of the left-hand value is the
LP minimum of \(-\alpha\) subject to
\[
 p\ge0,\quad \mathbf1^{\mathsf T}p=1,\quad
 \alpha\mathbf1-M^{\mathsf T}p\le0.
\]
Its value is finite and it is feasible, because all
matrix entries are finite and at least one mixed
strategy exists. Use a nonnegative multiplier \(q\)
for the inequality and a free multiplier \(\beta\)
for the equality. Infimizing the Lagrangian
\[
 -\alpha+q^{\mathsf T}(\alpha\mathbf1-M^{\mathsf T}p)
       +\beta(\mathbf1^{\mathsf T}p-1)
\]
over \(\alpha\in\R,\ p\ge0\) gives \(-\beta\) if
\(q\in\Delta_n\) and \(Mq\le\beta\mathbf1\),
and \(-\infty\) otherwise.
This is exactly the standard linear-program dual
after adding slack variables for the inequalities
and splitting \(\alpha\) into its positive and
negative parts. Theorem~\ref{thm:linear-programming-duality}
therefore identifies the primal minimum with
\(-\min_{q\in\Delta_n}\max_i(Mq)_i\), with both
optima attained. If \(p_*,q_*\) attain the values,
then \(p_*^{\mathsf T}Mq\) is at least the value
for all \(q\), and \(p^{\mathsf T}Mq_*\) is at most
the value for all \(p\). Taking \(p=p_*,q=q_*\)
gives the claimed saddle.
\end{proof}

This is the finite polyhedral minimax result. The
general topological minimax problem has different
compactness and continuity assumptions.
See \cite[Sections 36--37]{Rockafellar} and
\cite[Chapter VI]{EkelandTemam}.
\chapter{Smooth Optimality and Constraint Qualifications}

\section{Unconstrained first- and second-order tests}
\label{v06-01}

\begin{theorem}[Unconstrained optimality tests]\label{thm:unconstrained-tests}
Let \(f\) be real-valued on an open subset of a
finite-dimensional real inner-product space,
and let \(x_*\) belong to that open set.
If \(f\) is differentiable at a local minimizer \(x_*\),
then \(Df(x_*)=0\).
If it is twice Fr\'echet differentiable there, its
Hessian at the local minimizer is positive semidefinite.
Conversely, if \(f\) is twice Fr\'echet differentiable at \(x_*\),
\(Df(x_*)=0\), and the Hessian is positive definite,
there are \(c,\rho>0\) such that
\[
 f(x)\ge f(x_*)+c\|x-x_*\|^2
 \qquad(\|x-x_*\|<\rho).
\]
\end{theorem}

\begin{proof}
Restrict \(f\) to the line \(x_*+td\), using
both signs of \(t\). At a local minimum the first
derivative in every direction is zero.
For the second-order statements use
\[
 f(x_*+h)-f(x_*)=Df(x_*)h+
           \tfrac12D^2f(x_*)[h,h]+o(\|h\|^2).
\]
For clarity, twice Fr\'echet differentiability at
\(x_*\) gives
\(Df(x_*+z)=Df(x_*)+D^2f(x_*)[z,\cdot]+r(z)\)
with \(\|r(z)\|/\|z\|\to0\).
Along \(z=th\), \(0\le t\le1\), subtract
the constant, linear, and quadratic terms from
\(f(x_*+th)\). The derivative of the remainder is
\(r(th)h\), bounded in absolute value by
\(\epsilon\|h\|^2\) for sufficiently small \(h\),
uniformly in \(t\). The one-dimensional mean value
theorem gives the displayed expansion. Substituting
\(h=td\) at the minimum and dividing by \(t^2\)
gives \(D^2f(x_*)[d,d]\ge0\).
For the converse, the positive quadratic form has
a positive minimum \(m\) on the unit sphere,
by compactness. Its half in the expansion is at
least \(\tfrac12m\|h\|^2\); choose a neighborhood
where the remainder has modulus at most
\(\tfrac14m\|h\|^2\). In dimension zero the growth
inequality holds for every \(c>0\).
\end{proof}

At zero, \(x^4\) and \(-x^4\) have the same vanishing
gradient and Hessian. They give a strict minimum and
a strict maximum respectively. In the real Hilbert
space \(\ell^2\), the smooth function
\[
 F(x)=\tfrac12\sum_{j=1}^\infty 2^{-j}x_j^2-\|x\|^4
\]
has zero derivative at zero and a Hessian which is
strictly positive on every nonzero vector. Its
quadratic part is a bounded quadratic form, and
the fourth power of the norm is the square of the
inner product \(\langle x,x\rangle\), so both are
smooth. At \(x=2^{-j/2}e_j\) the value is
\(-2^{-2j}/2\). These vectors tend to zero.
Thus pointwise Hessian positivity in infinite
dimension requires a separate check for a uniform
positive lower bound if it is to yield this growth
argument. The convex first-order alternative is
given in Section~\ref{v03-03}.
Compare \cite[Theorems 2.1--2.5]{NocedalWright}.

\section{Equality manifolds and multipliers}
\label{v06-02}

\begin{lemma}[Local equality chart]\label{lem:equality-chart}
Let \(h:U\to Y\) be \(C^r\), \(r\ge1\), on an open subset
of a real Banach space \(X\), with \(Y\) a real Banach
space. Suppose \(h(x_*)=0\) and \(B=Dh(x_*)\) has a
bounded linear right inverse \(R:Y\to X\).
Near \(x_*\) the feasible set is a \(C^r\) graph
over \(K=\ker B\). In particular every \(d\in K\)
is the derivative at zero of a \(C^r\) feasible
curve through \(x_*\), defined for both signs of its
parameter.
\end{lemma}

\begin{proof}
Since \(BR=I\), the maps \(RB\) and \(I-RB\) are
bounded complementary projections, with ranges
\(M=R(Y)\) and \(K\). Hence \(X=K\oplus M\).
The restriction \(B|_M:M\to Y\) has bounded inverse
\(R\). Apply the implicit function theorem to
\((k,m)\mapsto h(x_*+k+m)\), solving the equation
for \(m=\psi(k)\) near \(k=0\).
The derivative in \(k\) at zero vanishes, so
differentiating the graph equation gives \(D\psi(0)=0\).
The curve \(x_*+td+\psi(td)\) is feasible and has
the stated derivative. Every sufficiently close
feasible point belongs to the graph by the
local uniqueness in the implicit function theorem.
\end{proof}

\begin{theorem}[Equality multiplier rule]\label{thm:equality-multipliers}
Under the assumptions of Lemma~\ref{lem:equality-chart},
if \(x_*\) locally minimizes \(f\) on \(h=0\) and \(f\)
is differentiable at \(x_*\), there is a unique
\(\nu\in Y'\) such that
\[
 Df(x_*)+Dh(x_*)^*\nu=0.
\]
In finite dimension surjectivity of \(Dh(x_*)\)
is sufficient for the graph and multiplier conclusions.
\end{theorem}

\begin{proof}
For any \(d\in K\), the two-sided feasible curve
forces \(Df(x_*)d=0\). Since \(x-RBx\in K\) for
every \(x\in X\), we have \(Df(x_*)x=Df(x_*)RBx\).
Thus \(\nu=-Df(x_*)R\) gives the equation.
Surjectivity of \(B\) makes \(B^*\) injective,
so the multiplier is unique. In finite dimension a
surjective linear map has a bounded linear right
inverse, by choosing a complement of its kernel.
\end{proof}

The split-surjectivity formulation also applies to
finitely many equality constraints on a Banach
space, since a basis of their finite-dimensional
target has finitely many preimages defining a
bounded right inverse. This is the chart used for
regular integral constraints in
Section~\ref{v07-04}. See
\cite[Section 1.2.2]{Liberzon} and
\cite[Section 12.3]{NocedalWright}.

\section{Active constraints, tangency, and qualification}
\label{v06-03}

Consider \(C^1\) data on an open subset of a
finite-dimensional real space, with feasible set
\[
 S=\{x:g_i(x)\le0\ (1\le i\le m),\
                  h_j(x)=0\ (1\le j\le k)\}.
\]
At a feasible point \(x_*\), let
\(I=\{i:g_i(x_*)=0\}\).
The Bouligand tangent cone \(T_S(x_*)\) consists of
vectors \(d\) obtained as limits of
\((x_n-x_*)/t_n\), with \(x_n\in S\) and
\(t_n>0\) tending to zero. The linearized cone is
\[
 C_{\rm lin}
 =\{d:Dh(x_*)d=0,\ Dg_i(x_*)d\le0\text{ for }i\in I\}.
\]
Here \(h\) denotes the vector of equalities. The
linear independence constraint qualification,
abbreviated LICQ, requires that the equality and
active inequality differentials be linearly independent.
The Mangasarian--Fromovitz constraint qualification,
abbreviated MFCQ, requires surjectivity of \(Dh(x_*)\)
and a vector \(w\in\ker Dh(x_*)\) with
\(Dg_i(x_*)w<0\) for every \(i\in I\).

\begin{proposition}[Tangent cone under MFCQ]\label{prop:tangent-mfcq}
Always \(T_S(x_*)\subseteq C_{\rm lin}\).
Under MFCQ these cones are equal.
LICQ implies MFCQ.
\end{proposition}

\begin{proof}
Expand the constraints at \(x_*\) along the
feasible sequence in the definition of \(T_S(x_*)\).
The remainder is \(o(\|x_n-x_*\|)\); divided by
\(t_n\) it tends to zero because the scaled
increments converge. This proves inclusion.
The tangent cone is closed: if \(d_\ell\to d\)
and each \(d_\ell\) is tangent, choose one
feasible scaled increment within \(1/\ell\) of
\(d_\ell\), with scale less than \(1/\ell\).
These increments witness tangency of \(d\).

Under MFCQ, for \(d\in C_{\rm lin}\) and
\(\epsilon>0\), \(d+\epsilon w\in\ker Dh(x_*)\).
It is the derivative of an equality-feasible
curve by Lemma~\ref{lem:equality-chart}.
Every active inequality has a strictly negative
derivative in this direction, so it holds for
small positive time. Every inactive inequality
has a strict slack at \(x_*\), so it also holds
for small time by continuity.
Hence \(d+\epsilon w\in T_S(x_*)\).
Use closedness and let \(\epsilon\downarrow0\).

Under LICQ, the linear map with coordinates
\(Dh(x_*)\) and \(Dg_i(x_*)\) for \(i\in I\)
is surjective. Prescribe zero equality derivatives
and \(-1\) for each active inequality derivative.
This produces the required \(w\). For empty
\(I\), the strict inequalities are vacuous.
\end{proof}

Compare \cite[Definitions 12.1 and 12.5 and
Sections 12.3 and 12.6]{NocedalWright} and
\cite[Section 3]{MangasarianFromovitz}.

\section{Fritz John and the first-order KKT rule}
\label{v06-04}

For this sign convention the KKT system is
primal feasibility, \(\lambda_i\ge0\),
\(\lambda_i g_i(x_*)=0\), and
\[
 Df(x_*)+\sum_i\lambda_iDg_i(x_*)+
                    \sum_j\nu_jDh_j(x_*)=0 .
\]
The equality multipliers \(\nu_j\) are free.

\begin{theorem}[Fritz John and KKT necessity]\label{thm:fritz-john}
At a \(C^1\) constrained local minimum \(x_*\)
there are \(\lambda_0,\lambda_i\ge0\) and free
\(\nu_j\), jointly nonzero, with
\[
 \lambda_0 Df(x_*)+\sum_i\lambda_iDg_i(x_*)+
                       \sum_j\nu_jDh_j(x_*)=0,
 \qquad \lambda_i g_i(x_*)=0.
\]
Under MFCQ, \(\lambda_0>0\), so normalization gives
the KKT system. In particular LICQ suffices.
\end{theorem}

\begin{proof}
If \(Dh(x_*)\) is not onto, choose a nonzero
\(\nu\) in the kernel of \(Dh(x_*)^*\) and
set \(\lambda_0=\lambda_i=0\).
Suppose instead that \(Dh(x_*)\) is onto and
put \(K=\ker Dh(x_*)\).
There is no \(d\in K\) for which
\(Df(x_*)d<0\) and \(Dg_i(x_*)d<0\) for all
\(i\in I\). Indeed the equality chart would
realize that \(d\) as a curve that is feasible
for small positive time and decreases \(f\).
The strict-direction Gordan alternative of
Section~\ref{v01-09}, applied to the linear map
\(d\mapsto(Df(x_*)d,(Dg_i(x_*)d)_{i\in I})\)
on \(K\), therefore gives a nonzero nonnegative
combination of these coordinate functionals
that vanishes on \(K\). A linear functional
vanishing on \(K\) factors through \(Dh(x_*)\):
evaluate it on a right inverse and use
\(d-RDh(x_*)d\in K\). Subtract the resulting
equality combination to obtain the stationarity
equation, and assign zero coefficients to inactive
inequalities.

Now assume MFCQ. If \(\lambda_0=0\), evaluate
stationarity on the strict inward vector \(w\).
The equality terms vanish and the active
inequality terms are nonpositive, strictly
negative whenever their coefficients are
positive. Hence all those coefficients vanish.
Surjectivity then forces \(\nu=0\), contrary to
the jointly nonzero condition. Thus
\(\lambda_0>0\) and division gives KKT.
\end{proof}

An equivalent qualified route uses tangency.
Every tangent direction at a differentiable
local minimum satisfies \(Df(x_*)d\ge0\), by
dividing the objective difference along its
witness sequence by the positive scales.
By Proposition~\ref{prop:tangent-mfcq},
\(-Df(x_*)\) lies in the polar of \(C_{\rm lin}\).
The polyhedral polar formula of Section~\ref{v01-10}
then gives a nonnegative combination of the
active inequality differentials plus a free
equality combination. This is the same KKT
system. For the abnormal alternative compare
\cite[Sections 2--3]{MangasarianFromovitz};
for the qualified tangent formulation see
\cite[Theorems 12.1 and 12.9]{NocedalWright}.

\section{Multiplier uniqueness, compactness, and degeneracy}
\label{v06-05}

\begin{proposition}[Qualified multiplier sets]\label{prop:qualified-multiplier-sets}
At a constrained local minimizer satisfying LICQ,
the KKT multiplier is unique.
At a constrained local minimizer satisfying MFCQ,
the multiplier set is nonempty and compact.
\end{proposition}

\begin{proof}
Existence in both cases follows from
Theorem~\ref{thm:fritz-john}. Under LICQ,
subtracting two stationarity equations gives
a linear dependence of the active and equality
differentials. Inactive multipliers are zero,
so all coefficients in the difference vanish.

For MFCQ, the multiplier set at \(x_*\) is closed
by its linear equations and inequalities.
If it were unbounded, take a sequence of multiplier
pairs with norms tending to infinity and divide
each pair by its norm. Pass to a subsequence
converging to a unit pair. The limit has
nonnegative inequality coefficients, inactive
coefficients zero, and homogeneous stationarity,
since the divided \(Df(x_*)\) tends to zero.
Test the limit on the strict inward vector as
in the abnormal case of Theorem~\ref{thm:fritz-john}.
It gives all inequality coefficients zero,
and equality independence gives the remaining
coefficients zero. This contradicts the unit norm.
The set is thus bounded and closed in finite
dimension.
\end{proof}

For \(\min(-x)\) subject to two copies of
\(x\le0\), MFCQ holds at zero with \(w=-1\).
The multipliers are all \((\lambda_1,\lambda_2)
\ge0\) with \(\lambda_1+\lambda_2=1\).
This shows how compactness can coexist with
nonuniqueness. Complementarity only requires
a zero product \(\lambda_i g_i(x_*)=0\).
Strict complementarity additionally requires
\(\lambda_i>0\) for every active inequality.
For \(\min x^2\) subject to \(x\le0\), LICQ
holds at zero, and its unique multiplier is zero.
Finally for \(\min x\) subject to \(x^2\le0\),
zero is the only feasible point and stationarity
cannot hold with an objective coefficient one,
because \(Dg(0)=0\) and \(Df(0)=1\).
It has an abnormal Fritz John multiplier,
with objective coefficient zero and inequality
coefficient one.

\section{The critical cone and second-order necessity}
\label{v06-06}

At a feasible KKT point put
\[
 C_{\rm crit}=\{d\in C_{\rm lin}: Df(x_*)d\le0\},
 \qquad
 L=f+\sum_i\lambda_i g_i+\sum_j\nu_j h_j.
\]
For \(d\in C_{\rm lin}\), stationarity gives
\[
 Df(x_*)d=-\sum_{i\in I}\lambda_i Dg_i(x_*)d\ge0 .
\]
Thus a critical direction has \(Df(x_*)d=0\)
and \(Dg_i(x_*)d=0\) whenever \(\lambda_i>0\).
Conversely those equalities for the positive
multipliers characterize criticality among
linearized directions.

\begin{theorem}[Second-order necessary condition under LICQ]\label{thm:licq-second-order-necessity}
Assume \(f,g_i,h_j\) are \(C^2\) near a constrained
local minimum \(x_*\) satisfying LICQ, and use its
KKT multipliers in \(L\). Then
\[
 D^2L(x_*)[d,d]\ge0\qquad(d\in C_{\rm crit}).
\]
\end{theorem}

\begin{proof}
Fix \(d\in C_{\rm crit}\). Let
\(J=\{i\in I:Dg_i(x_*)d=0\}\).
The differentials of \(h\) and \(g_i\) for \(i\in J\)
are an independent subfamily of the LICQ
differentials. The equality chart applied to this
family realizes \(d\) by a \(C^2\) curve \(\gamma\)
through \(x_*\) on which \(h=0\) and \(g_i=0\) for
\(i\in J\). The other active inequalities have
strictly negative derivative, and the inactive
ones keep their slack. Thus \(\gamma(t)\) is
feasible for sufficiently small \(t\ge0\).
Since every positive inequality multiplier
corresponds to an index in \(J\), the identity
\(L(\gamma(t))=f(\gamma(t))\) holds along this
curve. Also \(DL(x_*)=0\).
The one-sided local minimum, with first derivative
zero, gives \((L\circ\gamma)''(0)\ge0\).
The second derivative chain rule reduces it to
\(D^2L(x_*)[d,d]\), because the term
\(DL(x_*)\gamma''(0)\) vanishes.
\end{proof}

The use of LICQ is in realizing each critical
direction with the required active equalities.
See \cite[Theorem 12.5]{NocedalWright}.

\section{Fixed-multiplier second-order sufficiency}
\label{v06-07}

\begin{theorem}[Constrained quadratic growth]\label{thm:constrained-growth}
Let \(f,g_i,h_j\) be \(C^2\) near a feasible point
\(x_*\), and suppose a fixed multiplier pair
satisfies its KKT system. If
\[
 D^2L(x_*)[d,d]>0
 \qquad(d\in C_{\rm crit}\setminus\{0\}),
\]
then some \(c,\rho>0\) satisfy
\(f(x)\ge f(x_*)+c\|x-x_*\|^2\) at every feasible
\(x\) with \(\|x-x_*\|<\rho\).
In particular \(x_*\) is a strict local minimizer
relative to the feasible set.
\end{theorem}

\begin{proof}
If no such \(c,\rho\) exist, choose feasible
\(x_n\ne x_*\) with \(r_n=\|x_n-x_*\|<1/n\) and
\[
 f(x_n)-f(x_*)<r_n^2/n .
\]
Pass to a subsequence such that the unit directions
\(d_n=(x_n-x_*)/r_n\) tend to a unit vector \(d\).
Feasibility gives \(d\in C_{\rm lin}\), by the
first-order constraint expansions. Dividing the
displayed bound by \(r_n\) and differentiating
the objective gives \(Df(x_*)d\le0\). Hence
\(d\in C_{\rm crit}\).
At \(x_*\), complementarity gives \(L(x_*)=f(x_*)\);
at each feasible \(x_n\), the multiplier signs
give \(L(x_n)\le f(x_n)\).
Stationarity and Taylor's formula imply
\[
 \frac{f(x_n)-f(x_*)}{r_n^2}
 \ge
 \frac{L(x_n)-L(x_*)}{r_n^2}
 =\tfrac12D^2L(x_*)[d_n,d_n]+o(1).
\]
The right-hand side tends to a positive number,
contradicting the original bound. If
\(C_{\rm crit}=\{0\}\), the unit critical limit
is already impossible.
\end{proof}

The sufficient condition starts from a KKT pair
and uses only feasible secants and their finite-dimensional
compactness. It therefore applies without a constraint
qualification. Compare \cite[Theorem 12.6]{NocedalWright}.
\chapter{First and Second Variations}

\section{Fixed endpoints and the Euler--Lagrange equation}
\label{v07-01}

Let \(a<b\), let \(u:[a,b]\to\mathbb R^m\), and write
\[
 J(u)=\int_a^b L(t,u(t),u'(t))\,dt,\qquad
 p(t)=L_v(t,u(t),u'(t)),\qquad q(t)=L_z(t,u(t),u'(t)).
\]
Here \(z\) denotes the state and \(v\) the velocity. We use the
Euclidean pairing for their derivatives. For \(C^1\) curves with
fixed endpoint values, a \emph{weak local minimum} means a minimum
in a \(C^1\) neighborhood within the admissible class. A
\emph{strong local minimum} uses only the uniform distance between
curves. This convention for weak minima concerns the topology of
the admissible class in this chapter. It is separate from weak
convergence in a Banach space. Later sections specify when
piecewise \(C^1\) competitors are allowed.

\begin{lemma}[Fundamental lemma and distributional derivative]
\label{lem:variation-fundamental}
If \(h\in L^1_{\mathrm{loc}}(a,b)\) and
\(\int h\phi=0\) for all \(\phi\in C_c^\infty(a,b)\), then \(h=0\)
almost everywhere. In particular a continuous such \(h\) vanishes
everywhere. If \(p\in L^1_{\mathrm{loc}}(a,b;\mathbb R^m)\),
\(q\in L^1_{\mathrm{loc}}(a,b;\mathbb R^m)\), and
\(\int(p\cdot\phi'+q\cdot\phi)=0\) for all vector tests, then
\(p\) has a locally absolutely continuous representative with
derivative \(q\) almost everywhere.
\end{lemma}
\begin{proof}
For the first assertion, let \([c,d]\subset(a,b)\) with \(c<d\).
Smooth compactly supported functions in \((a,b)\), bounded by
\(1\), can approximate the indicator of \((c,d)\) pointwise
except at its endpoints. For instance, convolve the indicator
with a smooth unit-mass kernel supported in an interval whose
radius tends to zero. Dominated convergence gives
\(\int_c^d h=0\). A local primitive of \(h\) is absolutely
continuous and consequently constant on each compact
subinterval. Thus its derivative \(h\) is zero almost everywhere.
For the second assertion take a local primitive \(Q\) of \(q\)
and apply the hypothesis to obtain
\(\int(p-Q)\cdot\phi'=0\). For a scalar test \(\psi\) of zero
integral, its primitive starting to the left of its support is
again compactly supported, so each component of \(p-Q\)
annihilates \(\psi\). Choose a test \(\rho\) of integral one.
Subtracting \((\int\psi)\rho\) from an arbitrary test \(\psi\)
shows that each component of \(p-Q\) equals a constant by the
first assertion. The locally absolutely continuous
representatives agree on overlaps and have the stated derivative.
\end{proof}

\begin{theorem}[First variation at fixed endpoints]
\label{thm:first-variation}
Suppose \(u\in C^1([a,b];\mathbb R^m)\) is a weak or strong local
minimum under fixed endpoints and \(L\) is \(C^1\) on an open
neighborhood of the path \((t,u(t),u'(t))\). For every
\(\eta\in C^1([a,b];\mathbb R^m)\) with \(\eta(a)=\eta(b)=0\),
\[
 DJ(u)[\eta]=\int_a^b (q\cdot\eta+p\cdot\eta')\,dt=0.
\]
Thus \(p'=q\) in distributions. Under the given continuity,
\(p\) is \(C^1\) and this is a pointwise equation on \((a,b)\).
The same interior weak identity holds if the unperturbed energy is finite
and, for each compactly supported smooth \(\eta\), zero is a two-sided
local minimum of \(s\mapsto J(u+s\eta)\), the integrand is differentiable
in \(s\) near zero for almost every \(t\), and its derivative there is
bounded by one integrable function of \(t\).
\end{theorem}
\begin{proof}
For small \(|s|\), \(u+s\eta\) belongs to the chosen neighborhood
and fixes the endpoints. For the classical hypotheses the
perturbed graphs remain in a compact subset of the differentiability
domain, so the derivative of the integrand in \(s\) is bounded
by a constant times \(|\eta|+|\eta'|\).
The mean value formula and dominated convergence justify
differentiation under the integral. The integrable-bound version
uses the same argument in each tested direction. A differentiable
scalar function at a two-sided local minimum has derivative zero.
Tests supported in \((a,b)\) then give the weak identity.
Lemma~\ref{lem:variation-fundamental} identifies the derivative.
Continuity of \(p,q\) identifies \(p\) with the \(C^1\)
primitive plus its constant everywhere.
\end{proof}

See \cite[Section 1.3.3 and Theorem 2.1]{Dacorogna} and
\cite[Sections 2.2--2.3]{Liberzon}. The compactly supported test
form is also the starting point of Section~\ref{v08-07} below.

\section{Free, constrained, and moving endpoints}
\label{v07-02}

\begin{lemma}[Inner time variation]
\label{lem:inner-time-variation}
Let \(u\in C^1([a,b];\mathbb R^m)\) be stationary under compactly
supported changes of time for a \(C^1\) Lagrangian \(L\).
In particular, a \(C^1\) local minimum is stationary under those
changes. Then \(H(t)=L(t,u,u')-p(t)\cdot u'(t)\) is \(C^1\)
with \(H'=L_t(t,u,u')\).
\end{lemma}
\begin{proof}
Given a smooth compactly supported scalar \(\zeta\), set
\(t_s(t)=t+s\zeta(t)\) and \(u_s(t_s(t))=u(t)\).
For small \(|s|\), \(t_s'>0\), and \(t_s\) is a diffeomorphism
of the interval fixing its endpoints. Since \(t_s'\to1\)
uniformly and \(u'\) is uniformly continuous, \(u_s\to u\) in
\(C^1\). Thus these curves give admissible two-sided variations
at a local minimum. On the fixed reference interval the action
is \(\int_a^b L(t_s,u,u'/t_s')t_s'\,dt\).
The derivative of its integrand is uniformly bounded for small
\(|s|\) by compactness of the graph and continuity of \(DL\).
Stationarity at \(s=0\) gives
\[
 \int_a^b \bigl(L_t(t,u,u')\zeta+H\zeta'\bigr)\,dt=0 .
\]
Lemma~\ref{lem:variation-fundamental} applied to the continuous
functions \(H\) and \(L_t(t,u,u')\) gives \(H'=L_t(t,u,u')\).
\end{proof}

\begin{proposition}[Boundary term and transversality]
\label{prop:variation-boundary}
Let \(u\) be a \(C^1\) extremal with \(p'=q\), stationary under
compactly supported changes of time. Assume \(L\) is \(C^1\) on an
open neighborhood of \((t,u(t),u'(t))\) for all \(t\in[a,b]\),
including both endpoints. Consider a variation
\(u_s\) on intervals \([a_s,b_s]\), with \(u_0=u\) and \(a<b\).
Assume \(a_s,b_s\) are differentiable at zero and the rescaled
curves \(z_s(t)=u_s(t_s(t))\) are differentiable at zero in
\(C^1([a,b];\mathbb R^m)\), where \(t_s\) is the increasing
affine map of the reference interval to \([a_s,b_s]\).
Write \(\delta t\) for the changes of the endpoint times
and \(\delta z\) for the total changes of the endpoint states.
The first variation is the right endpoint value minus the
left endpoint value of
\[
 p\cdot\delta z+(L-p\cdot u')\,\delta t.
\]
At a \(C^1\) local minimum with independent endpoint constraints
given by smooth submanifolds of time--state space, this covector
annihilates their tangent spaces. At fixed time with free state
it gives \(p=0\). At fixed time with state constrained to a
submanifold it gives \(p\) annihilating its tangent space.
\end{proposition}
\begin{proof}
Set \(h(t)=\partial_s t_s(t)|_{s=0}\) and
\(w_0=\partial_s z_s|_{s=0}\).
Changing variables to the fixed interval and differentiating
as in Lemma~\ref{lem:inner-time-variation} gives
\[
 \left.\frac{d}{ds}\right|_{s=0}\!
   \int_a^b L(t_s,z_s,z_s'/t_s')t_s'\,dt
 =\int_a^b \bigl(L_t h+q\cdot w_0+p\cdot w_0'+Hh'\bigr)\,dt
 =[p\cdot w_0+Hh]_a^b.
\]
The second equality follows from \(p'=q\) and the inner time
variation identity. The endpoint values of \(w_0,h\) are
\(\delta z,\delta t\), respectively. At a local minimum the
interior Euler--Lagrange identity follows by compactly supported
state variations, as in Theorem~\ref{thm:first-variation}.
To realize any two endpoint tangent vectors, take differentiable
curves in the endpoint submanifolds with these derivatives and
write their values as \((a_s,z_a(s)),(b_s,z_b(s))\).
Let \(\theta=(t-a)/(b-a)\) and set
\[
 z_s(t)=u(t)+(1-\theta)(z_a(s)-u(a))
                         +\theta(z_b(s)-u(b)).
\]
These rescaled curves are differentiable in \(C^1\) and have the
prescribed endpoints. They are two-sided competitors in the
local-minimum test. Taking one endpoint curve constant at a time
yields the annihilator condition separately at the two ends.
\end{proof}

For example, at a free terminal time with terminal state held
fixed, the condition is \(L-p\cdot u'=0\) there.
See \cite[Section 2.3.5]{Liberzon} and
\cite[Section 2.2]{Dacorogna}.

\section{The Du Bois--Reymond identity and first integrals}
\label{v07-03}

\begin{proposition}[First integrals]
\label{prop:beltrami}
Let \(L\) be \(C^1\) near a curve \(u\) with
\(p=L_v(t,u,u')\) of class \(C^1\) and \(p'=L_z(t,u,u')\).
Assume \(u\) is \(C^2\), or that \(u\) is \(C^1\) and is
stationary under compactly supported changes of time.
Then
\[
 \frac{d}{dt}\bigl(L(t,u,u')-p(t)\cdot u'(t)\bigr)
     =L_t(t,u,u').
\]
In particular an autonomous \(L\) gives a constant
\(L-p\cdot u'\), and a coordinate absent from \(L\) gives a
constant corresponding component of \(p\).
\end{proposition}
\begin{proof}
For \(u\in C^2\), the derivative on the left is
\(L_t+(L_z-p')\cdot u'+(L_v-p)\cdot u''=L_t\).
For the \(C^1\) curve stationary under changes of time,
Lemma~\ref{lem:inner-time-variation} gives the derivative
directly. In the autonomous case \(L_t=0\).
For an ignorable coordinate the corresponding component of
\(p'=L_z\) vanishes.
\end{proof}

See \cite[Theorem 2.11]{Dacorogna} and
\cite[Section 2.3.4]{Liberzon}.


\section{Integral equality constraints and abnormal multipliers}
\label{v07-04}

\begin{theorem}[Isoperimetric multiplier rule]
\label{thm:v07-04-01}
Let \(J\) and \(G=(G_1,\ldots,G_k)\) be the integrals of \(C^1\)
integrands on the space of \(C^1\) curves with fixed endpoints,
near \(u\). Suppose \(u\) is a constrained local minimum for
\(G=G(u)\). Let \(V\) be the vector space of \(C^1\) variations
vanishing at the endpoints. There are multipliers
\((\lambda_0,\lambda)\ne0\) with \(\lambda_0\ge0\) such that
\[
 \lambda_0 DJ(u)[\eta]+\lambda\cdot DG(u)[\eta]=0
 \qquad(\eta\in V).
\]
If \(DG(u):V\to\mathbb R^k\) is surjective, one can set
\(\lambda_0=1\), and the resulting \(\lambda\) is unique.
A multiplier with \(\lambda_0=0\) is abnormal.
\end{theorem}
\begin{proof}
Derivatives of the integrals along each finite-dimensional family
are continuous by the same compact-graph estimate as in
Theorem~\ref{thm:first-variation}. If \(DG(u)\) is surjective,
choose \(\xi_j\in V\) with \(DG(u)\xi_j=e_j\).
For \(\eta\in\ker DG(u)\) apply the finite-dimensional implicit
function theorem to
\(G(u+t\eta+\sum_j s_j\xi_j)=G(u)\) at \(t=0,s=0\).
It gives \(s=s(t)\), with \(s(0)=s'(0)=0\).
The derivative of \(J\) along this feasible curve is zero,
so \(DJ(u)[\eta]=0\). Any \(\eta\) differs from
\(\sum_j(DG(u)[\eta])_j\xi_j\) by a vector in the kernel,
which proves the normal equation with
\(\lambda_j=-DJ(u)[\xi_j]\) and its uniqueness.
If the derivative is not onto, a nonzero covector annihilates
its range, yielding an abnormal multiplier. The argument in
this case supplies no normalization to \(\lambda_0=1\).
\end{proof}

Applying Lemma~\ref{lem:variation-fundamental} to the stationary
augmented integrand gives its Euler--Lagrange equation. See \cite[Section 2.5.1]{Liberzon}.

\section{Second variation and the Legendre conditions}
\label{v07-05}

\begin{theorem}[Legendre necessity in one dimension]
\label{thm:v07-05-01}
At a classical \(C^1\) weak local minimum with fixed endpoints
and \(C^2\) integrand, the second variation is
\[
 Q(\eta)=\int_a^b
 \bigl(L_{zz}[\eta,\eta]+2L_{zv}[\eta,\eta']
                         +L_{vv}[\eta',\eta']\bigr)\,dt\ge0
\]
for every \(C^1\) variation \(\eta\) vanishing at the endpoints.
All derivatives are evaluated along the curve. For every
interior \(t\), \(L_{vv}(t,u(t),u'(t))\) is positive semidefinite.
\end{theorem}
\begin{proof}
Twice differentiate \(J(u+s\eta)\) under the integral using
boundedness of the Hessian on the perturbed compact graphs.
The scalar second derivative at the minimum is nonnegative.
Fix \(t_0\in(a,b)\), a vector \(e\), and a nonzero smooth
scalar bump \(\psi\) supported in \((-1,1)\). Set
\(\eta_\epsilon(t)=\epsilon\psi((t-t_0)/\epsilon)e\).
After the substitution \(r=(t-t_0)/\epsilon\), divide
\(Q(\eta_\epsilon)\ge0\) by \(\epsilon\). The state and mixed
terms tend to zero by their extra factors \(\epsilon^2\) and
\(\epsilon\). Continuity gives the limiting inequality
\[
 L_{vv}(t_0,u(t_0),u'(t_0))[e,e]\int_{-1}^1|\psi'(r)|^2\,dr\ge0,
\]
and the integral is positive.
\end{proof}

\begin{theorem}[Legendre--Hadamard necessity]
\label{thm:v07-05-02}
Let \(\Omega\subset\mathbb R^d\) be open and \(u:\Omega\to\mathbb R^m\)
be \(C^1\). Suppose \(L(x,z,F)\) is \(C^2\) near the graph of
\((x,u,\nabla u)\), and that \(u\) is a weak local minimum under
compactly supported smooth variations of the finite energy
\(\int_\Omega L(x,u,\nabla u)\,dx\). Then for \(x\in\Omega\),
\[
 L_{FF}(x,u(x),\nabla u(x))[e\otimes b,e\otimes b]\ge0
 \qquad(e\in\mathbb R^m,\ b\in\mathbb R^d).
\]
\end{theorem}
\begin{proof}
On each compact support, the same differentiation gives a
nonnegative second variation with gradient variables replacing
velocity variables. For \(b\ne0\), \(\chi\in C_c^\infty(\Omega)\),
use \(\eta_k(x)=k^{-1}e\chi(x)\sin(k b\cdot x)\).
The state terms and the mixed terms tend to zero since
\(\eta_k=O(k^{-1})\) and \(\nabla\eta_k\) stays bounded there.
In the gradient term the remainder after
\(e\otimes b\,\chi\cos(k b\cdot x)\) is uniformly \(O(k^{-1})\).
The resulting inequality in the limit is
\[
 \frac12\int_\Omega \chi^2 L_{FF}(x,u,\nabla u)
                     [e\otimes b,e\otimes b]\,dx\ge0.
\]
To justify the limit, use \(\cos^2 s=(1+\cos 2s)/2\).
The integral of the oscillating term tends to zero for smooth
compactly supported coefficients by integration by parts in
direction \(b\). Approximation of the continuous compactly
supported coefficient in \(L^1\) by such smooth coefficients
extends the conclusion. If the displayed coefficient of
\(\chi^2\) were negative at a point, its continuity would give
a ball of negative values, contradicting a test supported in
that ball. The case \(b=0\) is immediate.
\end{proof}

This is a rank-one test. For example, the quadratic form
\(F\mapsto\det F\) on \(2\times2\) matrices vanishes on rank-one
matrices and takes both signs on general matrices. See
\cite[Section 2.6.1]{Liberzon} and \cite[Sections 3.4--3.5]{Dacorogna}.

\section{The strengthened Legendre--Jacobi sufficient condition}
\label{v07-06}

\begin{theorem}[Coercive scalar second variation]
\label{thm:v07-06-01}
Suppose a scalar \(C^2\) stationary curve has fixed endpoints
and \(C^3\) integrand near its graph. Set \(P=L_{vv}\),
\(B=L_{zv}\), \(C_0=L_{zz}\) along the curve and \(V=C_0-B'\).
Suppose \(P>0\) throughout \([a,b]\). Let \(v_0\) solve
\[
 (Pv_0')'=Vv_0,\qquad v_0(a)=0,\quad v_0'(a)=1.
\]
If \(v_0\) has no zero on \((a,b]\), there exists \(\kappa>0\)
such that the second variation satisfies
\(Q(\eta)\ge\kappa\|\eta\|_{H^1}^2\) for \(\eta\in H^1_0(a,b)\).
The curve is a strict weak local minimum, with the same quadratic
control up to a positive constant in a sufficiently small
\(C^1\) neighborhood.
\end{theorem}
\begin{proof}
All coefficients are continuous, \(P\) is \(C^1\) and
bounded below away from zero, and \(B'\) is continuous.
Integration by parts in the mixed term gives
\(Q(\eta)=\int_a^b(P|\eta'|^2+V|\eta|^2)\,dt\).
The Jacobi equation is the first-order linear system in
\((v,Pv')\) with continuous coefficients on the compact interval.
Its solutions with the given initial data exist and are unique
\cite[Chapter 7]{MathGenomeD13}.
Near \(a\), \(v_0>0\) to the right by the initial derivative,
and the absence of zeros gives \(v_0>0\) on \((a,b]\).
Let \(w\) solve the same equation with \(w(a)=1,w'(a)=0\).
For small \(\delta>0\), \(w\ge1/2\) on \([a,a+\delta]\).
On \([a+\delta,b]\), \(v_0\) has positive minimum and \(w\)
is bounded. Thus for sufficiently small \(\epsilon>0\),
\(v=v_0+\epsilon w\) is positive on all of \([a,b]\).

For a smooth endpoint-zero \(\eta\) write \(\eta=v\zeta\).
Expansion and integration by parts give
\[
 \begin{aligned}
 Q(v\zeta)
 &=\int_a^b\bigl(Pv^2|\zeta'|^2
       +2Pvv'\zeta\zeta'+(P(v')^2+Vv^2)\zeta^2\bigr)\,dt\\
 &=\int_a^b Pv^2|\zeta'|^2\,dt.
 \end{aligned}
\]
Indeed \(2Pvv'\zeta\zeta' = Pvv'(\zeta^2)'\), the boundary term
is zero, and \((Pvv')'=P(v')^2+Vv^2\).
By density and multiplication by \(v\) or \(1/v\), bounded
operators on \(H^1(a,b)\), the identity extends to \(H^1_0\).
For zero-endpoint \(\zeta\) the fundamental theorem of calculus
and Cauchy--Schwarz give
\(\|\zeta\|_{L^2}\le(b-a)\|\zeta'\|_{L^2}\).
The positive lower bound of \(Pv^2\) and the upper bounds on
\(v,v'\) now give \(Q(\eta)\ge\kappa\|\eta\|_{H^1}^2\).

For small \(C^1\) perturbations, Taylor's formula in the state
and velocity, the vanishing first variation, and uniform
continuity of the second derivatives on a compact neighborhood
give
\[
 J(u+\eta)-J(u)=\tfrac12 Q(\eta)+r(\eta),\qquad
 |r(\eta)|\le \omega(\|\eta\|_{C^1})\|\eta\|_{H^1}^2,
\]
with \(\omega(s)\to0\) as \(s\downarrow0\).
Choosing \(\omega<\kappa/4\) proves strict minimality and the
claimed control.
\end{proof}

See \cite[Section 2.6.2]{Liberzon}. The positive lower bound on
the quadratic form is the step that permits Taylor's remainder
to be absorbed.

\section{Needle variations and corner conditions}
\label{v07-07}

\begin{theorem}[Weierstrass excess]
\label{thm:v07-07-01}
Suppose \(u\) is \(C^1\) and a strong local minimum among
continuous piecewise \(C^1\) curves with fixed endpoints.
For \(t_0\in(a,b)\) and a tested velocity \(w\in\mathbb R^m\), assume
\(L\) is \(C^1\) on an open neighborhood of the reference graph and
of the graph obtained by inserting the velocity \(w\) near \(t_0\),
including the velocity segments joining these graphs. Then
\[
 L(t_0,u(t_0),w)-L(t_0,u(t_0),u'(t_0))
                  -p(t_0)\cdot(w-u'(t_0))\ge0.
\]
\end{theorem}
\begin{proof}
Choose a fixed compact interval to the right of \(t_0\) and
a smooth function \(\rho\) supported there, of integral one.
For small \(\epsilon>0\) let \(I_\epsilon=[t_0,t_0+\epsilon]\).
Put \(A_\epsilon=\int_{I_\epsilon}(w-u'(t))\,dt\) and define
the endpoint-zero perturbation \(\eta_\epsilon\) by integrating
from \(a\) the derivative
\[
 \eta_\epsilon'(t)=(w-u'(t)){\boldsymbol 1}_{I_\epsilon}(t)
                       -A_\epsilon\rho(t).
\]
It is continuous and piecewise \(C^1\), and
\(\|\eta_\epsilon\|_\infty=O(\epsilon)\).
The first variation of \(J\) at \(u\) in this direction is zero.
One can either use two-sided perturbations by a scalar
multiple of \(\eta_\epsilon\), or integrate by parts with
\(p'=q\) and the endpoint conditions.
Off \(I_\epsilon\) both the perturbation and its derivative
are \(O(\epsilon)\). Subtracting the first variation there
leaves an integral of \(o(\epsilon)\) by uniform continuity
of the first derivatives. On \(I_\epsilon\) the new velocity
equals \(w\) and \(\eta_\epsilon=O(\epsilon)\).
Divide the remaining action difference by \(\epsilon\) and
let \(\epsilon\to0\). Continuity gives exactly the displayed
expression, since the state term in the first variation
contributes \(O(\epsilon^2)\) on \(I_\epsilon\).
Strong local minimality gives nonnegative action differences
for these competitors.
\end{proof}

\begin{proposition}[Weierstrass--Erdmann corner conditions]
\label{prop:v07-07-02}
Let a continuous piecewise \(C^2\) curve have a single interior
corner at \(c\), be stationary on the two pieces, and be a
strong local minimum for a class allowing piecewise \(C^1\)
perturbations at the corner. Suppose \(L\) is \(C^1\) on an open
neighborhood of the two closed one-sided state--velocity graphs,
including both limiting velocities at \(c\). If the joined state at \(c\) is
free, the one-sided momenta agree. If the corner time is also
free, the one-sided values of \(L-p\cdot u'\) agree.
\end{proposition}
\begin{proof}
On each \(C^2\) piece, Proposition~\ref{prop:beltrami} gives
\(H'=L_t\), so the first variation under any compactly
supported time change vanishes by integration by parts.
Apply Proposition~\ref{prop:variation-boundary} on the two
pieces, using the same corner time and total corner state
in each. Their contributions at the join sum to
\[
 (p(c^-)-p(c^+))\cdot\delta z+
 \bigl((L-p\cdot u')(c^-)-(L-p\cdot u')(c^+)\bigr)\delta t.
\]
Independently prescribed changes \(\delta z,\delta t\)
are realized by linear interpolation of endpoint data on the
two rescaled pieces, with a common join. The exterior endpoints
stay fixed. For a fixed-time corner one varies only \(\delta z\).
Both signs of each allowed direction are possible, so their
coefficients vanish as claimed.
\end{proof}

See \cite[Sections 3.1.1--3.1.2]{Liberzon}.

\section{The planar isoperimetric inequality and graph area}
\label{v07-08}

\begin{lemma}[Periodic Wirtinger inequality]
\label{lem:v07-08-01}
For a scalar \(C^1\) function \(f\) of period \(T>0\) with zero
mean,
\[
 \int_0^T |f|^2\,dt
   \le \left(\frac{T}{2\pi}\right)^2\int_0^T |f'|^2\,dt.
\]
Equality holds exactly for the first-harmonic span
\(A\cos(2\pi t/T)+B\sin(2\pi t/T)\), including zero.
\end{lemma}
\begin{proof}
For \(c_k=T^{-1}\int_0^T f(t)e^{-2\pi ikt/T}\,dt\),
the derivative has coefficients \((2\pi ik/T)c_k\) by
periodic integration by parts. Parseval for the periodic Fourier
basis \cite[Unit U1]{MathGenomeP102} therefore gives the
inequality term by term because
\(c_0=0\) and \(k^2\ge1\) for \(k\ne0\).
Equality forces \(c_k=0\) for \(|k|\ne1\), which gives the
stated formula in \(L^2\) and then everywhere by continuity.
\end{proof}

\begin{theorem}[Planar isoperimetric inequality]
\label{thm:v07-08-02}
Let a regular \(C^1\) simple closed plane curve, parametrized by a
\(C^1\) embedding of the circle, bound area \(A\) and have length \(T\).
Then \(4\pi A\le T^2\).
Equality holds exactly for circles.
\end{theorem}
\begin{proof}
Reparametrize the positively oriented curve by arc length
\(\gamma:[0,T]\to\mathbb R^2\) and subtract its mean.
The speed is one \cite[Section 9]{MathGenomeD04}.
The regular \(C^1\) Jordan boundary has local \(C^1\) graphs with the
bounded interior on one side \cite[Chapter 21]{MathGenomeD13}.
Thus Green's theorem \cite[Section 12]{MathGenomeD04}, applied to
\((P,Q)=(-y/2,x/2)\), gives
\(A=\tfrac12\int_0^T J_0\gamma\cdot\gamma'\,dt\), where \(J_0(x,y)=(-y,x)\).
Subtracting a constant \(m\) from the parametrization leaves the boundary
integral unchanged because \(J_0m\cdot(\gamma(T)-\gamma(0))=0\).
Cauchy--Schwarz and Wirtinger in each coordinate give
\[
 A\le\tfrac12\|\gamma\|_{L^2}\|\gamma'\|_{L^2}
    \le \frac{T}{4\pi}\|\gamma'\|_{L^2}^2
    =\frac{T^2}{4\pi}.
\]
If equality holds, \(A>0\) and equality in Cauchy--Schwarz
gives \(\gamma'=cJ_0\gamma\) with a fixed \(c>0\).
Equality in Wirtinger fixes \(c=2\pi/T\). The solution of this
linear equation is a circle with constant nonzero radius.
Conversely a circle has the indicated equality.
\end{proof}

\begin{proposition}[First variation of graph area]
\label{prop:v07-08-03}
Let \(\Omega\subset\mathbb R^d\) be bounded and open.
For \(u\in W^{1,1}(\Omega)\), the graph-area functional
\[
 \mathcal A(u)=\int_\Omega\sqrt{1+|\nabla u|^2}\,dx
\]
has first variation along \(\phi\in C_c^\infty(\Omega)\)
equal to
\[
 \int_\Omega\frac{\nabla u\cdot\nabla\phi}
                   {\sqrt{1+|\nabla u|^2}}\,dx .
\]
A local minimum with these two-sided tests therefore satisfies
the weak minimal-surface equation.
\end{proposition}
\begin{proof}
The derivative of \(F\mapsto\sqrt{1+|F|^2}\) is
\(F/\sqrt{1+|F|^2}\), whose norm is at most one. The difference
quotients are bounded by \(|\nabla\phi|\), which is integrable.
Dominated convergence gives the formula, and the scalar first
derivative at a two-sided minimum vanishes.
\end{proof}

See \cite[Chapters 5--6]{Dacorogna} and
\cite[Section 2.1.1]{Liberzon}.
\chapter{The Sobolev Direct Method and Weak Equations}
\label{chap:sobolev-direct}

Let \(\Omega\subset\R^d\) be a bounded open set, \(d\ge1\), and
let the maps under consideration take values in \(\R^m\), \(m\ge1\).
The notation \(Du\) denotes the \(m\)-by-\(d\) matrix of weak first
derivatives. Matrix pairings are Frobenius pairings. A Carath\'eodory
integrand is measurable in the base point and continuous in its
finite-dimensional arguments for almost every base point. We may
choose a common exceptional null set and define a harmless
integrand there. Composing with measurable finite-dimensional
fields gives a measurable function: round each field to a
countable rational grid and pass to the pointwise limit by
continuity. The basic variational existence and stationarity
arguments are treated separately: a lower bound used in the
existence proof and an upper derivative bound used in the
stationarity proof have different roles.

\section{Admissible classes, zero boundary values, and compactness}
\label{v08-01}

For \(1\le p<\infty\), \(W^{1,p}(\Omega;\R^m)\) consists of \(L^p\)
classes whose distributional first derivatives are in \(L^p\).
We use the norm \((\norm{u}_p^p+\norm{Du}_p^p)^{1/p}\).
Define \(W^{1,p}_0(\Omega;\R^m)\) as the closure of
\(C^\infty_c(\Omega;\R^m)\) in this norm. For prescribed \(u_0\) in
\(W^{1,p}\), the affine class is \(u_0+W^{1,p}_0\). On a bounded
Lipschitz open set the ordinary trace theorem identifies
\(W^{1,p}_0\) with the kernel of the trace into \(L^p(\partial\Omega)\).
We use this identification only in that setting
\cite[Chapter 9]{Brezis}\cite[Unit U5, ordinary trace theorem]{MathGenomeP102}.

\begin{proposition}\label{prop:sobolev-affine}
The space \(W^{1,p}\) is Banach, and it is reflexive for \(1<p<\infty\).
The prescribed affine class is norm closed and weakly closed.
Zero extension of a function in \(W^{1,p}_0(\Omega)\) belongs to
\(W^{1,p}(\R^d)\) with the same norm, with zero-extended weak derivatives.
If \(\Omega\) lies in a strip \(a<x_1<b\), then
\[
 \norm{u}_p\le (b-a)\norm{\partial_1u}_p
 \quad (u\in W^{1,p}_0(\Omega;\R^m)).
\]
\end{proposition}

\begin{proof}
The graph \(u\mapsto(u,Du)\) is a closed linear subspace of the
product \(L^p\) space. Indeed, if \(u_j\to u\) and \(Du_j\to V\)
in the respective \(L^p\) norms, their distributional identities
pass to the limit by testing against bounded compactly supported
smooth functions. This proves closedness, completeness, and, for
\(1<p<\infty\), reflexivity by \(L^p\) reflexivity and the closed
subspace property. The affine class is closed by its definition,
and weakly closed by convex separation
(Section~\ref{v02-03}).

For a smooth compactly supported function, extension by zero has
the indicated norm and derivatives. Pass to the limit in a defining
sequence for \(W^{1,p}_0\), using the closed graph just proved for
\(\R^d\). For such a smooth function and \(x'\in\R^{d-1}\),
\[
 \lvert u(t,x')\rvert
 \le \int_a^b \lvert\partial_1u(s,x')\rvert\,ds .
\]
H\"older bounds the \(p\)-th power by \((b-a)^{p-1}\) times the
line integral of \(\lvert\partial_1u\rvert^p\), including \(p=1\).
Integrate over \(t\in(a,b)\) and \(x'\), then use density. No
connectedness or boundary regularity is needed for this estimate.
Compare \cite[Theorem 1.51]{Dacorogna}.
\end{proof}

\begin{theorem}\label{thm:sobolev-compact-affine}
For \(1\le p<\infty\), bounded sets of \(W^{1,p}(\Omega)\) are
relatively compact in \(L^p(\Omega)\) if \(\Omega\) is bounded
Lipschitz. On an arbitrary bounded open \(\Omega\), the same
conclusion holds for bounded sets in any fixed class
\(u_0+W^{1,p}_0(\Omega)\).
\end{theorem}

\begin{proof}
In the first case apply the bounded \(W^{1,p}\) extension for
Lipschitz domains and multiply by a smooth cutoff equal to one on
a neighborhood of \(\overline\Omega\). The resulting functions
are bounded in \(W^{1,p}(\R^d)\) and supported in one compact set.
In the second case subtract \(u_0\) and use zero extension.
For either resulting family, translation satisfies
\[
 \norm{v(\,\cdot-h)-v}_p\le |h|\norm{Dv}_p.
\]
For a smooth function integrate its derivative along the segment
and use the integral triangle inequality. The general inequality
follows by density in \(W^{1,p}(\R^d)\).

Let \(\rho_\delta\) be a nonnegative smooth mollifier supported in
the radius-\(\delta\) ball with integral one. The translation
estimate gives \(\norm{\rho_\delta*v-v}_p\le
\delta\sup\norm{Dv}_p\), uniformly on the family. For fixed
\(\delta>0\), the convolutions have common compact support and
uniform bounds on their values and first derivatives. These
bounds follow from H\"older applied to \(\rho_\delta\) and to
\(D\rho_\delta\), with the \(L^\infty\) kernel bound at \(p=1\).
Arzel\`a--Ascoli gives relative compactness in the uniform norm
there and hence in \(L^p\). For any prescribed \(L^p\) error, first
choose \(\delta\) small and then a finite net for the convolved
family. Thus the original family is totally bounded in the
complete \(L^p\) space. Restriction and addition of \(u_0\)
give the conclusions. See \cite[Section 1.4]{Dacorogna} and
\cite[Unit U5]{MathGenomeP102}.
\end{proof}

\section{Strong--weak lower semicontinuity of integral energies}
\label{v08-02}

This section supplies the step from convexity only in the gradient
variables to lower semicontinuity of the energy. The same argument
allows the different integrability exponents of the minors used
in Section~\ref{v08-06}. Compare
\cite[Theorem 3.3, Remark 3.4, and Section 3.3.2]{Dacorogna} and
\cite[Chapter VIII, Section 1]{EkelandTemam}.

\begin{theorem}\label{thm:strong-weak-integral}
Let \((S,\mu)\) have finite measure, let \(1\le q<\infty\), and let
\(1<r_i<\infty\) for \(1\le i\le k\). Suppose
\(s_n\to s\) strongly in \(L^q(S;\R^m)\) and
\(z_{i,n}\rightharpoonup z_i\) weakly in \(L^{r_i}(S;\R^{d_i})\).
Write \(z=(z_1,\ldots,z_k)\). Let \(F(x,s,z)\) be finite
Carath\'eodory and convex in the entire \(z\) vector for each
\(s\) at almost every \(x\). Assume
\[
 F(x,s,z)\ge -a(x)-C|s|^q,\qquad a\in L^1(S),\quad C\ge0,
\]
uniformly in \(s,z\). Then
\[
 \int_S F(x,s(x),z(x))\,d\mu
 \le \liminf_n\int_S F(x,s_n(x),z_n(x))\,d\mu .
\]
The energy may take the value \(+\infty\).
\end{theorem}

\begin{proof}
Add \(a(x)+C|s|^q\) to obtain a nonnegative integrand \(G\).
The integrals of the added terms converge, because the \(L^q\)
norms converge. It is enough to prove the assertion for \(G\).
Pass to a subsequence realizing the lower limit, and then to a
further subsequence with \(s_n(x)\to s(x)\) almost everywhere.
If the lower limit is \(+\infty\), there is nothing further to
prove. We describe the proof for a finite lower limit.

Fix a rational vector \(b\) in the full fiber space
\(\R^{d_1+\cdots+d_k}\) and a positive \(\epsilon=1/j\).
For each \(x,s\), let \(w(x,s)\) minimize
\[
 G(x,s,y)+\frac{|y-b|^2}{2\epsilon}
\]
over the fiber. This continuous strongly convex function has
compact sublevel sets since \(G\ge0\), so its minimizer exists
and is unique. When \(s=s(x)\) is measurable, the minimizer is
measurable as a function of \(x\). To see this, the infimum
equals the infimum over a fixed countable dense set. For each
integer \(l\), choose the first member of that set with value
within \(1/l\) of the infimum. The selection is measurable
because all the values are measurable. If \(Q\) denotes the
displayed objective, the strong-convexity midpoint inequality
and minimality imply
\[
 Q(y)-Q(w)\ \ge\ \frac{|y-w|^2}{4\epsilon}.
\]
Consequently the chosen approximate minimizers converge to
\(w(x,s(x))\).

For a fixed \(x\) with the Carath\'eodory properties, the
minimizer varies continuously with \(s\), and also with the
center \(b\) if arbitrary real centers are allowed. In fact,
for \(s_l\to s\) and \(b_l\to b\), comparison at \(y=b_l\)
bounds \(|w(x,s_l)-b_l|^2/(2\epsilon)\) by \(G(x,s_l,b_l)\).
This is a locally uniform bound by continuity. Any cluster
point minimizes the limiting objective, by comparing with each
fixed competitor and passing to the limit. Uniqueness proves
convergence of the whole sequence.

With \(w=w(x,s)\), define the affine function of \(y\)
\[
 \ell_{b,j}(x,s,y)
   =G(x,s,w)+\epsilon^{-1}\langle b-w,y-w\rangle.
\]
It is a lower bound for \(G\). Indeed minimality at
\(w+\theta(y-w)\), convexity of \(G\), and expansion of the
square give, for \(0<\theta<1\),
\[
 0\le G(x,s,y)-G(x,s,w)
       +\epsilon^{-1}\langle w-b,y-w\rangle
       +\frac{\theta}{2\epsilon}|y-w|^2 .
\]
Let \(\theta\downarrow0\). The countable affine family recovers
\(G\) by its supremum. At a fixed \(x,s,y\) and \(\epsilon\),
take rational \(b\to y\). The coefficients of the affine
function are continuous in \(b\) by the minimizer argument.
With center \(y\), the minimizing \(w\) satisfies
\(|w-y|^2\le 2\epsilon G(x,s,y)\), so \(w\to y\) as
\(\epsilon\downarrow0\). The affine value at \(y\) is then at
least \(G(x,s,w)\), which tends to \(G(x,s,y)\). This proves
the supremum identity.

Enumerate the affine functions, including zero. For a finite
initial list, select measurably an index attaining the maximum
at \((x,s(x),z(x))\), taking the first in a tie. Evaluate the
selected affine function at \((x,s_n(x),z_n(x))\). Write its
coefficients as \(c_n(x),b_{i,n}(x)\), so its value is
\(c_n+\sum_i\langle b_{i,n},z_{i,n}\rangle\). The coefficients
are measurable and converge almost everywhere to the
coefficients \(c,b_i\) evaluated at \(s\). There are increasing
measurable sets \(S_M\) covering almost all of \(S\) on which
these coefficients, for every \(n\) and the limit, are bounded
by \(M\). For example use the sublevel sets of the countable
supremum of their norms.

On \(S_M\), the intercepts converge in \(L^1\) and the slopes in
\(L^{r_i'}\) by dominated convergence. Extend the restricted
coefficients by zero. Then
\[
 \int_{S_M}\langle b_{i,n},z_{i,n}\rangle
   \longrightarrow \int_{S_M}\langle b_i,z_i\rangle .
\]
The difference with \(b_i\) as test is controlled by H\"older
and boundedness of the weakly convergent sequence. The remaining
term converges by weak convergence. Since \(G\) dominates the
selected affine function on \(S_M\) and is nonnegative outside
it, its lower integral limit is at least the integral over
\(S_M\) of the finite maximum at the limiting arguments. Let
\(M\to\infty\) and then enlarge the finite list. Each passage
uses monotone convergence, as the maximum includes zero.
The countable supremum identity yields the required integral.
\end{proof}

\section{The direct method with convex gradient dependence}
\label{v08-03}

\begin{theorem}\label{thm:gradient-convex-existence}
Let \(\Omega\) be bounded open, \(p>q\ge1\),
\(u_0\in W^{1,p}(\Omega;\R^m)\), and let \(F(x,s,Z)\) be finite
Carath\'eodory and convex in \(Z\). Suppose there are
\(\alpha>0\), \(C\ge0\) and \(a\in L^1(\Omega)\) such that
\[
 F(x,s,Z)\ge\alpha |Z|^p-C|s|^q-a(x).
\]
If \(J(u)=\int_\Omega F(x,u,Du)\,dx\) is finite at some point of
\(u_0+W^{1,p}_0\), then \(J\) attains a finite minimum in that
class.
\end{theorem}

\begin{proof}
On the class in question, Proposition~\ref{prop:sobolev-affine}
applied to \(u-u_0\) gives
\(\norm{u}_p\le C_0+C_1\norm{Du}_p\). Since \(|\Omega|<\infty\),
H\"older bounds \(\norm{u}_q\) by a fixed multiple of this
quantity. Since \(q<p\), for any \(\eta>0\) one has
\(t^q\le\eta t^p+C_\eta\) on \([0,\infty)\), as is seen by
splitting at a large \(t\). Absorb the negative \(q\)-term to get
\[
 J(u)\ge \frac{\alpha}{2}\norm{Du}_p^p-C_2 .
\]
The finite competitor and this lower bound make the infimum
finite and provide a bounded minimizing sequence. By reflexivity
and Section~\ref{v03-01}, pass to a weakly convergent subsequence.
The affine class is weakly closed. The compactness theorem gives
a further subsequence converging strongly in \(L^p\), and thus
in \(L^q\), to the same limit. Apply
Theorem~\ref{thm:strong-weak-integral} with fiber \(Du_n\) in
\(L^p\). The resulting inequality proves attainment. Compare
\cite[Theorems 3.1 and 3.3, and Exercise 3.3.4]{Dacorogna}.
\end{proof}

\begin{corollary}\label{cor:dirichlet-model}
For scalar \(u\), bounded open \(\Omega\), \(u_0\in H^1(\Omega)\)
and \(b\in L^2(\Omega)\), the energy
\(\tfrac12\int_\Omega |Du|^2-\int_\Omega bu\) has a unique
minimizer on \(u_0+H^1_0(\Omega)\). It is characterized by
\[
 \int_\Omega Du\cdot D\varphi=\int_\Omega b\varphi
 \qquad(\varphi\in H^1_0(\Omega)).
\]
\end{corollary}

\begin{proof}
By the zero-boundary Poincar\'e estimate the gradient norm is
equivalent to the \(H^1\) norm on \(H^1_0\). Writing \(u=u_0+v\),
the energy differs by a constant from a positive quadratic
functional in \(Dv\) plus a bounded linear functional of \(v\).
It is coercive and weakly lower semicontinuous by convexity and
continuity. Apply Section~\ref{v03-02}. Equality of gradients
forces equality of members of the affine class, so the quadratic
energy is strictly convex there and the minimizer is unique.
Expanding the energy on \(u+t\varphi\) gives the displayed
identity by the vanishing linear coefficient. Conversely that
identity gives \(J(u+\varphi)-J(u)=\norm{D\varphi}_2^2/2\).
For \(b=0\) this is the weak Dirichlet principle.
\end{proof}

\section{Uniqueness and the scope of the existence hypotheses}
\label{v08-04}

\begin{proposition}\label{prop:sobolev-unique}
On a bounded open set, consider the affine class
\(u_0+W^{1,p}_0\), \(1\le p<\infty\), and an energy with
integrable negative part at every admissible map. A finite
minimizer is unique if \(F\) is strictly convex in \((s,Z)\)
for almost every base point. The same holds if \(F\) is
independent of \(s\) and strictly convex in \(Z\).
\end{proposition}

\begin{proof}
If two finite minimizers differ as \(L^p\) classes, their value
arguments differ on a set of positive measure. Strict joint
convexity makes the energy of their midpoint strictly smaller:
the nonnegative difference between the mean integrand and the
midpoint integrand is positive on this set. Its integral is
well defined: the negative part of the midpoint integrand is
integrable by hypothesis and its positive part is bounded by
the sum of the two finite integrands' positive parts.
In the gradient-only case, equality of the
gradients would imply equality of the two functions by
Poincar\'e on their zero-boundary difference. Thus the same
argument applies on a positive-measure set of unequal gradients.
\end{proof}

\begin{example}[Three limitations]\label{ex:direct-limitations}
On \((0,1)\), consider \(W^{1,1}\) functions with endpoint values
\(0,1\). The convex-gradient functional
\(\int_0^1(1+x)|u'(x)|\,dx\) has infimum \(1\), which it does not
attain. Every competitor has \(\int|u'|\ge1\), and
\(\int x|u'|>0\) since \(u'\) cannot vanish almost everywhere.
The functions \(u_n(x)=\min(nx,1)\) give values \(1+1/(2n)\).
Here the gradient growth is linear and coercive in \(W^{1,1}\)
on the affine class.

On \(H^1_0(0,1)\), the continuous-valued integral
\(\int_0^1 e^{-u(x)}\,dx\) is positive at every competitor but
has infimum zero. Indeed take \(u_n(x)=nx(1-x)\) and apply
dominated convergence. This energy is convex, but it does not
bound the norms of minimizing sequences.

Finally, let \(u_n\) be the continuous tent function on each
of the \(n\) consecutive intervals of length \(1/n\), with
slopes \(+1,-1\) on their two halves and value zero at their
endpoints. Then \(u_n\to0\) uniformly and \(u_n'\rightharpoonup0\)
in \(L^4(0,1)\). The latter follows from distributional
convergence by integration by parts and uniform \(L^4\)
boundedness, by density in \(L^{4/3}\).
Yet \(\int((u_n')^2-1)^2=0\) whereas the limiting energy is
\(1\). This isolates the failure of lower semicontinuity for a
nonconvex gradient integrand. Compare
\cite[Examples 3.7--3.10]{Dacorogna}.
\end{example}

\section{Minors: integral identities and weak continuity}
\label{v08-05}

For ordered sets of distinct row and column indices
\(\alpha=(a_1,\ldots,a_j)\), \(I=(i_1,\ldots,i_j)\), let
\(M_{\alpha,I}(Z)\) be the corresponding \(j\)-by-\(j\) minor
of \(Z\), with the indices in increasing order. We write
\(\mathcal M_j(Z)\) for the vector of all \(j\)-minors.
The first vector is the list of matrix entries.
See \cite[Lemma 3.26 and Remark 3.23]{Dacorogna}.

\begin{lemma}\label{lem:minor-divergence}
Fix \(1\le j\le \min(m,d)\). For a smooth \(u\), write \(C_l(Du)\)
for the cofactor in the first row and \(l\)-th column of
\(M_{\alpha,I}(Du)\), with \(C_1=1\) when \(j=1\). Then
\[
 \sum_{l=1}^j\partial_{i_l}C_l(Du)=0,\qquad
 M_{\alpha,I}(Du)
   =\sum_{l=1}^j\partial_{i_l}\bigl(u^{a_1}C_l(Du)\bigr).
\]
The second identity holds distributionally for \(u\in W^{1,p}\)
whenever \(p\ge j\).
\end{lemma}

\begin{proof}
Expand the cofactor by permutations of the columns other than
\(l\), so each term corresponds to a permutation \(\pi\) with
\(\pi(1)=l\). In its derivative \(\partial_{i_l}\), consider
the term obtained by differentiating the factor belonging to
row \(a_r\), \(r\ge2\). Swap the two assigned columns
\(\pi(1)\) and \(\pi(r)\) to obtain a term in the derivative
of the corresponding other cofactor. The two terms have the
same undifferentiated factors and equal mixed derivatives
with opposite permutation signs. This pairs every term, proving
the cofactor identity. The \(j=1\) identity is the derivative
of a constant. Expand the divergence of the displayed products
and use the cofactor identity and expansion in the first row.

For the distributional statement, approximate by smooth maps
in \(W^{1,p}\) locally, using convolution on a neighborhood of
the compact support of each test. A product of \(r\) factors
converging in \(L^p\), with \(r\le j\le p\), converges in
\(L^{p/r}\) by telescoping the product and using H\"older.
This applies to each minor and to each product \(u^{a_1}C_l\).
Thus both sides of the divergence identity pass to the limit
in distributions.
\end{proof}

\begin{proposition}\label{prop:fixed-trace-minors}
If \(p\ge j\), \(u,v\in W^{1,p}(\Omega;\R^m)\), and
\(u-v\in W^{1,p}_0(\Omega;\R^m)\), then
\[
 \int_\Omega \mathcal M_j(Du)\,dx
 =\int_\Omega \mathcal M_j(Dv)\,dx .
\]
\end{proposition}

\begin{proof}
First assume \(u-v\) is smooth and compactly supported. The
difference of the fluxes in Lemma~\ref{lem:minor-divergence}
vanishes outside that compact support. Test the divergence
identity with a smooth function equal to one on a neighborhood
of the support. The pairing with the flux difference is zero,
and the left side is the full integral of the minor difference.
In the general case approximate \(u-v\) in \(W^{1,p}\) by
smooth compactly supported differences and apply the result to
each corresponding pair. Strong \(W^{1,p}\) convergence gives
strong \(L^{p/j}\) convergence of the minors by the same product
estimate. Since \(\Omega\) has finite measure and \(p/j\ge1\),
their integrals converge.
\end{proof}

\begin{theorem}\label{thm:weak-minors}
Suppose \(u_n\rightharpoonup u\) in \(W^{1,p}(\Omega;\R^m)\)
and \(u_n\to u\) in \(L^p\). For each \(1\le j<p\) with
\(j\le\min(m,d)\),
\[
 \mathcal M_j(Du_n)\rightharpoonup\mathcal M_j(Du)
 \quad\text{in }L^{p/j}.
\]
The strong convergence premise holds for such a weakly
convergent sequence in a bounded Lipschitz domain, or in a
fixed affine zero-boundary class on any bounded open domain.
\end{theorem}

\begin{proof}
For \(j=1\) the assertion is part of weak Sobolev convergence.
Induct on \(j\). The order \(j-1\) cofactors converge weakly in
\(L^{p/(j-1)}\), and the order \(j\) minors are bounded in
\(L^{p/j}\) by H\"older. For each smooth compactly supported
test \(\varphi\), write the distributional pairing from
Lemma~\ref{lem:minor-divergence}. The terms involving
\(u_n^{a_1}-u^{a_1}\) tend to zero, since
\[
 \left|\int_\Omega (u_n^{a_1}-u^{a_1})C_l(Du_n)
                              \partial_{i_l}\varphi\right|
 \le \norm{u_n-u}_p\norm{C_l(Du_n)}_{p/(j-1)}
                        \norm{\partial_{i_l}\varphi}_{p/(p-j)} .
\]
For the terms with \(u^{a_1}\) fixed, the test
\(u^{a_1}\partial_{i_l}\varphi\) belongs to
\(L^{p/(p-j+1)}\) since \(p\ge p/(p-j+1)\) and the support has
finite measure. They therefore converge by the induction
hypothesis. This proves distributional convergence of the
minors to the stated limit. The conjugate exponent
\((p/j)'\) is finite since \(j<p\). Approximate each test
in \(L^{(p/j)'}\) by smooth compactly supported tests and use
the uniform \(L^{p/j}\) bound, proving weak convergence.
Finally, compactness from Theorem~\ref{thm:sobolev-compact-affine}
gives a strongly convergent subsequence of every subsequence,
and its limit must be \(u\) by weak convergence in \(L^p\).
This yields strong convergence of the full sequence under
the stated compactness hypotheses.
\end{proof}

\section{An existence theorem for polyconvex integrands}
\label{v08-06}

Put \(r=\min(m,d)\). A polyconvex representation for the
integrand has the form
\[
 F(x,s,Z)=G(x,s,\mathcal M_1(Z),\ldots,\mathcal M_r(Z)),
\]
where
\(G\) is convex in the full vector of independent minor
coordinates. Convexity in \(Z\) gives such a representation
by ignoring the higher coordinates.

\begin{theorem}\label{thm:polyconvex-existence}
Suppose \(\Omega\) is bounded Lipschitz,
\(p>\max(r,q)\) with \(q\ge1\), and \(G\) is a finite
Carath\'eodory polyconvex representation as above. Suppose
uniformly in the independent coordinates \(z_1,\ldots,z_r\),
\[
 G(x,s,z_1,\ldots,z_r)
       \ge \alpha|z_1|^p-C|s|^q-a(x),
 \qquad \alpha>0,\quad C\ge0,\quad a\in L^1(\Omega).
\]
If the energy has a finite competitor in
\(u_0+W^{1,p}_0(\Omega;\R^m)\), it has a finite minimizer there.
This includes the \(2\)-by-\(2\) representation in the matrix
and determinant under \(p>\max(2,q)\).
\end{theorem}

\begin{proof}
The coercive estimate and the weak and strong subsequence
extractions are those of Theorem~\ref{thm:gradient-convex-existence};
they only used the displayed lower bound, the fixed affine
class and Sobolev compactness. For each order \(j\), use
Theorem~\ref{thm:weak-minors} to obtain weak convergence of the
minor vector in \(L^{p/j}\), where \(p/j>1\). In
Theorem~\ref{thm:strong-weak-integral} use the full list of
these fibers and strong \(L^q\) convergence of \(u_n\).
Its conclusion makes the admissible limit a minimizer.
Compare \cite[Theorem 3.22 and Remark 3.23]{Dacorogna}.
\end{proof}

\begin{proposition}\label{prop:polyconvex-test}
For a fixed base point and fixed value argument, a finite
convex minor representation implies the quasiconvex
test inequality on every cube \(Q\):
\[
 F(Z)\le \frac1{|Q|}\int_Q F(Z+D\varphi)\,dx
 \qquad(\varphi\in C^\infty_c(Q;\R^m)).
\]
\end{proposition}
\begin{proof}
Compare the affine map with derivative \(Z\) and that map
plus \(\varphi\). Proposition~\ref{prop:fixed-trace-minors}
shows that the mean of the full minor vector of
\(Z+D\varphi\) equals the minor vector of \(Z\).
A subgradient of the finite convex representation at this
mean exists by Section~\ref{v02-07}. Integrate its affine
support inequality. The linear term has mean zero, giving
the claimed bound.
\end{proof}

For example in two dimensions,
\(F(Z)=|Z|^p+(\det Z-\gamma)^2\), \(p>2\), \(\gamma>0\),
is polyconvex by its displayed representation. On \(Z=tI\)
its increment above \(F(0)\) is
\(2^{p/2}|t|^p-2\gamma t^2+t^4<0\) for small nonzero \(t\).
Evenness along this line rules out convexity. On a bounded
domain with zero boundary values the zero map is a finite
competitor. For \(2<p<4\) the energy need not be finite
at every Sobolev competitor, which is why attainment was
formulated for an extended-valued energy. The full theory
characterizing lower semicontinuity by quasiconvexity and the
critical minor exponent \(p=r\) have separate hypotheses.

\section{Weak Euler--Lagrange equations and test classes}
\label{v08-07}

For a differentiable integrand in its last two variables,
write \(F_s\) and \(F_Z\) for its derivatives. The weak
Euler--Lagrange identity for an admissible test \(\varphi\) is
\begin{equation}\label{eq:weak-el}
 \int_\Omega\bigl[F_s(x,u,Du)\cdot\varphi
                    +F_Z(x,u,Du):D\varphi\bigr]\,dx=0 .
\end{equation}

\begin{theorem}\label{thm:el-dominated}
Suppose \(J(u)\) is finite, and \(t=0\) is a local minimum of
\(J(u+t\varphi)\) for both signs of small \(t\), for each
\(\varphi\in C^\infty_c(\Omega;\R^m)\). Assume the integrand is
Carath\'eodory with continuous first derivatives in \((s,Z)\)
almost everywhere. For each such direction assume that for
\(|t|\) small the absolute value of
\[
 F_s(x,u+t\varphi,Du+tD\varphi)\cdot\varphi
   +F_Z(x,u+t\varphi,Du+tD\varphi):D\varphi
\]
is bounded by an integrable function of \(x\). Then
\eqref{eq:weak-el} holds for every such test. When the composed
\(F_s,F_Z\) are locally integrable, this says
\(F_s-\operatorname{div} F_Z=0\) in distributions, row by row.
If \(F_s\) and the divergence are continuous, the equation
holds pointwise.
\end{theorem}

\begin{proof}
For each \(x\) with the differentiability assumptions,
apply the fundamental theorem of calculus along the line
\(t\mapsto(u+t\varphi,Du+tD\varphi)\). The assumed domination
bounds the difference quotients and also gives integrability
of the change of energy. Dominated convergence differentiates
under the integral at zero. A two-sided minimum has zero
derivative, yielding \eqref{eq:weak-el}. The distributional
divergence is defined by pairing with \(-D\varphi\).
A continuous function which vanishes as a distribution is
zero pointwise, since near a nonzero component a nonnegative
test supported in a region of constant sign would give a
nonzero pairing. In particular this applies in the classical
setting of \(C^1\) integrand and \(C^2\) extremal when the
composed flux has continuous divergence.
\end{proof}

\begin{theorem}\label{thm:el-sobolev}
Let \(1\le p<\infty\). Suppose \(F\) and its first derivatives
have the Carath\'eodory regularity in
Theorem~\ref{thm:el-dominated} and there is a finite reference
energy on \(u_0+W^{1,p}_0\). For \(p>1\) assume
\[
 |F_s(x,s,Z)|+|F_Z(x,s,Z)|
 \le b(x)+C\bigl(|s|^{p-1}+|Z|^{p-1}\bigr),
 \qquad b\in L^{p'}(\Omega),
\]
where \(p'=p/(p-1)\). At \(p=1\) assume instead a common
\(L^\infty\) bound on the derivatives.
Then the energy is finite on the whole affine class and any
local minimizer in its \(W^{1,p}\) topology satisfies
\eqref{eq:weak-el} for every \(\varphi\in W^{1,p}_0\).
If only an upper \(p\)-growth bound on both derivatives is
available, of the form
\(b(x)+C(|s|^p+|Z|^p)\) with \(b\in L^1\),
a finite local minimizer still satisfies \eqref{eq:weak-el}
for compactly supported smooth tests.
\end{theorem}

\begin{proof}
In the \(p-1\) case the derivative of \(F\) along a segment
of two Sobolev competitors is bounded by
\(b+C_p(|u|^{p-1}+|v|^{p-1}+|Du|^{p-1}+|Dv|^{p-1})\)
times \(|u-v|+|Du-Dv|\). The coefficient is in \(L^{p'}\),
so the product is integrable. At \(p=1\) use the essential
bound times an \(L^1\) function. Thus a finite reference
gives finite energy everywhere on the class. Use the same
bound with \(v=u+\varphi\) and \(v=u-\varphi\) to dominate
difference quotients for \(|t|\le1\), or directly bound the
intermediate values by \(|u|+|\varphi|\) and
\(|Du|+|D\varphi|\). Differentiating the line and using
minimality gives the formula for the whole test class.
At \(p\)-growth, the intermediate coefficient is in
\(L^1\) for each smooth compactly supported test, and the
test and its derivative are bounded. This gives the required
domination for those tests. It supplies no \(L^{p'}\)
estimate for general Sobolev tests. Compare
\cite[Theorem 3.13 and Remark 3.14]{Dacorogna}.
\end{proof}

\section{Convex sufficiency and the natural boundary condition}
\label{v08-08}

\begin{proposition}\label{prop:weak-el-sufficient}
In the all-tests setting of Theorem~\ref{thm:el-sobolev},
assume \(F(x,\cdot,\cdot)\) is convex for almost every \(x\).
A weak stationary point is a global minimizer in the fixed
affine class.
\end{proposition}
\begin{proof}
For each competitor \(v\), convex differentiable support at
\((u,Du)\) yields
\[
 F(x,v,Dv)\ge F(x,u,Du)
       + F_s(x,u,Du)\cdot(v-u)
       + F_Z(x,u,Du):D(v-u).
\]
Integrate. The last two terms are integrable and sum to zero
by \eqref{eq:weak-el} with \(v-u\in W^{1,p}_0\).
The conclusion follows. For gradient-only power energies
the integrand is strictly convex in \(Du\) for \(p>1\);
Proposition~\ref{prop:sobolev-unique} then gives uniqueness.
\end{proof}

\begin{proposition}\label{prop:natural-boundary}
Suppose \(\Omega\) has \(C^1\) boundary, the composed flux
\(P=F_Z(x,u,Du)\) is \(C^1\) up to the boundary, and
the composed \(F_s\) is continuous. If the weak stationarity identity
\eqref{eq:weak-el} holds for all restrictions of smooth tests defined
near \(\overline\Omega\), then
\[
 F_s-\operatorname{div}P=0\quad\text{in }\Omega,\qquad
 P\nu=0\quad\text{on }\partial\Omega .
\]
Conversely these two equations give \eqref{eq:weak-el} against that
test class.
\end{proposition}
\begin{proof}
Compactly supported tests first give the interior equation.
Apply the divergence theorem to each row of \(P\) paired
with the test. The interior term cancels, leaving
\(\int_{\partial\Omega}(P\nu)\cdot\varphi=0\).
If one component of \(P\nu\) is nonzero at a boundary
point, its sign is fixed in a small relative neighborhood
by continuity of \(P\) and \(\nu\). Take a nonnegative
ambient bump supported there and positive near the point,
in that component of the test. The local \(C^1\) graph
has positive surface measure there, so its integral has
strict sign, a contradiction. The converse follows from
the same integration by parts. With prescribed boundary
values the test class instead has zero trace; the boundary
pairing then vanishes for every flux and gives no condition
on \(P\nu\).
\end{proof}
\chapter{Variational Duality and Fluxes}
\label{chap:flux-duality}

This chapter applies the Banach-space perturbation theorem to a
standard pair of function spaces. A dual flux is a vector field
whose weak divergence realizes the prescribed load. The
constraint on that field, the equality of the two values, and
the attainment statements will each be proved.

\section{Convex integral functionals and their conjugates}
\label{v09-01}

Let \(\Omega\) have finite Lebesgue measure and \(1<p<\infty\),
with conjugate exponent \(p'=p/(p-1)\). The continuous dual of
\(L^p(\Omega;\R^d)\) is represented by \(L^{p'}(\Omega;\R^d)\)
under the integral pairing. This is the usual \(L^p\) duality
on the Lebesgue measure space \cite[Section 4.3]{Brezis}.

\begin{theorem}\label{thm:integral-conjugacy}
Suppose \(F(x,z)\) is finite Carath\'eodory and convex in \(z\)
for almost every \(x\), and
\begin{equation}\label{eq:two-sided-p-growth}
 a_0|z|^p-a(x)\le F(x,z)\le a(x)+a_1|z|^p,
 \qquad a_0>0,\quad a_1\ge0,\quad 0\le a\in L^1(\Omega).
\end{equation}
Then \(G(v)=\int_\Omega F(x,v(x))\,dx\) is finite, convex and
continuous on \(L^p\). Its conjugate is
\[
 G^*(\sigma)=\int_\Omega F^*(x,\sigma(x))\,dx,\qquad
 F^*(x,w)=\sup_{z\in\R^d}\{\langle w,z\rangle-F(x,z)\}
 \quad(\sigma\in L^{p'}).
\]
The right-hand integral is finite.
\end{theorem}

\begin{proof}
The growth bound implies finiteness and
\(|F(x,z)|\le a(x)+a_1|z|^p\). Convexity follows pointwise.
For norm continuity, let \(v_n\to v\) in \(L^p\).
From any subsequence extract \(v_{n_j}\) with
\(\sum_j\norm{v_{n_j}-v}_p<\infty\). The pointwise sum
\(H=|v|+\sum_j|v_{n_j}-v|\) is in \(L^p\) by the triangle
inequality and monotone convergence, and
\(v_{n_j}\to v\) almost everywhere. The integrands are
dominated by \(a+a_1 H^p\), so their integrals converge by
dominated convergence. The subsequence criterion proves
continuity of \(G\) along the original sequence.

For conjugacy, the supremum defining \(F^*\) can be taken
over the rational vectors, by continuity in \(z\). Hence its
composition with a measurable \(\sigma\) is measurable.
The \(z=0\) value bounds it below by \(-a(x)\), and maximizing
\(|\sigma||z|-a_0|z|^p\) gives the upper bound
\(a(x)+C|\sigma(x)|^{p'}\). The integral is finite.
The pointwise Fenchel inequality gives
\(\int \sigma\cdot v-G(v)\le\int F^*(x,\sigma(x))\,dx\).

For the reverse inequality, enumerate the rational vectors
starting with zero. For the first \(n\) vectors choose at each
\(x\) the first maximizing index of
\(\sigma(x)\cdot z-F(x,z)\). This gives a measurable,
finitely valued \(v_n\in L^p\), using finite measure.
The values of the maximized integrand increase pointwise to
\(F^*(x,\sigma(x))\). They have the common integrable lower
bound \(-a(x)\), so monotone convergence after adding \(a\)
proves the reverse inequality. This justifies the interchange
of conjugation and integration used in the flux calculation.
Compare \cite[Chapter IV, Sections 1--2]{EkelandTemam}.
\end{proof}

\section{The Dirichlet--flux dual pair}
\label{v09-02}

\begin{theorem}\label{thm:dirichlet-flux}
Let \(\Omega\subset\R^d\) be bounded and open,
\(X=W^{1,p}_0(\Omega)\), \(1<p<\infty\), and \(\ell\in X'\).
Let \(F\) satisfy the hypotheses of Theorem~\ref{thm:integral-conjugacy}.
Then
\begin{equation}\label{eq:general-flux-dual}
 \min_{u\in X}\left\{\int_\Omega F(x,Du)\,dx-\ell(u)\right\}
  =\max_{\substack{\sigma\in L^{p'}(\Omega;\R^d)\\
                    \int_\Omega \sigma\cdot Dv=\ell(v)
                         \ \text{for all }v\in X}}
           \left\{-\int_\Omega F^*(x,\sigma)\,dx\right\}.
\end{equation}
Both extrema are finite and attained. The constraint on
\(\sigma\) is equivalently
\(-\operatorname{div}\sigma=\ell\) on smooth compactly
supported tests.
\end{theorem}

\begin{proof}
Let \(A:X\to L^p(\Omega;\R^d)\) be \(u\mapsto Du\).
Use \(f(u)=-\ell(u)\) and \(g=G\) from
Theorem~\ref{thm:integral-conjugacy}. By Poincar\'e the
gradient norm controls the norm of \(u\) in \(X\), so the
primal objective is bounded below by
\[
 a_0\norm{Du}_p^p-C\norm{Du}_p-\norm{a}_1 .
\]
It is finite at \(0\). The space \(X\) is reflexive by
Proposition~\ref{prop:sobolev-affine}. The objective is norm
continuous and convex, hence weakly lower semicontinuous
by Section~\ref{v02-03}. Coercivity and the direct method
of Section~\ref{v03-02} give a finite attained minimum.

For the dual, \(g\) is continuous at \(A0\) and the primal
value is finite, so the Banach-space theorem of
Section~\ref{v04-03} gives an attained maximum of
\(-f^*(-A^*\sigma)-g^*(\sigma)\).
Now \(f^*(q)=0\) when \(q=-\ell\) and \(+\infty\) otherwise,
since a nonzero continuous linear functional is unbounded
above on the whole space. Thus the dual constraint is
\(A^*\sigma=\ell\) with precisely the sign in
\eqref{eq:general-flux-dual}. The integral formula for \(g^*\)
gives the displayed dual. The distributional constraint is
the same identity on \(C^\infty_c(\Omega)\). Its extension
to \(X\) follows by the defining density and H\"older's
inequality for the flux. The argument does not need a
smooth boundary.
Compare \cite[Chapter III, Theorems 4.1--4.2, and Chapter IV]{EkelandTemam}.
\end{proof}

\begin{corollary}\label{cor:p-dirichlet-flux}
For \(F(z)=|z|^p/p\), the dual objective in
\eqref{eq:general-flux-dual} is
\(-\int_\Omega |\sigma|^{p'}/p'\), and the primal minimizer
is unique.
\end{corollary}
\begin{proof}
Young's inequality gives \(w\cdot z-|z|^p/p\le |w|^{p'}/p'\),
with equality by choosing \(z\) in the direction of \(w\) at
length \(|w|^{1/(p-1)}\), or \(z=0\) if \(w=0\).
This computes the conjugate. The function \(|z|^p\) is
strictly convex for \(p>1\), and the gradient map on \(X\)
is injective by Poincar\'e. Thus the primal objective is
strictly convex and its attained minimizer is unique.
\end{proof}

\section{Constitutive relations and the quadratic case}
\label{v09-03}

\begin{proposition}\label{prop:optimal-flux}
In Theorem~\ref{thm:dirichlet-flux}, a primal admissible \(u\)
and a dual feasible \(\sigma\) are both optimal if and only if
\[
 \sigma(x)\in\partial_z F(x,Du(x))\quad\text{almost everywhere}.
\]
For the power energy this is
\[
 \sigma=|Du|^{p-2}Du
 \quad\text{with value zero at }Du=0,
 \qquad
 \int_\Omega |Du|^{p-2}Du\cdot Dv=\ell(v)\quad(v\in X).
\]
\end{proposition}
\begin{proof}
Dual feasibility substitutes \(\ell(u)=\int\sigma\cdot Du\)
into the primal-minus-dual value. The resulting gap is
\[
 \int_\Omega
   \bigl[F(x,Du)+F^*(x,\sigma)-\sigma\cdot Du\bigr]\,dx.
\]
All terms are integrable by the growth assumptions and
H\"older's inequality. The integrand is nonnegative and
vanishes precisely at Fenchel equality, equivalently the
stated subgradient condition by Section~\ref{v02-06}.
Its integral vanishes exactly when it vanishes almost
everywhere. By the proved equality of optimal values,
zero gap for a feasible pair is equivalent to simultaneous
optimality. The derivative of \(|z|^p/p\) gives the power
constitutive law, including zero by differentiability there.
Dual feasibility yields the displayed weak \(p\)-Laplace
equation. Conversely a weak solution with this flux is
a zero-gap feasible pair, hence optimal.
\end{proof}

For \(p=2\), \(H^1_0\) with the gradient inner product is
Hilbert by Poincar\'e, and the weak equation is its
Riesz representation problem. If a bounded symmetric
coercive form \(a\) replaces the gradient inner product,
the quadratic energy \(\tfrac12 a(u,u)-\ell(u)\) has the
same weak characterization, by Section~\ref{v03-07}.
For a nonsymmetric coercive form the weak problem still
has the Lax--Milgram solution. Differentiating its diagonal
quadratic functional instead uses the symmetric part of the
form, since the mixed derivative is
\(\tfrac12(a(u,v)+a(v,u))\). In the complex convention of
Section~\ref{v03-07}, put
\(h(u,v)=\tfrac12(a(u,v)+\overline{a(v,u)})\).
The real derivative of
\(\tfrac12\operatorname{Re}a(u,u)-\operatorname{Re}\ell(u)\)
in direction \(v\) is \(\operatorname{Re}(h(u,v)-\ell(v))\).
Testing with \(v\) and \(iv\) recovers \(h(u,v)=\ell(v)\).
Real coercivity gives a solution for the full form by
Lax--Milgram, while the energy characterizes the solution
for its Hermitian part.

\section{An obstacle problem}
\label{v09-04}

\begin{theorem}\label{thm:hilbert-obstacle}
Let \(\Omega\subset\R^d\) be bounded open, \(\ell\in(H^1_0)'\),
and let \(\psi\) be measurable with
\[
 K=\{v\in H^1_0(\Omega):v\ge\psi\text{ almost everywhere}\}
\]
nonempty. Then the Dirichlet energy
\(\tfrac12\int_\Omega|Du|^2-\ell(u)\) has a unique minimizer
over \(K\). This minimizer is characterized by
\begin{equation}\label{eq:obstacle-vi}
 \int_\Omega Du\cdot D(v-u)\,dx\ge \ell(v-u)
 \qquad(v\in K).
\end{equation}
The distribution \(-\Delta u-\ell\) is nonnegative.
If \(u,\psi\) have continuous representatives in \(\Omega\),
then \(-\Delta u-\ell=0\) on the open set \(\{u>\psi\}\).
\end{theorem}

\begin{proof}
The pointwise constraint is convex and is closed under norm
convergence in \(H^1_0\): convergence in \(L^2\) supplies
an almost everywhere convergent subsequence, preserving
the inequality. Equip \(H^1_0\) with the equivalent gradient
Hilbert norm. The gradient form is bounded and coercive,
so Section~\ref{v03-06} gives a unique solution of
\eqref{eq:obstacle-vi}. For completeness, the energy
difference on a competitor \(v\) is
\[
 \tfrac12\norm{D(v-u)}_2^2+
 \int_\Omega Du\cdot D(v-u)-\ell(v-u).
\]
It is nonnegative by \eqref{eq:obstacle-vi}, and zero
forces \(v=u\) by Poincar\'e. Conversely, at a minimizer
differentiate the energy along \(u+t(v-u)\) for \(t\downarrow0\)
to obtain the inequality.

Taking \(v=u+\varphi\) for nonnegative smooth compactly
supported \(\varphi\) gives nonnegativity of
\(-\Delta u-\ell\) as a distribution. In the continuous case,
each compact subset of \(\{u>\psi\}\) has a strictly positive
minimum gap. Thus on the support of any test inside that open
set, both \(u+t\varphi\) and \(u-t\varphi\) are admissible for
small \(t>0\). The two inequalities imply equality of the
distributional pairing. Compare
\cite[Chapter II, Section 2, and Chapter IV, Section 3]{EkelandTemam}
and \cite[Section 5.3]{Brezis}.
\end{proof}
\chapter{Fixed Points and an Existence Principle}

\section{Sperner's lemma and Brouwer's theorem}
\label{v10-01}

Write \(\Delta^d=\{x\in\mathbb R^{d+1}:x_i\ge0,\,
\sum_{i=0}^d x_i=1\}\), and identify the \(i\)-th vertex with its
coordinate index. In a finite geometric triangulation of \(\Delta^d\),
a Sperner labeling assigns each vertex \(x\) an index \(i\)
such that \(x_i>0\). Thus every vertex in a face receives one of
the indices of that face.

\begin{lemma}[Sperner parity]
\label{lem:sperner-parity}
The number of \(d\)-simplices in a Sperner-labeled triangulation
whose vertices carry all \(d+1\) labels is odd.
\end{lemma}
\begin{proof}
In dimension zero the number is one. In dimension \(d\ge1\)
count pairs \((\sigma,\tau)\) where \(\sigma\) is a
\(d\)-simplex and \(\tau\) is one of its facets carrying
the \(d\) distinct labels \(0,\ldots,d-1\). Every interior
facet is incident to two \(d\)-simplices, and a boundary
facet to one. To see this, at a relative-interior point
of a facet a sufficiently small ball meets only simplices
containing that facet. A containing \(d\)-simplex occupies
one side of its affine span near the point. Disjoint
interiors allow at most one on each side, and coverage of
\(\Delta^d\) gives one on every side present in \(\Delta^d\).
A boundary facet counted here must lie on the face
\(x_d=0\), since a boundary face excluding another coordinate
forbids its label. The induced triangulation of \(x_d=0\)
has an odd number of fully labeled facets by induction.
Thus the total incidence count is odd.

From the \(d\)-simplex side, a fully labeled simplex
contributes one facet to the count. A simplex containing
all of \(0,\ldots,d-1\) but missing \(d\) has one repeated
label and contributes two facets. All other simplices
contribute none. Reducing the count modulo two proves the claim.
\end{proof}

\begin{lemma}[Fine triangulations]
\label{lem:fine-simplex-triangulations}
The simplex \(\Delta^d\) admits a sequence of finite geometric
triangulations with mesh tending to zero and face-compatible
restrictions.
\end{lemma}
\begin{proof}
For a simplex with vertex set \(I\), let \(b_F\) be the barycenter
of the vertices in a nonempty face \(F\subset I\).
The simplices spanned by the barycenters of full nested flags
give its barycentric subdivision. For completeness, if the
barycentric coordinates satisfy
\(\lambda_{i_0}\ge\cdots\ge\lambda_{i_d}\ge0\), set
\(F_k=\{i_0,\ldots,i_k\}\), and let \(\lambda_{i_{d+1}}=0\).
The point with these coordinates has the expansion
\[
 \sum_{k=0}^{d}(k+1)(\lambda_{i_k}-\lambda_{i_{k+1}})b_{F_k}.
\]
The coefficients are nonnegative and sum to one, and checking
each vertex coordinate gives the expansion. The \(b_{F_k}\)
are affinely independent, as their coordinate matrix in this
ordering is triangular with nonzero diagonal. Where orderings
tie, positive coefficients occur only at strict decreases.
Their associated prefixes are the same for any ordering
of the tied coordinates. Hence intersections are common
faces. This also proves compatibility with the original faces,
where the excluded coordinates are zero. Subdivisions of
neighboring simplices therefore fit together.

For \(F\subset G\) nonempty and distinct,
\(b_G=(|F|/|G|)b_F+(1-|F|/|G|)b_{G\setminus F}\).
If the original diameter is \(D\), the distance of any two
such nested barycenters is at most \(dD/(d+1)\). This bound
extends to the whole flag simplex by writing differences
of convex combinations as averages of pairwise differences.
When \(d\ge1\), repeated subdivisions reduce the mesh by
the factor \(d/(d+1)<1\). For \(d=0\) the singleton suffices.
\end{proof}

\begin{theorem}[Brouwer fixed point theorem]
\label{thm:brouwer-compact-convex}
Let \(K\) be nonempty, compact, and convex in a
finite-dimensional real normed space. Every continuous
map \(F:K\to K\) has a fixed point.
\end{theorem}
\begin{proof}
First consider \(\Delta^d\). In each fine triangulation,
if a vertex is already fixed, the conclusion holds.
Otherwise label a vertex \(x\) by an index with \(x_i>F_i(x)\).
Such an index exists since the coordinates of \(x\) and
\(F(x)\) each sum to one, and it has \(x_i>0\).
By Lemma~\ref{lem:sperner-parity} choose a fully labeled simplex
in each triangulation. By compactness take a subsequence on
which its vertices of label zero converge to \(x_*\).
The mesh bound in Lemma~\ref{lem:fine-simplex-triangulations}
makes the vertices of every other label converge to the
same point. Passing to limits gives \(x_{*,i}\ge F_i(x_*)\)
for each \(i\). Equal sums give equality in every coordinate.

For general \(K\), choose Euclidean coordinates and their
auxiliary inner product, with the same topology by finite
dimensionality. Dimension zero is immediate. In dimension
\(d>0\), choose \(M>0\) so that all coordinates of points
in \(K\) have absolute value at most \(M\).
The region defined by \(x_i\ge-M\) for \(1\le i\le d\)
and \(\sum_i x_i\le dM\) is a nondegenerate simplex \(S\)
containing \(K\). This follows also by translating the
coordinates by \(M\) and scaling the simplex with
nonnegative coordinates and sum at most \(2dM\).
The Hilbert projection \(P_K:S\to K\) is a continuous
retraction by Theorem~\ref{thm:hilbert-projection}.
Apply the simplex case to \(F\circ P_K:S\to S\),
after affine change of coordinates. A fixed point of
this map lies in \(K\), where \(P_K\) is the identity,
so it is fixed by \(F\).
\end{proof}

This is the compact convex form used in
\cite[Exercise 6.26]{Brezis}. The simplex argument also covers
lower-dimensional and singleton convex subsets via the retraction.

\section{Schauder through finite-dimensional approximation}
\label{v10-02}

\begin{theorem}[Schauder]
\label{thm:schauder-fixed-point}
Let \(C\) be nonempty, closed, and convex in a real Banach space
and \(F:C\to C\) be continuous. If \(F(C)\) is relatively compact,
then \(F\) has a fixed point.
\end{theorem}
\begin{proof}
Let \(K=\overline{F(C)}\), a nonempty compact subset of \(C\).
For \(\epsilon>0\), choose a finite \(\epsilon/2\)-net
\(y_1,\ldots,y_N\) in \(K\) and set
\[
 w_j(y)=(\epsilon-\|y-y_j\|)_+,\qquad
 P_\epsilon(y)=\frac{\sum_{j=1}^{N}w_j(y)y_j}
                     {\sum_{j=1}^{N}w_j(y)}\quad(y\in K).
\]
The denominator is positive by the net property.
This is a continuous map into the finite-dimensional compact
convex hull \(C_\epsilon=\operatorname{conv}\{y_1,\ldots,y_N\}
\subset C\). The hull is compact as the continuous image
of a finite-dimensional simplex. Positive weights occur
only at distances less than \(\epsilon\), so
\(\|P_\epsilon(y)-y\|\le\epsilon\).
Brouwer applied to \(P_\epsilon\circ F:C_\epsilon\to C_\epsilon\)
gives \(x_\epsilon\) with
\(\|x_\epsilon-F(x_\epsilon)\|\le\epsilon\).

Take \(\epsilon\) tending to zero through a sequence.
Compactness of \(K\) yields a subsequence with
\(F(x_\epsilon)\to x\in K\). The error bound gives
\(x_\epsilon\to x\), and continuity gives \(F(x)=x\).
The argument uses relative compactness of the image and
allows an unbounded \(C\).
\end{proof}

See \cite[Exercise 6.26]{Brezis}.

\section{A Schaefer continuation criterion}
\label{v10-03}

\begin{theorem}[Schaefer]
\label{thm:v10-03-01}
Let \(X\) be a real Banach space and let \(T:X\to X\) be
continuous and map bounded sets to relatively compact sets.
If the set
\[
 \{x\in X:x=tT(x)\text{ for some }t\in[0,1]\}
\]
is bounded, then \(T\) has a fixed point.
\end{theorem}
\begin{proof}
Choose \(R>0\) strictly greater than all norms in the
displayed set. The radial truncation \(r_R(y)=y\) for
\(\|y\|\le R\), and \(r_R(y)=Ry/\|y\|\) otherwise, is
continuous. On the closed ball of radius \(R\),
\(r_R\circ T\) is a continuous self-map with relatively
compact image, since the closure of \(T\)'s image of
this ball is compact and \(r_R\) is continuous.
Theorem~\ref{thm:schauder-fixed-point} supplies a fixed
point \(x\) of the composition.
If \(\|T(x)\|>R\), then \(x=(R/\|T(x)\|)T(x)\) and
\(\|x\|=R\), contrary to the choice of \(R\).
Thus \(r_R(T(x))=T(x)\), giving \(x=T(x)\).
\end{proof}

\section{Boundary examples and the contraction principle}
\label{v10-04}

\begin{example}[Why the hypotheses matter]
The rotation \(x\mapsto-x\) of the unit circle in \(\mathbb R^2\)
is a continuous fixed-point-free self-map of a compact nonconvex
set. The map \(x\mapsto (1+x)/2\) is a continuous self-map of the
bounded convex open interval \((0,1)\) with its only possible
fixed point at the omitted endpoint \(1\).
On the closed unit ball \(B\subset\ell^2\) the map
\[
 T(x)=\bigl(\sqrt{1-\|x\|^2},x_1,x_2,\ldots\bigr)
\]
is continuous and has image in the unit sphere. At a fixed
point all coordinates would equal the first, and square
summability would force all coordinates to vanish. Yet
\(T(0)\ne0\). Also \(T(e_n)=e_{n+1}\), so the image is not
relatively compact. Thus a closed bounded convex set in
infinite dimension alone gives no fixed point guarantee.
\end{example}

\begin{proposition}[Banach contraction principle]
\label{prop:v10-04-01}
If \(M\) is a nonempty complete metric space and
\(T:M\to M\) is Lipschitz with constant \(q<1\), then
\(T\) has a unique fixed point. From any \(x_0\) the iterates
\(x_{n+1}=T(x_n)\) converge to it with
\(d(x_n,x_*)\le q^n d(x_0,x_1)/(1-q)\).
\end{proposition}
\begin{proof}
The successive distances are bounded by
\(q^n d(x_0,x_1)\). Summing the geometric series shows that
the sequence is Cauchy and gives the stated tail bound.
Completeness gives a limit, and continuity of \(T\) gives
the fixed point identity. For two fixed points the
distance is at most \(q\) times itself, proving uniqueness.
\end{proof}

In Section~\ref{v03-06} the contraction comes from coercivity
and projection, so it provides uniqueness and an iteration.
The compactness arguments above provide existence.

\backmatter
\begin{thebibliography}{99}
\bibitem{Rockafellar}
R.~T. Rockafellar, \emph{Convex Analysis}. Princeton University Press,
1970. Second printing, 1972.
\bibitem{Boyd}
S. Boyd and L. Vandenberghe, \emph{Convex Optimization}.
Cambridge University Press, 2004. Corrected seventh printing, 2009.
\bibitem{EkelandTemam}
I. Ekeland and R. T\'emam, \emph{Convex Analysis and Variational
Problems}. SIAM, 1999.
\bibitem{Brezis}
H. Brezis, \emph{Functional Analysis, Sobolev Spaces and Partial
Differential Equations}. Springer, 2011.
\bibitem{Dacorogna}
B. Dacorogna, \emph{Introduction to the Calculus of Variations},
fourth edition. World Scientific, 2025.
\bibitem{NocedalWright}
J. Nocedal and S.~J. Wright, \emph{Numerical Optimization},
first edition. Springer, 1999.
\bibitem{Liberzon}
D. Liberzon, \emph{Calculus of Variations and Optimal Control Theory:
A Concise Introduction}. Princeton University Press, 2012.
\bibitem{MangasarianFromovitz}
O.~L. Mangasarian and S. Fromovitz, The Fritz John necessary optimality
conditions in the presence of equality and inequality constraints.
\emph{Journal of Mathematical Analysis and Applications}
17 (1967), 37--47.
\bibitem{MathGenomeP102}
MathGenome, \emph{Fourier Analysis, Distributions, Sobolev Spaces, and
Approximation}. Unified notes P1-02, 2026. Unit U1, periodic Fourier
series, and Unit U5, extension, compactness, and trace.
\bibitem{MathGenomeD04}
MathGenome, \emph{Multivariable Analysis}. Notes D04, Sections 9 and 12,
regular arc length and Green's theorem.
\bibitem{MathGenomeD13}
MathGenome, \emph{Ordinary Differential Equations and Dynamical Systems}.
Notes D13, Chapters 7 and 21, linear systems and smooth Jordan boundaries.
\end{thebibliography}
\end{document}
